%!TEX program = pdflatex
\documentclass[aps,%
twocolumn,%
floatfix,%
pra,%
amsfonts,%
groupedaddress,%
superscriptaddress]{revtex4-2}%

%%%%%%%%%%%%%%%%%%%%%%%%%%%%%%%%%%%%%%%%%%%%%%%%%%%%%%%%%%%%%%%%%%%%%%%%%%%%%%%%%%%%%%%%%%%%%%%%%%%
% Packages
%%%%%%%%%%%%%%%%%%%%%%%%%%%%%%%%%%%%%%%%%%%%%%%%%%%%%%%%%%%%%%%%%%%%%%%%%%%%%%%%%%%%%%%%%%%%%%%%%%%
\usepackage{amsmath}
\usepackage{amssymb}
\usepackage{amsthm}
\usepackage{bm,bbm} % to use the command \mathbbm{1}, \bm{x}
\usepackage[table,dvipsnames,usenames]{xcolor}
\usepackage{graphicx}
\usepackage{indentfirst}
\usepackage{enumitem}
\usepackage{algorithm2e}
\usepackage{algpseudocode}
\usepackage{hyperref}
% \usepackage{biblatex}
\usepackage{tikz}
\usetikzlibrary{graphs} 
\usepackage[titletoc,toc,page,title]{appendix}
\definecolor{beamer@blendedblue}{rgb}{0.2,0.2,0.7}      % color used in beamer

%%%%%%%%%%%%%%%%%%%%%%%%%%%%%%%%%%%%%%%%%%%%%%%%%%%%%%%%%%%%%%%%%%%%%%%%%%%%%%%%%%%%%%%%%%%%%%%%%%%
% General Commands
%%%%%%%%%%%%%%%%%%%%%%%%%%%%%%%%%%%%%%%%%%%%%%%%%%%%%%%%%%%%%%%%%%%%%%%%%%%%%%%%%%%%%%%%%%%%%%%%%%%
\newtheorem{theorem}{Theorem}
\newtheorem{proposition}[theorem]{Proposition}
\newtheorem{lemma}{Lemma}
\newtheorem{corollary}{Corollary}
\newtheorem{problem}{Problem}
\newtheorem{definition}{Definition}[section]
%%%%%%%%%%%%%%%%%%%%%%%%%%%%%%%%%%%%%%%%%%%%%%%%%%%%%%%%%%%%%%%%%%%%%%%%%%%%%%%%%%%%%%%%%%%%%
% TODO List
% https://tex.stackexchange.com/questions/247681/how-to-create-checkbox-todo-list
%% to use \ding{67} symbols
%% https://math.uoregon.edu/wp-content/uploads/2014/12/compsymb-1qyb3zd.pdf
\usepackage{pifont}
\newlist{todolist}{itemize}{2}
\setlist[todolist]{label=$\square$}
\newcommand{\cmark}{\ding{51}}%
\newcommand{\xmark}{\ding{55}}%
\newcommand{\done}{\rlap{$\square$}{\raisebox{2pt}{\large\hspace{1pt}\cmark}}\hspace{-2.5pt}}
\newcommand{\wontfix}{\rlap{$\square$}{\large\hspace{1pt}\xmark}}
% USAGE:
% \begin{todolist}
%   \item[\done] Frame the problem
%   \item Write solution
%   \item[\wontfix] profit
%  \end{todolist}

%%%%%%%%%%%%%%%%%%%%%%%%%%%%%%%%%%%%%%%%%%%%%%%%%%%%%%%%%%%%%%%%%%%%%%%%%%%%%%%%%%%%%%%%%%%%%
% Quantum Commands
%%%%%%%%%%%%%%%%%%%%%%%%%%%%%%%%%%%%%%%%%%%%%%%%%%%%%%%%%%%%%%%%%%%%%%%%%%%%%%%%%%%%%%%%%%%%%
\DeclareMathOperator{\Tr}{Tr}
\DeclareMathOperator{\tr}{Tr}
\DeclareMathOperator{\NOT}{NOT}
\DeclareMathOperator{\diag}{diag}
\DeclareMathOperator{\rank}{rank}
\DeclareMathOperator{\conv}{conv}
\DeclareMathOperator{\sgn}{sgn}
\DeclareMathOperator{\Sp}{Span}
\DeclareMathOperator{\id}{id}
\newcommand{\bra}[1]{\langle #1\rvert}
\newcommand{\ket}[1]{\lvert #1\rangle}
\newcommand{\proj}[1]{\lvert #1\rangle\!\langle #1\rvert}
\newcommand{\ox}{\otimes}
\newcommand{\1}{\mathbbm{1}}
\newcommand{\mean}[1]{\langle #1\rangle}
\newcommand{\gen}[1]{\langle #1\rangle}
\newcommand{\ketbra}[2]{\lvert #1\rangle\!\langle #2\rvert}
\newcommand{\braket}[2]{\langle #1\vert #2\rangle}
%%%%%%%%%%%%%%%%%%%%%%%%%%%%%%%%%%%%%%%%%%%%%%%%%%%%%%%%%%%%%%%%%%%%%%%%%%%%%%%%%%%%%%%%%%%%%
%%% Mathematical symbols
%%%%%%%%%%%%%%%%%%%%%%%%%%%%%%%%%%%%%%%%%%%%%%%%%%%%%%%%%%%%%%%%%%%%%%%%%%%%%%%%%%%%%%%%%%%%%
\newcommand*{\cF}{\mathcal{F}}
\newcommand*{\cH}{\mathcal{H}}
\newcommand*{\cO}{\mathcal{O}}
\newcommand*{\cP}{\mathcal{P}}
\newcommand*{\cS}{\mathcal{S}}
\newcommand*{\bF}{\mathbb{F}}
\newcommand*{\bP}{\mathbb{P}}
\newcommand*{\cA}{\mathcal{A}}
\newcommand*{\cD}{\mathcal{D}}
\newcommand*{\cM}{\mathcal{M}}
\newcommand*{\cU}{\mathcal{U}}
\newcommand*{\cJ}{\mathcal{J}}
\newcommand*{\cL}{\mathcal{L}}
\DeclareMathOperator{\sym}{sym}
\newcommand*{\cB}{\mathcal{B}}
\newcommand*{\cI}{\mathcal{I}}
\DeclareMathOperator{\Alg}{Alg}

% Colors
\definecolor{alizarin}{rgb}{0.82, 0.1, 0.26}
\definecolor{googleblue}{HTML}{4285F4}
\definecolor{googlered}{HTML}{DB4437}
\definecolor{googleyellow}{HTML}{F4B400}
\definecolor{googlegreen}{HTML}{0F9D58}
\definecolor{klevinblue}{HTML}{002FA7}
\definecolor{tiffanyblue}{HTML}{0ABAB5}

\newcommand{\KW}[1]{{\color{googlered}(KW: #1)}}
\newcommand{\SYC}[1]{\textcolor{blue}{(SYC: #1)}}

% Define the PRL section style
% USAGE: \prlsection{XXX.}
\newcommand\prlsection[1]{\textit{\textbf{#1}}---}
\newcommand\prlsubsection[1]{\paragraph*{#1\!\!\!\!}}

\begin{document}

% Optional titles: 
%%% Quantum state verification achiving global optimal efficiency
%%% Quantum memory assisted entangled state verification with local measurements
%%% Optimal Verification of Entangled States assisted by Quantum Memory
%%% Quantum Memory Assisted Entangled State Verification
%%% Optimal Verification of Entangled States Assisted by Quantum Memory
%%% Memory Effects in Quantum State Verification
%%% Optimal Verification of Entangled States with Quantum Memory
% Remark: Since we cannot declare the optimality in the general case, we'd better focus on 
% Advantage of quantum memory, which is commonly termed as "Memory effects in Quantum XXX":
%%% Memory Effects in Quantum Channel Discrimination
%%% Quantum channels and memory effects
%%% Memory effects in quantum processes
%%% Memory Effects in Quantum Metrology
%%% Loss of coherence and memory effects in quantum dynamics
%%% Memory Effects in Quantum Dynamics Modelled by Quantum Renewal Processes
%%% Memory effects in quantum information transmission across a Hamiltonian dephasing channel
\title{Multi-copy shadow tomography}
\author{Siyuan Chen}
\affiliation{Institute for Quantum Computing, Baidu Research, Beijing 100193, China}%
\affiliation{Physics Department, Jilin University}%
\begin{abstract}
\noindent
We study the problem of simultaneously estimating nonlinear polynomial functionals of a quantum state $\rho$ using multi-copy measurements on $\rho^{\ox k}$.
We introduce the \emph{f-compatibility} framework, which generalizes quantum measurement compatibility to functional estimation and serves as a complexity calibrator:
the f-compatibility robustness $s^*$ determines the sample complexity of any estimation strategy.
For local (product) measurements on $n$-qubit systems, we prove a \emph{tensor product factorization theorem}: the robustness decomposes as $s^*_{\mathrm{loc}} = \prod_l s^*_l$ over qubit sites, recovering the $3^w$ scaling of classical shadows from a per-site incompatibility of $\sqrt{3}$.
Crucially, nonlinear functionals can be strictly easier to estimate than their parent observables: for squared Pauli expectations $\{|\tr(P\rho)|^2\}$, the ``squaring trick'' $[P\ox P, Q\ox Q] = 0$ yields $s^* = 1$ via Bell measurements, achieving $\mathcal{O}(1)$ sample complexity.
We construct an explicit nonlocal $2k$-copy strategy based on a commuting operator family $\{A_P\}$ that simultaneously estimates all $|\tr(P\rho^k)|$ via a single projective measurement.
\color{red}Finally, we establish a local obstruction: the relevant lifted operators have non-product structure, and we prove that two-copy local 3-design strategies incur exponential cost $\sim (11/3)^n$ for product states (up to $5^n$ for entangled states) for estimating $\tr(P\rho^2)$.\color{black}
\end{abstract}
\date{\today}
\maketitle

%%%%%%%%%%%%%%%%%%%%%%%%%%%%%%%%%%%%%%%%%%%%%%%%%%%%%%%%%%%%%%%%%%%%%%%%%%%%%%%%%%%%%%%%%%%%%%%%%%%
%%%%%%%%%%%%%%%%%%%%%%%%%%%%%%%%%%%%%%%%%%%%%%%%%%%%%%%%%%%%%%%%%%%%%%%%%%%%%%%%%%%%%%%%%%%%%%%%%%%
\prlsection{Introduction.}
Classical shadow tomography~\cite{huang2020predicting} enables predicting many linear properties $\tr(O_i\rho)$ of an $n$-qubit state $\rho$ from few randomized measurements. Under local Clifford measurements, estimating a weight-$w$ Pauli observable $P$ requires $\mathcal{O}(3^w)$ samples, originating from the per-qubit depolarization factor $(d+1) = 3$ of the inverse shadow channel. This scaling is optimal among single-copy local strategies~\cite{chen2024optimal}.

A natural extension uses \emph{multi-copy} measurements on $\rho^{\ox k}$, exploiting quantum correlations across replicas. A striking example is Bell sampling on $\rho^{\ox 2}$: measuring in the Bell basis simultaneously estimates all squared Pauli expectations $|\tr(P\rho)|^2$ with $\mathcal{O}(1)$ samples~\cite{hangleiter2024bell,montanaro2017learning}, an exponential improvement over the single-copy $3^w$ scaling.

What determines when multi-copy strategies help? For \emph{nonlinear} functionals such as $\tr(P\rho^2)$, the question is particularly rich: two-copy local strategies already incur exponential cost $\sim (11/3)^n$--$5^n$ depending on the state (Appendix~\ref{app:local_shadow}), yet an explicit nonlocal four-copy circuit achieves $\mathcal{O}(1)$~\cite{ekert2002direct}. This locality gap calls for a unified framework that simultaneously captures:
(i)~the intrinsic hardness of multi-copy estimation,
(ii)~the factorization of hardness across qubit sites under local measurements, and
(iii)~the advantage of entangled (nonlocal) strategies.

In this Letter, we develop the \emph{f-compatibility} framework, bridging quantum measurement compatibility theory~\cite{heinosaari2015noise,designolle2019incompatibility,Busch1986} with shadow estimation complexity. Our main results are: (1)~a general framework linking estimation robustness $s^*$ to sample complexity (Sec.~\ref{sec:framework}); (2)~a tensor product factorization theorem explaining the $3^w$ scaling from per-site $\sqrt{3}$ incompatibility (Sec.~\ref{sec:tensor}); (3)~the squaring trick showing that nonlinear functionals can be exponentially easier (Sec.~\ref{sec:nonlinear}); (4)~an explicit nonlocal $2k$-copy strategy (Sec.~\ref{sec:nonlocal}); and (5)~a local no-go establishing exponential cost for $\tr(P\rho^2)$ under product measurements (Sec.~\ref{sec:local_nogo}). Related concurrent work~\cite{chen2025simultaneous,ye2025replica,du2025optimal} studies nonlinear estimation from complementary perspectives.

%%%%%%%%%%%%%%%%%%%%%%%%%%%%%%%%%%%%%%%%%%%%%%%%%%%%%%%%%%%%%%%%%%%%%%%%%%%%%%%%%%%%%%%%%%%%%%%%%%%
\prlsection{Framework.}\label{sec:framework}
Consider $K$ observables $\{A_i\}_{i=1}^K$ on an $n$-qubit Hilbert space $\cH$, and a family of polynomial functionals $\cF = \{f_\alpha\}_{\alpha \in \cA}$ of degree $\leq d$ in the outcome probability vectors $\vec{p}_i(\rho) = (\tr[A_i^{(r)}\rho])_r$. By the classical estimation condition~\cite{wu2020polynomial}, unbiased estimation of degree-$d$ polynomials requires at least $d$ i.i.d.\ samples.

In the quantum setting, $k$ copies of $\rho$ enable a joint POVM $\{\Pi_j\}$ on $\cH^{\ox k}$, whose outcomes are post-processed by classical estimators $h_\alpha(j)$.

\begin{definition}[f-compatibility]\label{def:f-compat}
$\{A_i\}$ is \emph{$k$-copy $\cF$-compatible} if there exists a POVM $\{\Pi_j\}_{j\in\cJ}$ on $\cH^{\ox k}$ and bounded estimators $h_\alpha:\cJ\to[-1,1]$ such that
\begin{align}\label{eq:f-compat}
f_\alpha\big(\vec{p}_1(\rho),\ldots,\vec{p}_K(\rho)\big) = \sum_{j} h_\alpha(j)\;\tr[\Pi_j\;\rho^{\ox k}]
\end{align}
for all $\rho$ and all $\alpha \in \cA$.
\end{definition}

Any degree-$d$ polynomial $f_\alpha$ admits a unique \emph{linearization}: a Hermitian operator $F_\alpha$ on $\cH^{\ox k}$ ($k \geq d$) such that $f_\alpha(\vec{p}(\rho)) = \tr[F_\alpha \rho^{\ox k}]$ for all $\rho$. Condition~\eqref{eq:f-compat} becomes $\sum_j h_\alpha(j)\Pi_j = F_\alpha$.

The \emph{f-compatibility robustness} $s^*(\cF)$ is the minimum super-normalization $s \geq 1$ satisfying $\sum_j \Pi_j = s\cdot\1^{\ox k}$ such that the above conditions are feasible. This optimal $s^*$ is computed by the following optimization (a bilinear program in $h_\alpha(j)$ and $\Pi_j$, which admits an SDP relaxation; see Appendix~\ref{app:f-compat}):
\begin{align}\label{eq:f-compat-SDP}
s^* = \min\; s \;\;\text{s.t.}\;\; &\sum_j h_\alpha(j)\Pi_j = F_\alpha,\; \sum_j \Pi_j = s\1, \notag\\
&\Pi_j \geq 0,\; |h_\alpha(j)| \leq 1.
\end{align}
The \emph{copy-complexity trade-off} follows: estimating $|\cA|$ functionals within error $\varepsilon$ with confidence $1-\delta$ requires
\begin{align}\label{eq:copy-complexity}
N = \mathcal{O}\!\left(k\cdot [s^*(k)]^2\cdot \varepsilon^{-2}\cdot \log(|\cA|/\delta)\right)
\end{align}
copies of $\rho$. The optimal copy number minimizes $k\cdot [s^*(k)]^2$ over $k$.

%%%%%%%%%%%%%%%%%%%%%%%%%%%%%%%%%%%%%%%%%%%%%%%%%%%%%%%%%%%%%%%%%%%%%%%%%%%%%%%%%%%%%%%%%%%%%%%%%%%
\prlsection{Tensor product structure.}\label{sec:tensor}
For $n$-qubit systems, a \emph{local} $k$-replica strategy uses product POVMs $\Pi_{\vec{j}} = \bigotimes_{l=1}^n \Pi_{j_l}^{(l)}$ on $(\mathbb{C}^2)^{\ox k}$ per site, with factored estimators $h_\alpha(\vec{j}) = \prod_l h_{\alpha_l}^{(l)}(j_l)$.

\begin{theorem}[Tensor product factorization]\label{thm:tensor}
If the linearizations have product form $F_\alpha = \bigotimes_{l=1}^n F_{\alpha_l}^{(l)}$, then the local f-compatibility robustness factorizes:
\begin{align}\label{eq:tensor-factor}
s^*_{\mathrm{loc}}(\cF) = \prod_{l=1}^n s^*_l,
\end{align}
where $s^*_l$ is the single-site robustness from~\eqref{eq:f-compat-SDP} restricted to site $l$.
\end{theorem}

This reduces the $n$-qubit problem to $n$ independent single-qubit optimizations of dimension $2^k$. The sample complexity~\eqref{eq:copy-complexity} with local measurements becomes $N = \mathcal{O}\!\left((\prod_l s^*_l)^2\varepsilon^{-2}\log(|\cA|/\delta)\right)$.

\emph{Recovering $3^w$.}---For single-copy ($k=1$) linear Pauli estimation $f_i(\rho) = \tr(P_i\rho)$ with $P_i = \bigotimes_l P_{i_l}$, the linearization is $F_i = P_i$ itself. At each non-identity site, the single-qubit optimization asks for joint measurability of unsharp Paulis $\{\frac{\1 + \eta P}{2}\}_{P\in\{X,Y,Z\}}$, which requires $\eta \leq 1/\sqrt{3}$~\cite{Busch1986,heinosaari2015noise}. This gives $s^*_l = \sqrt{3}$ per non-identity site, yielding $s^*_{\mathrm{loc}} = (\sqrt{3})^w$ and variance $\sim 3^w$ for weight-$w$ Paulis, matching the classical shadow result~\cite{huang2020predicting}.

%%%%%%%%%%%%%%%%%%%%%%%%%%%%%%%%%%%%%%%%%%%%%%%%%%%%%%%%%%%%%%%%%%%%%%%%%%%%%%%%%%%%%%%%%%%%%%%%%%%
\prlsubsection{When nonlinearity helps.}\label{sec:nonlinear}
For \emph{linear} functionals of binary (two-outcome) POVMs, f-compatibility is equivalent to standard joint measurability~\cite{designolle2019incompatibility}: estimation is exactly as hard as compatibility.

\begin{theorem}[No advantage for linear binary]\label{thm:linear-equiv}
For $K$ two-outcome POVMs $\{A_i^{(\pm)}\}$ and $K$ linearly independent linear functionals $f_i(\vec{p}) = \tr(P_i\rho)$, the f-compatibility copy number equals the standard compatibility number.
\end{theorem}

The situation changes dramatically for \emph{nonlinear} functionals. Consider the squared expectations $f_i(\rho) = |\tr(P_i\rho)|^2$ for Pauli observables $\{P_i\} \subseteq \{X,Y,Z\}$. The two-copy linearizations are $F_i = P_i \ox P_i$. The key observation---the \emph{squaring trick}---is:
\begin{align}\label{eq:squaring}
[P_i\ox P_i,\; P_j\ox P_j] = (P_iP_j)\ox(P_iP_j) - (P_jP_i)\ox(P_jP_i) = 0,
\end{align}
even when $[P_i, P_j] \neq 0$. Since $\{X\ox X, Y\ox Y, Z\ox Z\}$ all commute, they share a common eigenbasis---the Bell basis. Bell measurement thus achieves $s^* = 1$, giving $\mathcal{O}(1)$ sample complexity for all squared Pauli expectations simultaneously.

By the tensor product theorem, the per-site robustness is $s^*_l = 1$ at \emph{every} site, giving $s^*_{\mathrm{loc}} = 1$ for $n$-qubit squared Paulis. This explains the efficiency of Bell sampling~\cite{hangleiter2024bell}: it is an instance of f-compatible estimation with zero robustness penalty.



%%%%%%%%%%%%%%%%%%%%%%%%%%%%%%%%%%%%%%%%%%%%%%%%%%%%%%%%%%%%%%%%%%%%%%%%%%%%%%%%%%%%%%%%%%%%%%%%%%%
\prlsection{Nonlocal multi-copy strategies.}\label{sec:nonlocal}
For the more challenging nonlinear functionals $\tr(P\rho^k)$ (not merely $|\tr(P\rho)|^k$), we construct an explicit nonlocal strategy.

\begin{theorem}[Nonlocal Pauli shadow]\label{thm:nonlocal}
Let $\bP^d_\pi$ be the permutation operator on $2k$ replicas with $\pi = (\mathrm{odds})(\mathrm{even})$ acting as $k$-cycles on odd and even labels respectively, and let $\sym_\pi[A] := \frac{1}{k}\sum_{j=0}^{k-1} \pi^j A\, (\pi^{j})^\dagger$ denote symmetrization under the cyclic group generated by $\pi$. Define
\begin{align}\label{eq:AP}
A_P = \sym_\pi\!\left[P^{\ox 2}\ox(\1^{\ox 2})^{\ox(k-1)}\right]\cdot\bP^d_\pi.
\end{align}
Then $\cA = \{A_P\}_{P\in\cP_n}$ is a commuting family with $\tr(A_P\rho^{\ox 2k}) = [\tr(P\rho^k)]^2$, so a single projective measurement in its common eigenbasis simultaneously estimates all $|\tr(P\rho^k)|$.
\end{theorem}

The commutativity follows from $[S_P, S_Q] = 0$ since each pair of terms either acts on disjoint subsystems or involves $[P^{\ox 2}, Q^{\ox 2}] = 0$ (the squaring trick). The expectation value factorizes as $\tr[(P\rho\ox\rho^{\ox(k-1)})\bP^d_{(\mathrm{odds})}] \cdot \tr[(P\rho\ox\rho^{\ox(k-1)})\bP^d_{(\mathrm{even})}] = [\tr(P\rho^k)]^2$.

For $k=2$, the eigenbasis states have the form
$\ket{\Psi^s_{(\mathbf{a},\mathbf{b})}} = 2^{-1/2}[\ket{\Phi_\mathbf{a}}\ket{\Phi_\mathbf{b}} + s\ket{\Phi_\mathbf{b}}\ket{\Phi_\mathbf{a}}]$
for $\mathbf{a} < \mathbf{b}$ with eigenvalues $\lambda = s\cdot 2^{-1}[(-1)^{[\mathbf{r},\mathbf{a}]} + (-1)^{[\mathbf{r},\mathbf{b}]}]$. We realize this measurement with a circuit of transversal CNOTs followed by adaptive computational/GHZ-basis measurements (Appendix~\ref{app:UB_nonlin_P_shad}, Fig.~\ref{fig:4rep_strategy}).

Recovering signed values $\tr(P\rho^k)$ from $|\tr(P\rho^k)|$ requires $\mathcal{O}(\varepsilon^{-4}\log|\cO|)$ additional copies with ancilla, via the signal estimation technique (Appendix, Lemma~\ref{lemma:signal_estim}).

%%%%%%%%%%%%%%%%%%%%%%%%%%%%%%%%%%%%%%%%%%%%%%%%%%%%%%%%%%%%%%%%%%%%%%%%%%%%%%%%%%%%%%%%%%%%%%%%%%%
\prlsubsection{Local no-go.}\label{sec:local_nogo}
For estimating the full nonlinear shadow $\{\tr(P\rho^2)\}_{P\in\cP_n}$, the relevant two-copy operators are $\tilde{O}_P = \tfrac{1}{2}(P\ox\1 + \1\ox P)\cdot\mathrm{SWAP}$, which have eigenvalues $\{+1, 0, -1\}$ on the symmetric subspace---they are \emph{spin-1} (non-binary). By Theorem~\ref{thm:linear-equiv}, the no-advantage result does not apply; estimation may in principle be easier than compatibility. Nevertheless, local strategies remain exponentially costly.

The obstruction is structural: the four-copy operator $A_P = \tfrac{1}{2}(P^{\ox 2}\ox\1^{\ox 2} + \1^{\ox 2}\ox P^{\ox 2})\cdot\pi$ is a \emph{sum} of two product operators, not itself a product. Theorem~\ref{thm:tensor} does not apply, so the per-site factorization $s^*_{\mathrm{loc}} = \prod_l s^*_l$ fails. Explicit computation (Appendix~\ref{app:local_shadow}, Lemma~\ref{lemma:covar_Pauli}) shows that two-copy local 3-design shadows incur second moment $\sim (11/3)^n \cdot (5/11)^{\mathrm{wt}(P)}$ for product states, with operator norm up to $5^n$ for entangled states---exponential in $n$ for any Pauli weight. \color{red}The structural obstruction (non-product form of the lifted operators) suggests that higher local copy numbers cannot circumvent the exponential cost, though a rigorous proof for arbitrary $k$ remains open (see Discussion).\color{black}

This stands in sharp contrast to the nonlocal four-copy strategy of Theorem~\ref{thm:nonlocal}, which achieves $\mathcal{O}(1)$ variance for all $|\tr(P\rho^2)|$. The gap quantifies nonlocality as a measurement resource for multi-copy shadow tomography.

%%%%%%%%%%%%%%%%%%%%%%%%%%%%%%%%%%%%%%%%%%%%%%%%%%%%%%%%%%%%%%%%%%%%%%%%%%%%%%%%%%%%%%%%%%%%%%%%%%%
\prlsection{Discussion.}\label{sec:discussion}
We have introduced the f-compatibility framework, which unifies quantum measurement compatibility~\cite{heinosaari2015noise,designolle2019incompatibility} with multi-copy shadow estimation. The tensor product factorization theorem (Theorem~\ref{thm:tensor}) explains the $3^w$ scaling of classical shadows as emergent from per-site $\sqrt{3}$ incompatibility, while the squaring trick reveals that nonlinear functionals can have dramatically lower estimation complexity ($s^*=1$ vs.\ $s^*=\sqrt{3}$). The contrast between the efficient nonlocal strategy (Theorem~\ref{thm:nonlocal}) and the local exponential lower bound highlights multi-copy entanglement as a genuine resource.


Several directions remain open. First, determining the optimal copy-complexity trade-off $\min_k k\cdot[s^*(k)]^2$ for general functional families is an algorithmic challenge. Second, proving a rigorous exponential lower bound for local $k$-copy strategies estimating $\{\tr(P\rho^2)\}_P$ at arbitrary $k$ would strengthen the local no-go. Third, the commuting structure of $\{A_P\}$ connects to graph-theoretic approaches~\cite{graph2024nonlinear} and fermionic joint measurements~\cite{mcnulty2024fermionic}, whose interplay with f-compatibility deserves further exploration.

%%%%%%%%%%%%%%%%%%%%%%%%%%%%%%%%%%%%%%%%%%%%%%%%%%%%%%%%%%%%%%%%%%%%%%%%%%%%%%%%%%%%%%%%%%%%%%%%%%%
\prlsection{Acknowledgements.}We thank Kun Wang for helpful discussions.

%%%%%%%%%%%%%%%%%%%%%%%%%%%%%%%%%%%%%%%%%%%%%%%%%%%%%%%%%%%%%%%%%%%%%%%%%%%%%%%%%%%%%%%%%%%%%%%%%%%
% References
%%%%%%%%%%%%%%%%%%%%%%%%%%%%%%%%%%%%%%%%%%%%%%%%%%%%%%%%%%%%%%%%%%%%%%%%%%%%%%%%%%%%%%%%%%%%%%%%%%%

% \input{references.bbl}
\bibliographystyle{apsrev4-2}
\bibliography{references}

%%%%%%%%%%%%%%%%%%%%%%%%%%%%%%%%%%%%%%%%%%%%%%%%%%%%%%%%%%%%%%%%%%%%%%%%%%%%%%%%%%%%%%%%%%%%%%%%%%%
% Appendices
%%%%%%%%%%%%%%%%%%%%%%%%%%%%%%%%%%%%%%%%%%%%%%%%%%%%%%%%%%%%%%%%%%%%%%%%%%%%%%%%%%%%%%%%%%%%%%%%%%%
% \appendix
% \onecolumngrid
% https://tex.stackexchange.com/questions/488302/referencing-an-appendix-while-using-prl-style
\makeatletter
\newcommand{\appendixtitle}[1]{\gdef\@title{#1}}
\newcommand{\appendixauthor}[1]{\gdef\@author{#1}}
\newcommand{\appendixaffiliation}[1]{\gdef\@affiliation{#1}}
\newcommand{\appendixdate}[1]{\gdef\@date{#1}}
\makeatother

% Define the custom title block
\makeatletter%
\newcommand{\appendixmaketitle}{%
\begin{center}%
\vspace{0.4in}%
{\Large \@title \par}%
% {\normalsize
%   \lineskip .5em%
%   \begin{tabular}[t]{c}%
%     \@author
%   \end{tabular}\par}%
% \vspace{5pt}
% {\itshape \@affiliation \par}%
\end{center}%
\par%
}%
\makeatother%

\setcounter{secnumdepth}{2}
\appendix
\widetext
\newpage

\appendixtitle{\bf
Supplemental Material for ``Multi-copy shadow tomography''}
% \appendixauthor{Siyuan Chen, Wei Xie, and Kun Wang} 
% \appendixaffiliation{
% Institute for Quantum Computing, Baidu Research, Beijing 100193, China \\
% Physics Department, Jilin University \\
% School of Computer Science and Technology, University of Science and Technology of China}
\appendixmaketitle
\vspace{0.2in}

%%%%%%%%%%%%%%%%%%%%%%%%%%%%%%%%%%%%%%%%%%%%%%%%%%%%%%%%%%%%%%%%%%%%%%%%%%%%%%%%%%%%%%%%%%%%%%%%%%%
%%%%%%%%%%%%%%%%%%%%%%%%%%%%%%%%%%%%%%%%%%%%%%%%%%%%%%%%%%%%%%%%%%%%%%%%%%%%%%%%%%%%%%%%%%%%%%%%%%%
%%%%%%%%%%%%%%%%%%%%%%%%%%%%%%%%%%%%%%%%%%%%%%%%%%%%%%%%%%%%%%%%%%%%%%%%%%%%%%%%%%%%%%%%%%%%%%%%%%%
The contents of the supplementary material are structured as follows:
\begin{itemize}
\item Section~A: f-Compatibility framework (SDP formulation, relaxation gap, proofs)
\item Section~B: Tensor product structure (factorization proof, per-site analysis, spin-1 operators)
\item Section~C: Algebraic framework for multi-copy estimation (tensor algebra, Fisher information, AJD, SDP relaxation)
\item Section~D: Multi-copy local estimation of fixed $\tr(P\rho^2)$
\item Section~E: Multi-copy local shadows (3-design variance, covariance operators, nonlocal 4-replica strategy)
\item Section~F: Lower bounds for nonlinear Pauli shadow
\end{itemize}

%%%%%%%%%%%%%%%%%%%%%%%%%%%%%%%%%%%%%%%%%%%%%%%%%%%%%%%%%%%%%%%%%%%%%%%%%%%%%%%%%%%%%%%%%%%%%%%%%%%
%%%%%%%%%%%%%%%%%%%%%%%%%%%%%%%%%%%%%%%%%%%%%%%%%%%%%%%%%%%%%%%%%%%%%%%%%%%%%%%%%%%%%%%%%%%%%%%%%%%
%%%%%%%%%%%%%%%%%%%%%%%%%%%%%%%%%%%%%%%%%%%%%%%%%%%%%%%%%%%%%%%%%%%%%%%%%%%%%%%%%%%%%%%%%%%%%%%%%%%
\section{f-Compatibility framework}\label{app:f-compat}

In this section we provide the full SDP formulation for f-compatibility and key proofs omitted from the main text.

\subsection{Review of compatibility}

In this section, we briefly review previous result about compatibility and its related results.
Given a quantum state $\rho$, different measurement strategies $x \in \mathcal{X}$ may lead to different distributions $p(a|x)$ of the measurement outcomes. 
The distribution of the measurement outcomes is given by $p(a|x) = \tr(E_{a|x}\rho)$.
The ordinary definition of compatibility is meant to extract 
\begin{definition}
    A set of POVMs $\{E_{a|x}\}$ is jointly measurable if there exists a parent POVM $\{G_\lambda\}$ such that $E_{a|x} = \sum_\lambda p(a|x,\lambda)\,G_\lambda$.
\end{definition}

For example, the Pauli measurements on a qubit can be represented as a set of POVMs $\{E_{a|x}\}$ with $E_{a|x = X} = (1/2)\cdot (\mathbb{I} \pm  X)$, $E_{a|x = Y} = (1/2)\cdot (\mathbb{I} \pm  Y )$  and $E_{a|x = Z} = (1/2)\cdot (\mathbb{I}\pm Z)$, where $a\in \{+,-\}$ and $x\in \{X, Y, Z\}$. 
It has been shown that they are the sharpest Pauli measurements and not jointly measurable. But the unsharp Pauli measurements are jointly measurable. 
$$
E^{\eta_{P}}_{a|x = P} = \eta_P \frac{\mathbb{I} \pm P}{2} + (1 - \eta_P) \frac{\mathbb{I}}{2} =  \frac{\mathbb{I} \pm P\eta_P}{2}, P \in \{X, Y, Z\}
$$
The mixing of $\mathbb{I}$ can be interpreted as $\eta_P$ fractional random guessing or a noise on Pauli measurement $P$.
As shown in~\cite{busch1986unsharp}, the unsharp Pauli measurements are jointly measurable if and only if $(\eta_x)^2 + (\eta_y)^2 +(\eta_z)^2 \leq 1$.

\begin{theorem}[Unsharp Pauli Measurements Are Jointly Measurable]
Define the noisy unsharp versions of Pauli measurements as $E^{\eta_x}_{a|x = X} = (1/2)\cdot (\mathbb{I} + a \eta_x X)$, $E^{\eta_y}_{a|x = Y} = (1/2)\cdot (\mathbb{I} + a \eta_y Y )$  and $E^{\eta_z}_{a|x = Z} = (1/2)\cdot (\mathbb{I}+a\eta_z Z)$, with outcomes $a\in \{+,-\}$  and $0 \leq \eta_x, \eta_y, \eta_z \leq 1$.
The triple is jointly measurable if and only if $(\eta_x)^2 + (\eta_y)^2 +(\eta_z)^2 \leq 1$. The corresponding joint measurement, whose marginals are the unsharp Pauli observables, is  $G(x, y, z) = \frac{1}{8}(\mathbb{I} + x\eta_x X + y\eta_y Y + z\eta_z Z)$.
\end{theorem}

\subsubsection{Robustness of Incompatibility}
It is unnecessary to always use $\mathbb{I}$ mixedness to cancel the non-compatibility. In \cite{skrzypczyk2019all}, the authors use arbitrary $N_{a|x}$ mixed with $M_{a|x}$ to cancel the non-compatibility, that is, 
\begin{align}
    r N_{a|x} + M_{a|x} &= (1 + r)\sum_{\lambda} p(a|x,\lambda) G_\lambda, \forall a,x.\\
     &\Rightarrow (1 + r)\sum_{\lambda} p(a|x,\lambda) G_\lambda \geq M_{a|x} \quad \forall a,x.
\end{align}
This induces the definition below:

\begin{definition}[Genralized Robustness of compatiablity]\label{def:generalized_robustness}
    Given a set of POVMs $\{M_{a|x}\}_{x}$, its generalized robustness of compatiability is defined as the minimization over $r \in [0,\infty]$, the joint POVMs $ \{G_{\lambda}\}_{\lambda}$, and post-processsing distribution $p(a|x,\lambda)$, such that condition below is satisfied:
\begin{align}
    (1 + r)\sum_{\lambda} p(a|x,\lambda) G_\lambda \geq M_{a|x} \quad \forall a,x.
\end{align}
\end{definition}

In our case, $x = P,Q$ and $a = +1,-1$. Thus $G_{\lambda}$ can be coarse-grained into $G_{00}, G_{01}, G_{10}, G_{11}$ as: $G_{j,k} = \sum_\lambda p(a_P=j, a_Q=k | \lambda) G_\lambda$, where,
\begin{align*}
p(a|P, \lambda) = \sum_{a_Q \in \{0, 1\}} p(a_P = a, a_Q | \lambda),  \quad p(a|Q, \lambda) = \sum_{a_P \in \{0, 1\}} p(a_P, a_Q = a | \lambda).
\end{align*}
Then
\begin{align*}
\sum_\lambda p(a|P, \lambda) G_\lambda &= \sum_\lambda \left( \sum_{a_Q} p(a, a_Q | \lambda) \right) G_\lambda \\
&= \sum_{a_Q} \left( \sum_\lambda p(a, a_Q | \lambda) G_\lambda \right) = \sum_{a_Q} G_{a, a_Q}.
\end{align*}

By denoting the super-normalization factor as $s = 1+r$, where $r$ is the generalized compatibility robustness. Defining the super-normalized POVM elements $\tilde{G}_{ij} = s G_{ij}$, we can completely eliminate $\lambda$ and formulate the exact Primal SDP for the Robustness of Incompatibility:
$$
\boxed{
\begin{aligned}
\text{Minimize:} \quad & s \\
\text{Subject to:} \quad & \sum_{k = 0}^{1}\tilde{G}_{j,k} \geq M_{j} \quad \forall j \in \{0,1\} \\
& \sum_{j = 0}^{1}\tilde{G}_{j,k} \geq N_{k} \quad \forall k \in \{0,1\} \\
& \sum_{i,j \in \{0,1\}} \tilde{G}_{ij} = s\mathbb{I} \\
& \tilde{G}_{ij} \geq 0 \quad \forall i,j \in \{0,1\}
\end{aligned}
}
$$
Once the SDP is solved for $s^*$, the minimum noise robustness is directly obtained as $r^* = s^* - 1$.

\begin{lemma}[Dual Discrimination~\cite{skrzypczyk2019all}]
By finding the dual problem, the generalized compatibility robustness of $\{P,Q\}$ is equivalent to the problem of finding the best mixed ensemble $\mathcal{E} = \mathcal{E}_P = \{\rho_0,\rho_1\} \cup \mathcal{E}_Q := \{\sigma_0,\sigma_1\}$ such that case (i) has $1 +s^*$ times better success rate than case (ii):
\begin{itemize}
    \item  Given $P,Q$ specified, using the $\{M_0, M_1\}$ or $\{N_0, N_1\}$ measurements to discriminate the index $0$ or $1$.
    \item  Without getting the information that the state is coming from $\mathcal{E}_P$ or $\mathcal{E}_Q$, distinguish index $0$ or $1$.
\end{itemize}
\end{lemma}
\color{red}
[TODO: add new lemma]\
\begin{lemma}[Discrimination reduction]
Given $\{P,Q\}$ specified and $\mathcal{E} = \mathcal{E}_P = \{\rho_0,\rho_1\} \cup \mathcal{E}_Q := \{\sigma_0,\sigma_1\}$  solved as the optimal ensemble that achieves the best discrimination success rate, then the problem of simultaneously estimation of $\{P,Q\}$ can be reduced to discriminating $\{P,Q\}$ with $\{M_0, M_1\}$ or $\{N_0, N_1\}$.
\end{lemma}
\color{black}

\subsubsection{Multicopy Compatibility}

In \cite{carmeli2016quantum}, the authors define the multicopy compatibility as follows:
Given $K$ POVMs $\{E_{a|x}\}_{x = 1}^K, a \in \Omega_x$ on outcome sets $\{\Omega_1,\ldots,\Omega_K\}$, we ask the following: given $k$ copies of a quantum state $\rho^{\otimes k}$, can we design a single POVM $G_{\lambda}$ with outcomes $(x_1,\ldots,x_K) \in \Omega_1\times\cdots\times\Omega_K$ on $\rho^{\otimes k}$ such that, for each $i$, the marginal statistics of $G$ on $\Omega_i$ reproduce the measurement statistics of $A_i$ on $\rho$?
In the case $k \geq K$, this is always achievable: implementing $A_1\otimes \cdots \otimes A_K \otimes \mathbb{I}^{\otimes (N-K)}$ on $\rho^{\otimes N}$ measures each $A_i$ on a separate copy, and the marginals correctly reproduce the individual statistics.
In the case $k < K$, the problem becomes non-trivial.

\begin{definition}[Multicopy Compatibility]\label{def:multicopy-compatibility}
Given $k$ POVMs $\{E_{a|x}\}_{x = 1}^K$ on outcome sets $\{\Omega_i\}_{i=1}^K$, respectively.
They are $k$-replica compatible if there exists an POVMs $G_{\lambda}$ with outcomes $\lambda = (x_1,\ldots,x_K) \in \Omega_1\times\cdots\times\Omega_K$ on the product state $\rho^{\otimes k}$, such that
\begin{align}
\sum_{\lambda/x_i} \operatorname{tr}\!\left[\rho^{\otimes k}\, G_\lambda \right] = \operatorname{tr}\!\left[\rho\, E_{a|x_i}\right],\forall a \in \Omega_i.
\end{align}
\end{definition}

For any operator $O$ on $\cH$, we define the \emph{symmetric embedding} into $\cH^{\ox k}$ by
\begin{align}\label{eq:sym-embed}
\1^{\ox(k-1)}\odot O := \frac{1}{k}\sum_{i=1}^{k} \1^{\ox(i-1)}\ox O \ox \1^{\ox(k-i)},
\end{align}
which satisfies $\tr[(\1^{\ox(k-1)}\odot O)\,\rho^{\ox k}] = \tr(O\rho)$ for all $\rho$.

\begin{lemma}[Structure Lemma]\label{lem:structure-lemma}
The POVMs $\{E_{a|x}\}_{x = 1}^K$ are $k$-replica compatible if and only if
$\tilde{E}_{a|x} := \mathbb{I}^{\otimes(k-1)}\odot E_{a|x} $ is compatible.
\end{lemma}

\begin{proof}
Let $\Sigma_k(B) = (k!)^{-1}\sum_{\pi \in S_k} P_\pi B P_\pi^\dagger$ denote the symmetrization map on $\mathcal{L}(\mathcal{H}^{\otimes k})$. Since $\rho^{\otimes k}$ is permutation-symmetric, $\tr[\rho^{\otimes k}\, B] = \tr[\rho^{\otimes k}\, \Sigma_k(B)]$ for all $\rho$.

($\Leftarrow$) Suppose $\{\tilde{E}_{a|x}\}$ are compatible with parent POVM $\{G_\lambda\}$ on $\mathcal{H}^{\otimes k}$ and conditional distributions $p(a|x,\lambda)$ satisfying $\sum_\lambda p(a|x,\lambda)\, G_\lambda = \tilde{E}_{a|x}$. Then $\sum_\lambda p(a|x,\lambda)\, \tr[\rho^{\otimes k}\, G_\lambda] = \tr[\rho^{\otimes k}\, \tilde{E}_{a|x}] = \tr[\rho\, E_{a|x}]$ for all $\rho$, so $\{G_\lambda\}$ realizes $k$-replica compatibility.

($\Rightarrow$) Suppose $\{E_{a|x}\}$ are $k$-replica compatible via a POVM $\{G_\lambda\}$ on $\mathcal{H}^{\otimes k}$, with marginals $M_{a|x} := \sum_{\lambda: x\text{-th outcome}=a} G_\lambda$ satisfying $\tr[\rho^{\otimes k}\, M_{a|x}] = \tr[\rho\, E_{a|x}]$ for all $\rho$.
Symmetrize: $\tilde{G}_\lambda := \Sigma_k(G_\lambda)$. This preserves positivity and normalization, and $\tr[\rho^{\otimes k}\, \tilde{G}_\lambda] = \tr[\rho^{\otimes k}\, G_\lambda]$. The symmetrized marginals $\tilde{M}_{a|x} = \Sigma_k(M_{a|x})$ satisfy $\tr[\rho^{\otimes k}\, \tilde{M}_{a|x}] = \tr[\rho\, E_{a|x}] = \tr[\rho^{\otimes k}\, \tilde{E}_{a|x}]$ for all $\rho$. Since products $\rho^{\otimes k}$ span the symmetric subspace and $\tilde{M}_{a|x}$, $\tilde{E}_{a|x}$ are both permutation-symmetric, we conclude $\tilde{M}_{a|x} = \tilde{E}_{a|x}$ on the symmetric subspace. (The identification extends to all of $\mathcal{H}^{\otimes k}$ by a polynomial expansion argument: writing $\rho = \mathbb{I}/d + t\Delta$ and matching coefficients of $t^j$ for all traceless $\Delta$ shows that only the $j=0$ and $j=1$ components contribute, forcing $\tilde{M}_{a|x} = \mathbb{I}^{\otimes(k-1)}\odot E_{a|x}$.) Thus $\{\tilde{G}_\lambda\}$ is a parent POVM for $\{\tilde{E}_{a|x}\}$, establishing standard compatibility.
\end{proof}

\subsection{The definition of f-compatibility}


Given $K$ unknown discrete probability distributions $(\vec{p}_1, \vec{p}_2, \dots, \vec{p}_K)$ on finite alphabets of sizes $|\Omega_1|, |\Omega_2|, \dots, |\Omega_K|$. Suppose we draw $(n_1, n_2, \dots, n_K)$ i.i.d. samples independently from each distribution with $\sum_{i=1}^K n_i = k$, yielding frequency vectors $(N_{i,1}, \ldots, N_{i,|\Omega_i|})$ for each $i$, where $\sum_j N_{i,j} = n_i$. 
The concept of nonparametric statistics not interested in directly estimating the distributions $(\vec{p}_1, \vec{p}_2, \dots, \vec{p}_K)$, but rather a function of those distributions. This includes the estimation of Shannon entropy $\sum_{i=1}^K \sum_{j \in \Omega_i} p_{i,j} \log p_{i,j}$, the collision probability $\sum_{i=1}^K \sum_{j=1}^K p_{i,j} p_{j,i}$, the variance of the distribution $\sum_{i=1}^K \sum_{j \in \Omega_i} (p_{i,j} - \sum_{j \in \Omega_i} p_{i,j})^2$, the support size $\sum_{i=1}^K \mathbf{1}(p_{i,j} > 0)$, etc.

A monomial in these probability variables takes the form $m_{\mathbf{x}} = \prod_{i=1}^K \prod_{j \in \Omega_i} p_{i,j}^{x_{i,j}}$, where $\mathbf{x} = (x_{i,j})$ is a multi-index of non-negative integers. The degree with respect to distribution $i$ is $r_i(\mathbf{x}) = \sum_{j \in \Omega_i} x_{i,j}$, and we defined the maximum degree as $r_{m}:=$ $\max_{i} r_i(\mathbf{x})$.
A multivariate polynomial of degree at most $r$ in $(\vec{p}_1, \vec{p}_2, \dots, \vec{p}_K)$ is a function of the form $F(\vec{p}_1, \ldots, \vec{p}_K) = \sum_{\mathbf{x} \in \mathcal{X}_r} c_{\mathbf{x}}\, m_{\mathbf{x}}$, where $\mathcal{X}_r$ is the set of all multi-indices $\mathbf{x}$ with $|\mathbf{x}| \leq r$ and $c_{\mathbf{x}}$ are constants.
It has been shown that given $r$ fixed~\cite{wu2020polynomial}, only functionals that can be expressed as a multivariate polynomial of degree at most $k$ in $(\vec{p}_1, \vec{p}_2, \dots, \vec{p}_K)$ have unbiased estimators based on the observation $\boldsymbol{N}$.


\begin{lemma}[Classical Estimation Condition~\cite{wu2020polynomial}]
A functional $F(\vec{p}_1, \ldots, \vec{p}_K)$ admits an unbiased estimator from $(n_1, \ldots, n_K)$ i.i.d. samples if and only if $F$ is a polynomial of degree array $r = (r_1, \ldots, r_K)$ in $(\vec{p}_1, \vec{p}_2, \dots, \vec{p}_K)$ that satisfies the condition $r_i \leq n_i$.
\end{lemma}

\begin{proof}
Any estimator $\hat{F}$ is a function of the frequency vectors $(N_{i,j})$. Each frequency vector follows a multinomial distribution, so:
\begin{align}
\mathbb{E}_P[\hat{F}] = \sum_{\{N_{i,j}\}} \hat{F}(N_{i,j}) \prod_{i=1}^K \frac{n_i!}{\prod_j N_{i,j}!} \prod_j p_{i,j}^{N_{i,j}}.
\end{align}
Here summation is over all possible frequency vectors $(N_{i,j})$.
The right-hand side is manifestly a polynomial in $\{p_{i,j}\}$ of degree at most $n_i$ in the variables of distribution $i$. 
Therefore, for $\mathbb{E}[\hat{F}] = F$ to hold for all valid probability vectors, $F$ must itself be such a polynomial.

Conversely, for a monomial $m_{\mathbf{x}} = \prod_{i,j} p_{i,j}^{x_{i,j}}$ with $\sum_j x_{i,j} = n_i$, to build an unbiased estimator, we first build a product of i.d.d. ditributions of $n_i$ numbers of random variables distributed according to $\{p_{i,j}\}_{j \in \Omega_i}$.
This forms a distribution on space 
$$
\{s^{(r)}_{j}\}:= (s^{(1)}_{1}, \ldots, s^{(1)}_{n_{1}}) \times  \cdots \times (s^{(K)}_{1}, \ldots, s^{(K)}_{n_{K}}) \in \Omega_1^{\times n_1} \times \cdots \times \Omega_K^{\times n_K}.
$$
Then for a monomial $m_{\mathbf{x}} = \prod_{i} \prod_{j \in \Omega_i} p_{i,j}^{x_{i,j}}$ with $\sum_j x_{i,j} = n_i$, we can build an unbiased estimator by setting up an function $\Omega_1^{n_1} \times \cdots \times \Omega_K^{n_K} \to \{0,1\}$ that 
\begin{align}
\mathbb{I}(\{s^{(r)}_{j}\}) = \begin{cases}
1 & \text{if } s^{(r)}_{1} = \cdots = s^{(r)}_{n_r}, \\
0 & \text{otherwise},
\end{cases}
\end{align}
and conclude that the estimator is unbiased: $\mathbb{E}_{\{s^{(r)}_{j}\}}\!\left[\mathbb{I}(\{s^{(r)}_{j}\})\right] = m_{\mathbf{x}}$.
\end{proof}


In the quantum setting, the $K$ distribution vectors are generated from the quantum state $\rho$ and different POVMs $\{A_1,\ldots,A_K\}$ on the outcome sets $\{\Omega_1,\ldots,\Omega_K\}$, i.e., $\vec{p}_i(\rho) = (\mathrm{tr}[A_i^{(r)}\rho])_{r \in \Omega_i}$.
The functional set $\mathcal{F}:=\{f_\alpha\}_{\alpha \in \mathcal{A}}$ can then be replaced by a set of matrix valued functions $\{f_\alpha\}_{\alpha \in \mathcal{A}}$ such that $f_\alpha(\rho)$.
We then define the $k$-copy f-compatibility as follows:

\begin{definition}[Multicopy f-compatibility]\label{def:f-compatibility}
 We say that a set of state properties $\mathcal{F}$ is $k$-copy compatible if there exists a POVM $\{\Pi_j\}_{j \in \mathcal{J}}$ on $\mathcal{H}^{\otimes k}$ and classical functions $\{h_\alpha: \mathcal{J} \to [-1,1]\}$ such that
$$
f_\alpha(\rho) = \sum_{j \in \mathcal{J}} h_\alpha(j)\;\mathrm{tr}[\Pi_j\;\rho^{\otimes k}], \quad \forall\,\rho,\;\forall\,\alpha
$$
We denote $k(\mathcal{F})$ as the smallest copy number $k$ such that $\mathcal{F}$ is $k$-copy compatible.
\end{definition}

The condition $h_\alpha(j) \in [-1,1]$ ensures that the random variable is bounded. 
By Hoeffding's inequality, by implementing the POVM $\{\Pi_j\}_{j \in \mathcal{J}}$ on $\mathcal{H}^{\otimes k}$ on the state $\rho^{\otimes k}$ for $N$ times and generate $\{j_1,\cdots,j_N\}$, then
\begin{align}
\mathrm{Pr}\left[\left|N^{-1}\sum_{i=1}^N h_\alpha(j_i) - f_{\alpha}(\rho)\right| \geq \varepsilon\right] \leq 2\exp\left(-\frac{\varepsilon^2 N}{2}\right).
\end{align}
Considering all $\alpha \in \mathcal{A}$, the probability that the random variable deviates from the true value by more than $\varepsilon$ is at most $2\exp\left(-\frac{\varepsilon^2 N}{2}\right)|\mathcal{A}|$. This means that $N \sim \frac{2}{\varepsilon^2} \log(2|\mathcal{A}|\delta^{-1})$ is the sufficient number of copies to estimate all the functionals in $\mathcal{F}$ with accuracy $\varepsilon$ and confidence $1-\delta$.
It is also a generalization of the Definition~\ref{def:multicopy-compatibility} of $k$-compatible, which is a special case when $f_{(m,r)}(\rho) = \mathrm{tr}[A_m^{(r)}\rho], \alpha = (m,r) \in \{(m,r) : m \in \{1,\cdots,K\},\, r \in \Omega_m\}$.

A great advantage of this definition is that the compatible learning of multiple fixed properties of quantum states might be easier than the compatible learning of the whole POVMs distributions. 
Here below we provide a description of $\mathcal{F}$ where this simplification holds.

\color{red}
[todo: Does ordinary k-compatibility can upper bound f-compatibility?]
\begin{lemma}[Upper bound of Multicopy f-compatibility]\label{lem:f-lowerbound}
For any set of POVMs $\{A_i\}_{i=1}^K$ that is $k$-compatible and any polynomial functional set $\mathcal{F}$ of maximum degree at most $r$ and evaluates to bounded values in $[-1,1]$. then $k(\mathcal{F}) \leq r_{m}k$. Specifically, the inequality reaches the tight bound in $r = 1$ case when $f_{\alpha}(\vec{p}_1(\rho),\ldots,\vec{p}_K(\rho)) = p_i(a_i) = \mathrm{tr}[A_i^{(a_i)} \rho]$, $\alpha = (i,a_i) \in \{(i,a_i) : i \in \{1,\cdots,K\},\, a_i \in \Omega_i\}$.
\end{lemma}
\begin{proof}

    First we show that $k$-compatibility implies $\mathcal{F}$-compatibility ($k(\mathcal{F}) \leq k$).
    Assume the set of POVMs $\{A_i\}_{i=1}^K$ is $k$-compatible. By Lemma~\ref{lem:structure-lemma}, there exists a joint mother POVM $G(\vec{x}) = G(x_1, \ldots, x_K)$ acting on $\mathcal{H}^{\otimes k}$ such that measuring $G(\vec{x})$ on $\rho^{\otimes k}$ yields a joint classical probability distribution $P(\vec{x}|\rho) = \mathrm{tr}[G(\vec{x})\rho^{\otimes k}]$, from which marginal distributions are recovered: $\sum_{\vec{x} \setminus x_i} P(\vec{x}|\rho) = p_i(x_i|\rho)$ for $x_i \in \Omega_i$.
    
    Now suppose $\mathcal{F} = \{f_\alpha\}_{\alpha}$ consists of degree-$\leq k$ polynomial functionals that admit bounded classical estimators from the joint outcome distribution. That is, there exist classical functions $C_\alpha(\vec{x}) \in [-1, 1]$ such that:
    $$
    f_\alpha(\vec{p}_1, \dots, \vec{p}_K) = \sum_{\vec{x}} C_\alpha(\vec{x})\, P(\vec{x}|\rho) = \sum_{\vec{x}} C_\alpha(\vec{x})\, \mathrm{tr}[G(\vec{x})\,\rho^{\otimes k}].
    $$
    Setting $\Pi_{\vec{x}} := G(\vec{x})$ and $h_\alpha(\vec{x}) := C_\alpha(\vec{x})$ satisfies all conditions of Definition~\ref{def:f-compatibility}, yielding $k(\mathcal{F}) \leq k$.
    
\end{proof}
\color{black}

Another critical question is, whether the f-compatiablity indeed replects the inherient hardness to simultanous estimation a set of properties $\{f_1(\rho),\cdots,f_\alpha(\rho)\}$.
Given two set of function of $\rho$, $\mathcal{S}:= (S_\alpha(\rho))_{\alpha \in \mathcal{A}}$ and $\mathcal{T}:=(T_\beta(\rho))_{\beta \in \mathcal{B}}$, If there exists a function $f: \mathcal{A} \to \mathcal{B}$ such that $T_\beta(\rho) = f(\S_\alpha(\rho))$, then the estimation complexity of $\{T_\beta(\rho)\}_{\beta \in \mathcal{B}}$ is the same as that of $\{S_\alpha(\rho)\}_{\alpha \in \mathcal{A}}$.


\subsection{Application on linear functionals}
For Hermitian operators $O_i$ with bounded eigenvales, $\|O_i\|_{\infty} \leq 1$, it splits the Hilbert space into two eigenspaces with POVMs as $\{M_i^{(+)}:= \frac{\mathbb{I} + O_i}{2},M_i^{(-)}:= \frac{\mathbb{I} - O_i}{2}\}$.
An familiar example is the Pauli compatibility problem, given a set of Pauli observables $\{P_i\}_i$, getting its average is as hard as compatible realization its 2-outcome POVMs $\frac{\mathbb{I} \pm P_i}{2}$.
We then noted that the f-compatibility problem of the set of properties $\mathcal{F}:= \{\mathrm{tr}[O_i\rho]\}_i$ is equivalent to the standard joint measurability problem of the POVMs $\{M_i^{(+)},M_i^{(-)}\}$.
 
\begin{lemma}[Equivalence for bounded operators]\label{thm:linear_functional_compatibility}
   Given a set of operators $\{O_i\}_i$ with $\|O_i\|_\infty \leq 1$, and set up $\mathcal{F}:= \{\mathrm{tr}[O_i\rho]\}_i$, then the set of linear functional $\mathcal{F}$ is $k$-copy compatible if and only if the POVMs $\{M_i^{\pm}:= \frac{\mathbb{I} \pm O_i}{2}\}_i$ is $k$-copy compatible.
\end{lemma}

\begin{proof}
Suppose $\{M_i^+, M_i^-\}_{i=1}^K$ are jointly measurable. By definition there exists a parent POVM $\{G_\lambda\}_\lambda$ that supports on $\mathcal{H}^{\otimes k}$ and a conditional distributions $p(a|i,\lambda)$ with $a \in \{+1,-1\}$ such that:
\begin{align}
\mathbb{I}^{\otimes(k-1)}\odot \frac{\mathbb{I} + a\, O_i}{2} = \sum_\lambda p(a|i,\lambda)\, G_\lambda, \quad \forall\, a, i.
\end{align}
Taking the difference $a=+1$ minus $a=-1$:
\begin{align}
    \mathbb{I}^{\otimes(k-1)}\odot  O_i = \sum_\lambda \bigl[p(+1|i,\lambda) - p(-1|i,\lambda)\bigr]\, G_\lambda =: \sum_\lambda h_i(\lambda)\, G_\lambda.
\end{align}
Use the fact that $\mathrm{tr}\left[\left(\mathbb{I}^{\otimes(k-1)}\odot O_i\right) \rho^{\otimes k} \right]= \mathrm{tr}(O \rho)$ and $p(+1|i,\lambda) + p(-1|i,\lambda) = 1$ both are non-negative, $|h_i(\lambda)| \leq 1$.
 Therefore $\mathrm{tr}[O_i\rho] = \sum_\lambda h_i(\lambda)\,\mathrm{tr}[G_\lambda\rho]$ for all $\rho$, which is exactly f-compatibility with bounded estimators.

Inversely, suppose there exists a POVM $\{\Pi_j\}_j$ and estimators $\{h_i(j)\}$ with $|h_i(j)| \leq 1$ such that
\begin{align}
     \mathrm{tr}\left[ \sum_j h_i(j)\,\Pi_j \rho^{\otimes k}\right] = \mathrm{tr}\left[ O_i \rho \right] = \mathrm{tr}\left[ \mathbb{I}^{\otimes(k-1)}\odot  O_i \rho^{\otimes k} \right] .
\end{align}
Since this is satisfied for general $\rho$, we conclude that the symmetric version of $\mathrm{sym}\left[\Pi_j\right] := (k!)^{-1}\sum_{\pi \in S_k} S_{k}\Pi_i S_{k}^\dagger$ such that $\sum_j h_i(j)\,\mathrm{sym}\left[\Pi_j\right] =\mathbb{I}^{\otimes(k-1)}\odot O_i$ for all $i$. 
Define conditional probabilities:
$$
p(\pm 1|i,j) := \frac{1 \pm h_i(j)}{2}.
$$
Since $|h_i(j)|\leq 1$, both $p(+1|i,j) \geq 0$ and $p(-1|i,j) \geq 0$, and they sum to 1. Now verify the marginals:

\begin{align}
\sum_j p(\pm 1|i,j)\,\mathrm{sym}\left[\Pi_j\right] &= \frac{1}{2}\sum_j \mathrm{sym}\left[\Pi_j\right] \pm \frac{1}{2}\sum_j h_i(j)\,\mathrm{sym}\left[\Pi_j\right]  = \frac{\mathbb{I} \pm \mathbb{I}^{\otimes(k-1)}\odot O_i}{2}\\
&=\mathbb{I}^{\otimes(k-1)}\odot \frac{\mathbb{I} + a\, O_i}{2}.
\end{align}

This is precisely a joint measurement realization of $\{\mathbb{I}^{\otimes(k-1)}\odot \frac{\mathbb{I} \pm \, O_i}{2}\}_{i=1}^K$ with parent POVM $\{\mathrm{sym}\left[\Pi_j\right]\}$, thus implies the $k$-replica compatiabilty of $\{\frac{\mathbb{I} \pm \, O_i}{2}\}_i$ based on Lemma~\ref{lem:structure-lemma}.
\end{proof}
This also implies that a set of linear property $\{\mathrm{tr}\left(\rho O_i\right)\}_{i = 1}^K$ is $k$-replica f-compatiable for $\rho$ if and only if its symmetric $\{\mathbb{I}^{\otimes (k-1)} \odot O_i\}$ is single copy f-compatiable on $\rho^{\otimes k}$.

% \color{red}
% (todo: A generalization to linear combinations of binary POVMs:
%  \begin{lemma}[The linear functional compatibility problem]\label{thm:general_linear_functional_compatibility}
% Given $K$ 2-outcome POVMs $\{A_i^{(+)},A_i^{(-)}\}$ and a set of $K$ functionals $f_1, \ldots, f_K$ that are linearly independent combinations of $2K$-variables $\{\mathrm{tr}[A_i^{(+)}\rho],\mathrm{tr}[A_i^{(-)}\rho]\}_{i=1}^K$, then $k(\mathcal{F}) = k$. That is, the f-compatibility problem is as hard as the standard joint measurability problem of the POVMs $\{A_i^{(+)},A_i^{(-)}\}$.
% \end{lemma}
% )
% \color{black}

\subsection{Application on non-linear functionals}
The functional set $\mathcal{F}$ are in general not linear in $\rho$, but includes many non-linear polynomial functions such as $|\mathrm{tr}[P_i\rho]|^2$ and $ \mathrm{tr}[P_i\rho^2]$.
For a set of rank-$r$ properties, we define the $r$-replica linearization of $\mathcal{F}$ as a set of operators $\{F_\alpha\}_{\alpha \in \mathcal{A}}$ such that the non-linear functions can be quantified in linear way.

\begin{definition}[Linearization of degree-$r$ polynomials]\label{def:linearization}
    Given a set of degree-$r$ non-linear functions $\{f_\alpha(\rho)\}_{\alpha \in \mathcal{A}}$, its $l$-linearization (for $l \geq r$) is defined as a set of Hermitian operators $\{F_\alpha\}_{\alpha \in \mathcal{A}}$ on $\cH^{\ox l}$ such that $\mathrm{tr}(F_{\alpha}\, \rho^{\otimes l}) = f_{\alpha}(\rho)$ for all states $\rho$.
\end{definition}
An $l \geq r$ linearization of a degree-$r$ function of $\rho$ can always be done since its components satisfy $\left[\mathrm{tr}(O \rho^{m})\right]^{p} = \mathrm{tr}([(O\otimes \mathbb{I}^{m-1})\cdot S_{m}]^{\otimes p} \cdot \rho^{\otimes m p})$.
 Thus we can always set $F = \mathbb{I}^{\otimes (l-mp)} \otimes [(O\otimes \mathbb{I}^{m-1}) \cdot S_{m}]^{\otimes p}$ or its symmetric over $l$ copies of state $\rho^{\otimes l}$.

\begin{lemma}[Linearization lemma]\label{lem:structure_fcomp_lemma}
Given a set of rank-$r$ funciton $\mathcal{F}$, it is $l$-replica compatiable if and only if its $l$-linearization is possible and is single copy f-compatiable on $\rho^{\otimes l}$.
\end{lemma}
\begin{proof}
    \color{red}
    (todo: complete the proof)
    \color{black}
\end{proof}

\subsection{f-Compatibility SDP: primal and dual}

Since now $k$-replica f-compatiablity on a set of rank-$r$ properties can be reduced to a single-replica f-compatiablity problem, in this section, we reduced the single-replica f-compatiablity problem as an SDP problem. 
The f-compatibility robustness $s^*(\cF)$ is computed by the following SDP. The optimization variables are the POVM elements $\{\Pi_j\}_{j\in\cJ}$, the estimator scalars $\{h_\alpha(j)\}$, and the super-normalization factor $s$.
\begin{equation}\label{eq:primal-SDP}
\boxed{
\begin{aligned}
\text{Minimize:} \quad & s \\
\text{Subject to:} \quad & \sum_{j} h_\alpha(j)\,\Pi_j = F_\alpha, \quad \forall\,\alpha \in \cA \\
& \sum_j \Pi_j = s\,\1^{\ox k},\\
& \Pi_j \geq 0, \quad |h_\alpha(j)| \leq 1.\\
\end{aligned}
}
\end{equation}


The constraint $|h_\alpha(j)| \leq 1$ can be rewritten as $\Pi_j \pm h_\alpha(j)\Pi_j \geq 0$, or equivalently by introducing matrix variables $\Gamma_{\alpha,j} := h_\alpha(j)\Pi_j$ subject to $\Pi_j \pm \Gamma_{\alpha,j} \geq 0$.
\color{red}
(todo: fix the relaxation gap. Refer to `_resource/JMvsQSD_main2.tex`''s proof of dual problem. Here my case should be a special case of their proof.)
\color{black}
\emph{Relaxation gap.} When $\Pi_j$ has rank $> 1$, replacing the scalar constraint $h_\alpha(j)\Pi_j$ by a free matrix $\Gamma_{\alpha,j}$ with $\Pi_j \pm \Gamma_{\alpha,j} \geq 0$ is a \emph{relaxation}: $\Gamma_{\alpha,j}$ need not equal a scalar times $\Pi_j$. The relaxed SDP is exact for rank-1 (projective) POVMs but can underestimate $s^*$ otherwise. For example, for $\{X,Y,Z\}$ on $\mathbb{C}^2$, the relaxed SDP gives $s^*_{\mathrm{relax}} = 1$ while the true scalar-constrained answer is $s^* = \sqrt{3}$.

\paragraph{Dual problem.}
The dual of the primal SDP~\eqref{eq:primal-SDP} reads:
\begin{align}
s^* = \max\; & \sum_\alpha \tr[\Lambda_\alpha F_\alpha] \notag\\
\text{s.t.}\; & \sum_\alpha |\Lambda_\alpha| \leq W, \quad \tr[W] = 1, \quad W \geq 0, \label{eq:dual-SDP}
\end{align}
where $|\Lambda_\alpha| = \sqrt{\Lambda_\alpha^\dagger \Lambda_\alpha}$ denotes the operator absolute value. The dual variable $W$ acts as a ``worst-case'' state that maximizes the total estimation demand, normalized to unit trace.

\subsection{Examples}
\paragraph{$\mathcal{F} = \{|\tr(X\rho)|^2, |\tr(Y\rho)|^2, |\tr(Z\rho)|^2\}$.}
We show that for the squared Pauli expectations $f_i(\rho) = |\tr(P_i\rho)|^2$ with $\{P_i\} = \{X,Y,Z\}$ on a single qubit, Bell measurement achieves $s^* = 1$.
The two-copy linearizations are $F_i = P_i \ox P_i$. These operators mutually commute since $[P_i\ox P_i, P_j\ox P_j] = (P_iP_j)\ox(P_iP_j) - (P_jP_i)\ox(P_jP_i) = 0$. Their common eigenbasis is the Bell basis $\{\ket{\Phi_{ab}}\}_{a,b\in\{0,1\}}$ with eigenvalues:
\begin{align}
(P_i \ox P_i)\ket{\Phi_{ab}} = (-1)^{[P_i, (a,b)]}\ket{\Phi_{ab}},
\end{align}
where $[\cdot,\cdot]$ denotes the symplectic inner product. Setting $\Pi_{ab} = \proj{\Phi_{ab}}$ and $h_i(a,b) = (-1)^{[P_i,(a,b)]}$, we have $\sum_{ab} h_i(a,b)\Pi_{ab} = P_i \ox P_i = F_i$ and $\sum_{ab}\Pi_{ab} = \1^{\ox 2}$, so $s = 1$. Since $|h_i(a,b)| = 1 \leq 1$, all constraints are satisfied with $s^* = 1$.

\paragraph{$\mathcal{F} = \{\tr(X\rho), \tr(Y\rho), \tr(Z\rho)\}$.}

\paragraph{$\mathcal{F} = \{ \tr(\rho^2),\tr(X\rho^2), \tr(Y\rho^2), \tr(Z\rho^2)\}$.}
%%%%%%%%%%%%%%%%%%%%%%%%%%%%%%%%%%%%%%%%%%%%%%%%%%%%%%%%%%%%%%%%%%%%%%%%%%%%%%%%%%%%%%%%%%%%%%%%%%%
%%%%%%%%%%%%%%%%%%%%%%%%%%%%%%%%%%%%%%%%%%%%%%%%%%%%%%%%%%%%%%%%%%%%%%%%%%%%%%%%%%%%%%%%%%%%%%%%%%%
\section{Tensor product Extension}\label{app:tensor}
We note that it is impossible to solve the SDP for f-compatibility for a $n$-qubit state $\rho^{\otimes n}$ in polynomial time, since the number of variables scales exponentially with $n$.   
We provide a tensor product extension of the f-compatibility SDP to handle this problem.
To utilize the f-compatibility SDP for a $n$-qubit state $\rho^{\otimes n}$, we proposed that we can tensor product the optimal solution of f-compatibility SDP for a $k$-qubit state and forms a $k$-replica strategy to simultaneously estimate the $n$-qubit obsevables.

\subsection{Proof of Theorem~\ref{thm:tensor} (factorization)}

We prove that for product-form linearizations $F_\alpha = \bigotimes_{l=1}^n F_{\alpha_l}^{(l)}$ with $\alpha = (\alpha_1,\ldots,\alpha_n)$, the local f-compatibility robustness satisfies $s^*_{\mathrm{loc}} = \prod_{l=1}^n s^*_l$.

\emph{Upper bound ($s^*_{\mathrm{loc}} \leq \prod_l s^*_l$).}
Let $\{\Pi_{j_l}^{(l)}\}$ and $\{h_{\alpha_l}^{(l)}(j_l)\}$ be optimal solutions at site $l$ with $\sum_{j_l}\Pi_{j_l}^{(l)} = s^*_l\,\1_{2^k}$. Define the product POVM $\Pi_{\vec{j}} = \bigotimes_l \Pi_{j_l}^{(l)}$ and product estimators $h_\alpha(\vec{j}) = \prod_l h_{\alpha_l}^{(l)}(j_l)$. Then:
\begin{align}
\sum_{\vec{j}} h_\alpha(\vec{j})\Pi_{\vec{j}} &= \bigotimes_{l=1}^n \sum_{j_l} h_{\alpha_l}^{(l)}(j_l)\Pi_{j_l}^{(l)} = \bigotimes_{l=1}^n F_{\alpha_l}^{(l)} = F_\alpha,\\
\sum_{\vec{j}} \Pi_{\vec{j}} &= \bigotimes_{l=1}^n s^*_l\,\1 = \textstyle\left(\prod_l s^*_l\right)\1^{\ox kn}.
\end{align}
Since $|h_\alpha(\vec{j})| = \prod_l |h_{\alpha_l}^{(l)}(j_l)| \leq 1$, the product solution is feasible with $s = \prod_l s^*_l$.

\emph{Lower bound ($s^*_{\mathrm{loc}} \geq \prod_l s^*_l$).}
We use the dual SDP~\eqref{eq:dual-SDP}. Let $\{\Lambda_{\alpha_l}^{*(l)}, W_l^*\}$ be optimal dual variables at site $l$, achieving $s^*_l = \sum_{\alpha_l} \tr[\Lambda_{\alpha_l}^{*(l)} F_{\alpha_l}^{(l)}]$ subject to $\sum_{\alpha_l} |\Lambda_{\alpha_l}^{*(l)}| \leq W_l^*$ and $\tr[W_l^*] = 1$.

Construct product dual variables for the global (unrestricted) SDP: $\Lambda_\alpha := \bigotimes_{l=1}^n \Lambda_{\alpha_l}^{*(l)}$ and $W := \bigotimes_{l=1}^n W_l^*$. These are feasible: $\tr[W] = \prod_l \tr[W_l^*] = 1$, and the constraint holds since $|\bigotimes_l \Lambda_{\alpha_l}| = \bigotimes_l |\Lambda_{\alpha_l}|$ for Hermitian operators (the absolute value of a tensor product is the tensor product of absolute values), giving
\begin{align}
\textstyle\sum_\alpha |\Lambda_\alpha| = \sum_{\alpha_1,\ldots,\alpha_n} \bigotimes_l |\Lambda_{\alpha_l}^{*(l)}| = \bigotimes_l \bigl(\sum_{\alpha_l} |\Lambda_{\alpha_l}^{*(l)}|\bigr) \leq \bigotimes_l W_l^* = W.
\end{align}
The dual objective evaluates to $\sum_\alpha \tr[\Lambda_\alpha F_\alpha] = \prod_l \sum_{\alpha_l} \tr[\Lambda_{\alpha_l}^{*(l)} F_{\alpha_l}^{(l)}] = \prod_l s^*_l$.
Therefore $s^*_{\mathrm{global}} \geq \prod_l s^*_l$. Since restricting to local (product) POVMs can only increase the primal optimum, $s^*_{\mathrm{loc}} \geq s^*_{\mathrm{global}} \geq \prod_l s^*_l$.

\subsection{Per-site analysis for Pauli observables}

\emph{Linear case ($k=1$).} For Pauli observables $\{X,Y,Z\}$ on $\mathbb{C}^2$, the per-site SDP asks: minimize $s$ subject to $\sum_j h_P(j)\Pi_j = P$ for $P \in \{X,Y,Z\}$, $\sum_j \Pi_j = s\,\1$, $\Pi_j \geq 0$, $|h_P(j)| \leq 1$. By Theorem~\ref{thm:linear-equiv}, this equals the joint measurability robustness of unsharp Paulis $\{\frac{\1+\eta P}{2}\}$, which gives $\eta_{\max} = 1/\sqrt{3}$ and $s^* = 1/\eta_{\max} = \sqrt{3}$~\cite{Busch1986,heinosaari2015noise}. For a weight-$w$ Pauli, $s^*_{\mathrm{loc}} = (\sqrt{3})^w$, yielding variance $3^w$.

\emph{Quadratic case ($k=2$).} For squared expectations $\{|{\tr(P\rho)}|^2\}$, the linearizations $\{P\ox P\}$ commute (Sec.~A.3). Bell measurement gives $s^*_l = 1$ at every site. Hence $s^*_{\mathrm{loc}} = 1^n = 1$, and the sample complexity is $\mathcal{O}(1)$ independent of $n$ and $w$.

\subsection{Spin-1 structure of nonlinear Pauli operators}

For estimating $\tr(P\rho^2)$, the natural two-copy operator is $\tilde{O}_P = \frac{1}{2}(P\ox\1 + \1\ox P)\cdot\mathrm{SWAP}$. On the symmetric subspace $\mathrm{Sym}^2(\mathbb{C}^2)$, this operator has eigenvalues $\{+1, 0, -1\}$: it is a spin-1 observable. Since it is not binary (eigenvalues $\pm 1$ only), Theorem~\ref{thm:linear-equiv} does not apply, and f-compatibility may in principle be strictly easier than standard 2-copy compatibility.

The four-copy operators $A_P$ from Theorem~\ref{thm:nonlocal} with $k=2$ take the form $A_P = \frac{1}{2}(P^{\ox 2}\ox\1^{\ox 2} + \1^{\ox 2}\ox P^{\ox 2})\cdot\bP_\pi^d$. Critically, $A_P$ is a \emph{sum} of two product operators, not itself a product. Therefore the tensor product factorization (Theorem~\ref{thm:tensor}) does not apply to $\{A_P\}$, and one cannot decompose the global estimation problem into independent per-site SDPs. This non-product structure is the fundamental obstruction preventing local strategies from achieving the efficiency of the nonlocal four-copy measurement.


%%%%%%%%%%%%%%%%%%%%%%%%%%%%%%%%%%%%%%%%%%%%%%%%%%%%%%%%%%%%%%%%%%%%%%%%%%%%%%%%%%%%%%%%%%%%%%%%%%%
%%%%%%%%%%%%%%%%%%%%%%%%%%%%%%%%%%%%%%%%%%%%%%%%%%%%%%%%%%%%%%%%%%%%%%%%%%%%%%%%%%%%%%%%%%%%%%%%%%%

\section{Algebraic framework for multi-copy estimation}\label{app:alg_framework}

This section develops the algebraic and information-theoretic foundations underlying the multi-copy estimation framework of the main text. We introduce the $k$-copy tensor algebra, characterize its commutative subalgebras, and connect the optimal estimation problem to approximate joint diagonalization (AJD) and SDP relaxations.

%%%%%%%%%%%%%%%%%%%%%%%%%%%%%%%%%%%%%%%%%%%%%%%%%%%%%%%%%%%%%%%%%%%%%%%%%%%%%%%%%%%%%%%%%%%%%%%%%%%
\subsection{$k$-copy tensor algebra and estimable families}\label{app:tensor_algebra}

Let $\{O_1, \ldots, O_n\}$ be Hermitian operators on $\cH$ ($\dim\cH = d$), and let $G = \gen{O_1, \ldots, O_n}$ denote the group (or $*$-algebra) they generate. On $k$ copies, define the \emph{$k$-copy tensor algebra}:
\begin{align}
\Alg_k(G) := \Sp\{O_{i_1}\ox\cdots\ox O_{i_k} : i_j \in \{0,1,\ldots,n\}\} \subset \cL(\cH^{\ox k}),
\end{align}
where $O_0 := \1$.

\begin{definition}[$k$-Estimable family]\label{def:k-estimable}
A set of Hermitian operators $\{R_\alpha\}_{\alpha \in \cS} \subset \Alg_k(G)$ is a \emph{$k$-estimable family} if:
\begin{enumerate}[label=(\roman*)]
\item $[R_\alpha, R_\beta] = 0$ for all $\alpha,\beta \in \cS$ (mutual commutativity),
\item $R_\alpha^\dagger = R_\alpha$ for all $\alpha$ (Hermiticity).
\end{enumerate}
The expectation values $\{\tr[R_\alpha\,\rho^{\ox k}]\}_\alpha$ can be simultaneously estimated from a single projective measurement in the common eigenbasis of $\{R_\alpha\}$.
\end{definition}

\begin{definition}[Polynomial map]\label{def:poly-map}
Given a $k$-estimable family $\{R_\alpha\}_\alpha$, the \emph{polynomial map} $\Phi: \mathbb{R}^n \to \mathbb{R}^{|\cS|}$ is defined by
\begin{align}
\Phi_\alpha(\vec{t}) := \tr[R_\alpha\,\rho^{\ox k}] = g_\alpha(t_1,\ldots,t_n,\; \{t_{ij}\}_{i<j},\; \ldots),
\end{align}
where $t_i = \tr(O_i\rho)$ are the target parameters and $t_{ij} = \tr(O_iO_j\rho)$ etc.\ are auxiliary parameters arising from non-closure of $\{O_i\}$.
\end{definition}

\emph{Remark (non-closure and auxiliary parameters).}
When $\{O_i\}$ does not form a closed group, $g_\alpha$ depends on moments $\tr(O_{i_1}\cdots O_{i_m}\rho)$ that are not simple functions of $\{t_i\}$. The NPA moment matrix $M_l(y) \succeq 0$~\cite{navascues2008convergent,pironio2010convergent} constrains these auxiliary parameters: given observed values $\{g_\alpha\}_\alpha$ and the PSD constraint, the feasible set for $\{t_i\}$ is a projection of the spectrahedron $\{y : M_l(y) \succeq 0\}$. When $G$ is a group, all higher moments reduce to linear moments via structure constants: $\tr(O_iO_j\rho) = \sum_k c_{ij}^k t_k$, and $g_\alpha$ becomes a polynomial in $\{t_i\}$ alone.


%%%%%%%%%%%%%%%%%%%%%%%%%%%%%%%%%%%%%%%%%%%%%%%%%%%%%%%%%%%%%%%%%%%%%%%%%%%%%%%%%%%%%%%%%%%%%%%%%%%
\subsection{Three levels of commutative structure}\label{app:three_levels}

For a group $G = \{O_1, \ldots, O_N\}$, the commutative subalgebras of the $k$-copy algebra organize into a hierarchy.

\begin{proposition}[Hierarchy of commutative subalgebras]\label{prop:three-levels}
The following three levels are ordered by inclusion:

\emph{Level~1 (Group center, $k=1$, degree $0$).}
The center $Z(G) = \{g \in G : [g,h]=0\;\forall h\in G\}$ gives observables measurable alongside all others. Map: $g_\alpha(\vec{t}) = c$ (constants only). Size: $|\cS| = |Z(G)|$.

\emph{Level~2 (MASA of $\mathbb{C}[G]$, $k=1$, degree $1$).}
A maximal abelian subalgebra (MASA) of $\mathbb{C}[G]$ gives the largest set of simultaneously measurable linear combinations. Map: $g_\alpha(\vec{t}) = \vec{c}_\alpha \cdot \vec{t}$ (linear, no information loss). Size: $|\cS| \leq d$.

\emph{Level~3 ($k$-copy tensor MASA, degree $\leq k$).}
A MASA of $\Alg_k(G)$ gives polynomial functions of degree $\leq k$. Map: $g_\alpha(\vec{t}) = \sum c_{i_1\cdots i_k} t_{i_1}\cdots t_{i_k}$, potentially covering all observables but with information loss from the fiber structure ($\S$\ref{app:fiber}). Size: $|\cS| \leq d^k$.
\end{proposition}

\begin{table}[h]
\centering
\begin{tabular}{c|c|c|c|c}
\textbf{Level} & \textbf{Commutative in} & $|\cS|_{\max}$ & $\deg(g)$ & \textbf{Info loss} \\ \hline
1 & $G$ & $|Z(G)|$ & 0 & N/A \\
2 & $\mathbb{C}[G]$ & $d$ & 1 & None \\
3 & $\Alg_k(G)$ & $d^k$ & $\leq k$ & Fiber size
\end{tabular}
\caption{Three levels of commutative structure for a group $G$ of operators.}
\label{tab:three_levels}
\end{table}

\emph{Key example ($\cP_n$ Pauli group).}
Level~2 gives the stabilizer group: $|\cS| = 2^n$ out of $4^n$ Paulis, with a linear map.
Level~3 at $k=2$ gives $\{P\ox P\}_{P\in\cP_n}$, covering all $|\cS| = 4^n$ Paulis via the quadratic map $g_P = t_P^2$, exploiting the squaring trick~\eqref{eq:squaring}.


%%%%%%%%%%%%%%%%%%%%%%%%%%%%%%%%%%%%%%%%%%%%%%%%%%%%%%%%%%%%%%%%%%%%%%%%%%%%%%%%%%%%%%%%%%%%%%%%%%%
\subsection{Fiber structure and information loss}\label{app:fiber}

\begin{definition}[Fiber and information loss]\label{def:fiber}
Given a polynomial map $\Phi: \mathbb{R}^n \to \mathbb{R}^{|\cS|}$ induced by a $k$-estimable family, the \emph{fiber} over an observed value $\vec{v}$ is
\begin{align}
\Phi^{-1}(\vec{v}) = \{\vec{t} \in \cB : \Phi(\vec{t}) = \vec{v}\},
\end{align}
where $\cB = \{\vec{t} : \rho(\vec{t}) \geq 0\}$ is the Bloch body. The \emph{fiber cardinality} $D(\vec{v}) = |\Phi^{-1}(\vec{v})|$ measures information loss: $D = 1$ means perfect recovery, $D > 1$ means ambiguity.
\end{definition}

\begin{proposition}[B\'ezout bound]\label{prop:bezout}
If $g_1,\ldots,g_n$ have degrees $d_1,\ldots,d_n$ and the system $g_\alpha(\vec{t}) = v_\alpha$ is zero-dimensional, then the generic fiber satisfies
\begin{align}\label{eq:bezout}
|\Phi^{-1}(\vec{v})| \leq \prod_{\alpha=1}^n d_\alpha,
\end{align}
with equality generically over $\mathbb{C}$. The real fiber over $\cB$ may be smaller.
\end{proposition}

\emph{Example.} For $\{X\ox X, Y\ox Y, Z\ox Z\}$ on a single qubit, $g_P(\vec{t}) = t_P^2$, giving $\Phi^{-1}(\vec{v}) = \{(\pm\sqrt{v_1}, \pm\sqrt{v_2}, \pm\sqrt{v_3})\} \cap \cB$, with up to $2^3 = 8$ points reduced by the Bloch ball constraint $t_X^2+t_Y^2+t_Z^2 \leq 1$.


%%%%%%%%%%%%%%%%%%%%%%%%%%%%%%%%%%%%%%%%%%%%%%%%%%%%%%%%%%%%%%%%%%%%%%%%%%%%%%%%%%%%%%%%%%%%%%%%%%%
\subsection{Fisher information as pinching overlap}\label{app:fisher_pinching}

\begin{theorem}[Fisher information as pinching overlap]\label{thm:fisher-pinching}
Let $\{R_\alpha\}$ be a $k$-estimable family with common eigenbasis $\{|\lambda\rangle\}$, and let $\cP(\cdot) = \sum_\lambda \proj{\lambda}(\cdot)\proj{\lambda}$ be the associated pinching map. For the parametrization $\rho(\vec{t}) = \frac{1}{d}(\1 + \sum_s t_s O_s)$, the classical Fisher information matrix from measuring $\{R_\alpha\}$ on $\rho_0^{\ox k} = (\1/d)^{\ox k}$ is
\begin{align}\label{eq:fisher-pinching}
\cI_{ss'} = \frac{k^2}{d^k}\,\tr\!\big[\cP(O_s^{(k)})\cdot\cP(O_{s'}^{(k)})\big] = \frac{k^2}{d^k}\,\langle \cP(O_s^{(k)}),\, \cP(O_{s'}^{(k)})\rangle_{\mathrm{HS}},
\end{align}
where $O_s^{(k)} = \sym[O_s \ox \1^{\ox(k-1)}]$ is the $k$-replica symmetric extension.
\end{theorem}

\begin{proof}
Parametrize $\rho = \frac{1}{d}(\1 + \sum_s t_s O_s)$ so that $\partial\rho/\partial t_s = O_s/d$. The derivative of the outcome probability $p_\lambda = \tr[\proj{\lambda}\,\rho^{\ox k}]$ at $\rho_0 = \1/d$ is
\begin{align}
\frac{\partial p_\lambda}{\partial t_s}\bigg|_{\rho_0} = \frac{k}{d^k}\,\bra{\lambda}O_s^{(k)}\ket{\lambda},
\end{align}
where the factor $k$ arises from differentiating the product $\rho^{\ox k}$ and the symmetrization of $O_s$ across copies. At $\rho_0$, all outcomes are equiprobable: $p_\lambda^{(0)} = d^{-k}$. Thus:
\begin{align}
\cI_{ss'} &= \sum_\lambda \frac{1}{p_\lambda^{(0)}}\frac{\partial p_\lambda}{\partial t_s}\frac{\partial p_\lambda}{\partial t_{s'}} = \frac{k^2}{d^k}\sum_\lambda \bra{\lambda}O_s^{(k)}\ket{\lambda}\bra{\lambda}O_{s'}^{(k)}\ket{\lambda} = \frac{k^2}{d^k}\,\tr[\cP(O_s^{(k)})\cdot O_{s'}^{(k)}],
\end{align}
where the last step uses $\cP^2 = \cP = \cP^\dagger$ (orthogonal projection on Hilbert-Schmidt space).
\end{proof}

\emph{Corollary.} The total Fisher information is
\begin{align}\label{eq:total-fisher}
\tr[\cI] = \frac{k^2}{d^k}\sum_s \|\cP(O_s^{(k)})\|_{\mathrm{HS}}^2.
\end{align}
Since $\cP$ is an orthogonal projection, $\|\cP(A)\|^2 \leq \|A\|^2$ with equality iff $A$ lies in the commutative algebra. The information loss per observable is
\begin{align}\label{eq:info-loss}
\Delta_s = \|O_s^{(k)}\|_{\mathrm{HS}}^2 - \|\cP(O_s^{(k)})\|_{\mathrm{HS}}^2 = \|O_s^{(k)} - R_s\|_{\mathrm{HS}}^2 \geq 0,
\end{align}
where $R_s = \cP(O_s^{(k)})$ is the commuting projection.


%%%%%%%%%%%%%%%%%%%%%%%%%%%%%%%%%%%%%%%%%%%%%%%%%%%%%%%%%%%%%%%%%%%%%%%%%%%%%%%%%%%%%%%%%%%%%%%%%%%
\subsection{Equivalence to approximate joint diagonalization}\label{app:ajd}

Maximizing $\tr[\cI]$ from Eq.~\eqref{eq:total-fisher} is equivalent to minimizing $\sum_s \Delta_s$, i.e., finding the commuting family closest to $\{O_s^{(k)}\}$ in total Hilbert-Schmidt distance:
\begin{align}\label{eq:min-distance}
\min_{\{R_s\}} \sum_s \|R_s - O_s^{(k)}\|_{\mathrm{HS}}^2 \quad\text{s.t.}\quad [R_s, R_{s'}] = 0\;\;\forall\,s,s'.
\end{align}
Since commuting Hermitian operators share an eigenbasis $\{|\lambda\rangle\}$, substituting $R_s = \sum_\lambda r_{s,\lambda}\proj{\lambda}$ and optimizing over $r_{s,\lambda}$ at fixed basis gives $r_{s,\lambda}^* = \bra{\lambda}O_s^{(k)}\ket{\lambda}$ (the pinching). The problem reduces to optimizing over the eigenbasis.

\begin{theorem}[Optimal estimation $\Leftrightarrow$ AJD]\label{thm:ajd}
The problem of finding the $k$-estimable family maximizing total Fisher information is equivalent to
\begin{align}\label{eq:ajd}
\max_{U \in \mathrm{U}(d^k)} \;\sum_{s=1}^n \|\diag(U^\dagger O_s^{(k)} U)\|^2,
\end{align}
i.e., finding the unitary $U$ that \emph{approximately jointly diagonalizes} the symmetric extensions $\{O_s^{(k)}\}$.
\end{theorem}

This is the well-studied \emph{Approximate Joint Diagonalization} (AJD) problem in signal processing~\cite{cardoso1996joint}, solvable by Jacobi-type rotation algorithms. The Jacobi AJD algorithm iterates over pairs $(i,j)$ with $i < j$, finding the optimal $2\times 2$ Givens rotation angle $\theta^*$ that maximizes the diagonal energy $\sum_s [|\cos\theta\,(G_s)_{ii} + \sin\theta\,(G_s)_{ij}|^2 + |\cos\theta\,(G_s)_{jj} - \sin\theta\,(G_s)_{ji}|^2]$, where $G_s$ is the $2\times 2$ submatrix of $U^\dagger O_s^{(k)} U$. Each sub-problem has a closed-form solution. The complexity is $\mathcal{O}(D^2 n)$ per sweep with $D = d^k$, and typically $\mathcal{O}(10)$ sweeps suffice for convergence.


%%%%%%%%%%%%%%%%%%%%%%%%%%%%%%%%%%%%%%%%%%%%%%%%%%%%%%%%%%%%%%%%%%%%%%%%%%%%%%%%%%%%%%%%%%%%%%%%%%%
\subsection{SDP relaxation via NPA moment matrix}\label{app:sdp_npa}

The optimization~\eqref{eq:min-distance} is a QCQP (quadratic objective with quadratic commutativity constraints). We relax it to an SDP via the NPA moment matrix approach~\cite{navascues2008convergent,pironio2010convergent}.

Let $\{B_a\}_{a=1}^{D^2}$ be a Hermitian orthonormal basis of $\cL(\cH^{\ox k})$ with $\tr[B_a B_b] = \delta_{ab}$ and $D = d^k$. Write $R_s = \sum_a r_{s,a}\,B_a$ and define:
\begin{itemize}
\item Overlap coefficients: $\omega_{s,a} = \tr[O_s^{(k)} B_a]$,
\item Commutator structure constants: $C_{abc} = \tr[B_c(B_aB_b - B_bB_a)]$.
\end{itemize}

\begin{theorem}[SDP relaxation]\label{thm:sdp-relaxation}
The following SDP provides an upper bound on the optimal Fisher information:
\begin{align}\label{eq:sdp-relax}
\max\; & \sum_s \sum_a \omega_{s,a}\, r_{s,a} \notag\\
\text{s.t.}\; & \begin{pmatrix} 1 & \vec{r}^{\,T} \\ \vec{r} & Y \end{pmatrix} \succeq 0, \notag\\
& \sum_{a,b} C_{abc}\, Y_{(s,a),(s',b)} = 0 \quad \forall\,c,\;\forall\,s < s', \notag\\
& \sum_a Y_{(s,a),(s,a)} \leq \tr[(O_s^{(k)})^2] \quad \forall\,s,
\end{align}
where $\vec{r} = (r_{s,a}) \in \mathbb{R}^{nD^2}$ and $Y \in \mathbb{S}_+^{nD^2}$.
\end{theorem}

The constraints encode: (i)~\emph{moment matrix PSD} --- enforces $Y_{(s,a),(s',b)} \approx r_{s,a}\,r_{s',b}$, exact when $Y = \vec{r}\,\vec{r}^T$ (rank-1); (ii)~\emph{linearized commutativity} --- the operator identity $[R_s, R_{s'}] = 0$ means $\sum_{a,b} C_{abc}\,r_{s,a}\,r_{s',b} = 0$ for all $c$, relaxed to linear constraints on $Y$; (iii)~\emph{Archimedean bound} --- ensures $\tr[R_s^2] \leq \tr[(O_s^{(k)})^2]$, guaranteeing bounded operators.

\emph{Tightness.} The SDP is tight when the moment matrix admits a rank-1 completion $Y = \vec{r}\,\vec{r}^T$. By the flatness criterion of the NPA hierarchy~\cite{pironio2010convergent}, this is guaranteed in the limit of increasing relaxation order.


%%%%%%%%%%%%%%%%%%%%%%%%%%%%%%%%%%%%%%%%%%%%%%%%%%%%%%%%%%%%%%%%%%%%%%%%%%%%%%%%%%%%%%%%%%%%%%%%%%%
\subsection{Recovery of known results}\label{app:recovery}

\begin{proposition}[Recovery of known results]\label{prop:recovery}
The algebraic framework recovers the following as special cases:

\emph{(a) Standard compatibility ($k=1$, linear).}
With $k=1$, the $k$-estimable families are exactly the jointly measurable sets. The Fisher information reduces to $\cI_{ss'} = d^{-1}\tr[\cP(O_s)O_{s'}]$. The SDP~\eqref{eq:sdp-relax} reduces to the standard compatibility SDP (Sec.~\ref{app:f-compat}).

\emph{(b) Bell measurement ($k=2$, Pauli $\{X,Y,Z\}$).}
The AJD problem~\eqref{eq:ajd} with $n=3$, $d=2$, $k=2$ has optimal solution $U = U_{\mathrm{Bell}}$, giving $R_s = \sigma_s \ox \sigma_s$ with Fisher information $\cI_{ss'} = \delta_{ss'}\cdot t_s^2$ (vanishes at $\rho_0$, reflecting degree-2 information loss via the fiber $\S$\ref{app:fiber}).

\emph{(c) Nonlocal Pauli shadow ($k=2$, all $\cP_n$).}
The commuting family $\{A_P\}$ from Theorem~\ref{thm:nonlocal} is the AJD solution for the full Pauli group, with polynomial map $g_P(\vec{t}) = [\tr(P\rho^k)]^2$.

\emph{(d) f-Compatibility SDP as dual.}
The f-compatibility SDP~\eqref{eq:primal-SDP} is the dual perspective: POVM elements $\{\Pi_j\}$ correspond to common eigenprojectors of the optimal $\{R_\alpha\}$, and estimators $h_\alpha(j)$ correspond to eigenvalues $r_{\alpha,\lambda}$.
\end{proposition}

\begin{theorem}[Tensor commutativity uniqueness, informal~\cite{king2024exponential}]\label{thm:king-nogo}
The technique of achieving commutativity via simple tensor products ($O \ox O$ commuting when the originals anticommute) is essentially unique to displacement operators arising from the Heisenberg-Weyl group. Pauli operators are a special case.
\end{theorem}

This means for general $\{O_i\}$ beyond the Heisenberg-Weyl group, achieving $|\cS| = n$ with degree-2 polynomials may not be possible. The SDP framework~\eqref{eq:sdp-relax} then quantifies how much information can be extracted as a function of $k$ and the algebraic structure of $\{O_i\}$.


%%%%%%%%%%%%%%%%%%%%%%%%%%%%%%%%%%%%%%%%%%%%%%%%%%%%%%%%%%%%%%%%%%%%%%%%%%%%%%%%%%%%%%%%%%%%%%%%%%%
\subsection{Three-level information recovery for the nonlocal strategy}\label{app:three_level_recovery}

For the nonlocal strategy of Theorem~\ref{thm:nonlocal}, the Fisher information framework reveals a three-level structure of information recovery.

\emph{Layer~1: Direct symmetric extension ($\cI = 0$).}
The $2k$-replica symmetric extension $O_P^{(2k)} = \sym[P \ox \1^{\ox(2k-1)}]$ contains the linear signal $\tr(P\rho)$, but it does not lie in any large commutative subalgebra. The polynomial map $g_P(\vec{t}) = [\tr(P\rho^k)]^2$ has vanishing derivative at $\rho_0 = \1/d$, since $\tr(P\rho_0^k) = d^{-k}\tr(P) = 0$ for traceless $P$. Hence $\cI_{PP}|_{\rho_0} = 0$.

\emph{Layer~2: Commuting projection ($\cI \propto t^2$).}
The optimal commuting projection $R_P = A_P = \sym_\pi[P^{\ox 2} \ox \1^{\ox(2k-2)}]\cdot\bP^d_\pi$ gives $\tr(R_P\,\rho^{\ox 2k}) = [\tr(P\rho^k)]^2$. The Fisher information becomes
\begin{align}
\cI_{PP}(\vec{t}) = 4\,[\tr(P\rho^k)]^2 \cdot \left(\frac{\partial\,\tr(P\rho^k)}{\partial t}\right)^2,
\end{align}
which is nonzero when $\tr(P\rho^k) \neq 0$ but vanishes as $t^2$ near $\rho_0$. This is the root cause of the $\varepsilon^{-4}$ sample complexity (rather than $\varepsilon^{-2}$).

\emph{Layer~3: Signal estimation correction ($\cI \propto \mathrm{const}$).}
The Signal Estimation lemma (Lemma~\ref{lemma:signal_estim}) replaces $\rho^{\ox 2k}$ with $\rho^{\ox k}\ox\sigma^{\ox k}$, giving $\tr(A_P\,\rho^{\ox k}\ox\sigma^{\ox k}) = \tr(P\rho^k)\cdot\tr(P\sigma^k)$. When $\sigma$ is chosen such that $|\tr(P\sigma^k)| \geq \varepsilon$, this becomes a \emph{linear} function of $\tr(P\rho^k)$ with nonvanishing Fisher information. The cost is needing a preliminary estimate of $|\tr(P\rho^k)|$ to construct $\sigma$, resulting in $\varepsilon^{-4}$ total sample complexity.

\begin{table}[h]
\centering
\begin{tabular}{c|c|c|c}
\textbf{Layer} & \textbf{Observable} & $\cI$ \textbf{at} $\rho_0$ & \textbf{Samples} \\ \hline
Ideal (incompatible) & $\tr(P\rho)$ & $\mathcal{O}(1)$ & $\varepsilon^{-2}$ \\
2 (commuting) & $[\tr(P\rho^k)]^2$ & $0$ & $\varepsilon^{-4}$ \\
3 (+ ancilla) & $\tr(P\rho^k)\cdot\tr(P\sigma^k)$ & $\mathcal{O}(1)$ & $\varepsilon^{-4}$
\end{tabular}
\caption{Three layers of information recovery. The Fisher information at $\rho_0$ determines the achievable scaling.}
\label{tab:three_levels_recovery}
\end{table}


%%%%%%%%%%%%%%%%%%%%%%%%%%%%%%%%%%%%%%%%%%%%%%%%%%%%%%%%%%%%%%%%%%%%%%%%%%%%%%%%%%%%%%%%%%%%%%%%%%%
\subsection{Minimum copies for $\tr(P\rho^2)$: the 4-copy barrier}\label{app:min_copies}

We analyze the minimum number of copies $N$ needed to simultaneously estimate $\{\tr(P\rho^2)\}_{P\in\cP_n}$ for all $4^n$ Paulis.

\emph{Operator structure.}
The natural two-copy operator is $O_P = (P \ox \1) \cdot \mathrm{SWAP}$ with $\tr[O_P\,\rho^{\ox 2}] = \tr(P\rho^2)$. Its Hermitian symmetrization $\tilde{O}_P = \frac{1}{2}(P\ox\1 + \1\ox P)\cdot\mathrm{SWAP}$ satisfies the commutation relation
\begin{align}\label{eq:tilde-commutator}
[\tilde{O}_P,\, \tilde{O}_Q] = \tfrac{1}{2}[P,Q]^{(2)},
\end{align}
\color{red}where $[P,Q]^{(2)} := \tfrac{1}{2}([P,Q]\ox\1 + \1\ox [P,Q])$ is the symmetrized lift of the commutator. Since $[\tilde{O}_P,\tilde{O}_Q] = \tfrac{1}{2}[P,Q]^{(2)} \neq 0$ whenever $[P,Q]\neq 0$, the operators $\{\tilde{O}_P\}$ inherit the non-commutativity structure of the original Paulis.\color{black}

\emph{Why tensor squaring of $\tilde{O}_P$ fails.}
The squaring trick~\eqref{eq:squaring} succeeds for Paulis because they satisfy strict anticommutation: $\{P,Q\} = 0 \Rightarrow (PQ)^{\ox 2} = (QP)^{\ox 2}$. However, \color{red}for non-commuting (and hence anticommuting) Paulis $P,Q$ with $\{P,Q\}=0$,\color{black} $\tilde{O}_P$ has a nonzero anticommutator:
\begin{align}
\{\tilde{O}_P, \tilde{O}_Q\} = \tfrac{1}{2}(P\ox Q + Q\ox P) \neq 0.
\end{align}
Writing $\tilde{O}_P\tilde{O}_Q = S + C$ (symmetric + antisymmetric parts), the commutator of the tensor-squared operators is
\begin{align}
[\tilde{O}_P\ox\tilde{O}_P,\; \tilde{O}_Q\ox\tilde{O}_Q] = 2(S\ox C + C\ox S) \neq 0.
\end{align}
The nonzero anticommutator $S$ is an irreducible obstruction.

\emph{Why 3 copies are insufficient.}
On three copies, any pair of operator pairs $(ij)$ and $(kl)$ chosen from $\{1,2,3\}$ must share at least one index, since there are only $\binom{3}{2} = 3$ distinct pairs and no two are disjoint. The symmetrized operator $B_P = \tilde{O}_P^{[12]} + \tilde{O}_P^{[13]} + \tilde{O}_P^{[23]}$ gives $3\,\tr(P\rho^2)$ but $[B_P, B_Q] \neq 0$. Furthermore, $[\tr(P\rho^2)]^2$ is degree-4 in $\rho$, which cannot be accessed from 3 copies (only degree $\leq 3$ polynomials are available).

\emph{4 copies suffice.}
The construction of Theorem~\ref{thm:nonlocal} with $k=2$ gives a commuting family $\{A_P\}$ on 4 copies with $\tr(A_P\,\rho^{\ox 4}) = [\tr(P\rho^2)]^2$. This construction bypasses the anticommutation obstruction by using $P^{\ox 2}$ directly (which preserves the original Pauli anticommutation) combined with the permutation $\pi = (13)(24)$.

\begin{table}[h]
\centering
\begin{tabular}{c|c|c|c}
$N$ \textbf{copies} & \textbf{Achievable} & \textbf{Coverage} & \textbf{Info loss} \\ \hline
2 & $\tr(P\rho^2)$ (linear) & $2^n/4^n$ & None \\
3 & $\tr(P\rho^2)$ (linear) & $2^n/4^n$ & None \\
4 & $[\tr(P\rho^2)]^2$ & all $4^n$ & Sign lost \\
4 + ancilla & $\tr(P\rho^2)$ (via Lem.~\ref{lemma:signal_estim}) & all $4^n$ & $\varepsilon^{-4}$ cost
\end{tabular}
\caption{Minimum copies needed to estimate $\{\tr(P\rho^2)\}_{P\in\cP_n}$. Linear estimation covers only stabilizer-size subsets; 4 copies with the squaring trick cover all Paulis.}
\label{tab:min_copies}
\end{table}

This analysis demonstrates the power of the algebraic framework: by identifying the Lie algebra structure of $\{\tilde{O}_P\}$ (isomorphic to the original Paulis), one can invoke the Level~3 degree counting from Proposition~\ref{prop:three-levels} to determine the minimum copy number without exhaustive search over all possible commuting polynomials.

\section{Multi-copy local estimation of fixed $\mathrm{tr}(P \rho^2)$}\label{app:local_est}
\subsection{Bell measurement and its restriction}
Contrast to conclusion in the last section, we noted that the purity estimation strategy using Bell measurements already reach the $\mathcal{O}(1)$ sampling complexity.
Bell measurement is a two-copy local measurement strategy.
 The Bell measurement projector satisfied
\begin{align}
    \ket{\Phi_{00}}\!\bra{\Phi_{00}}&:= \frac{(II + XX - YY + ZZ)}{4},\ket{\Phi_{10}}\!\bra{\Phi_{10}}:= \frac{(II + XX + YY - ZZ)}{4},\\
    \ket{\Phi_{01}}\!\bra{\Phi_{01}}&:= \frac{(II - XX + YY + ZZ)}{4},\ket{\Phi_{11}}\!\bra{\Phi_{11}}:= \frac{(II - XX - YY - ZZ)}{4}.
\end{align}
Thus, in terms of sympletic representation, we have $\ket{\Phi_{\mathbf{r}}}\!\bra{\Phi_{\mathbf{r}}} = 4^{-n}\sum_{\mathbf{p}} (-1)^{[\mathbf{r}, \mathbf{p}] + \mathbf{p}_x\cdot \mathbf{p}_z}\sigma_{\mathbf{p}} \otimes \sigma_{\mathbf{p}} $ and
\begin{align}
\label{eq:Bell_deco}
    \mathcal{M}_{I}(\rho^{\otimes 2}) &= 16^{-n}\sum_{\mathbf{r}}  \sum_{\mathbf{p},\mathbf{q}} (-1)^{[\mathbf{r}, \mathbf{p} + \mathbf{q}] + \mathbf{p}_x\cdot \mathbf{p}_z +  \mathbf{q}_x\cdot \mathbf{q}_z}\mathrm{tr}\left[(\sigma_{\mathbf{p}} \otimes \sigma_{\mathbf{p}})\cdot\rho^{\otimes 2}\right] \sigma_{\mathbf{q}} \otimes \sigma_{\mathbf{q}}\\
    &= 4^{-n}\sum_{\mathbf{p}}\mathrm{tr}\left[(\sigma_{\mathbf{p}} \otimes \sigma_{\mathbf{p}})\cdot\rho^{\otimes 2}\right]\sigma_{\mathbf{p}} \otimes \sigma_{\mathbf{p}}.
\end{align}
In the langurage of shadow tomography, this channel on the two-replica space $\mathcal{H}(d^2)$ is not invertible since $\mathcal{M}_{I}(\sigma_{\mathbf{p}} \otimes \sigma_{\mathbf{q}}) = 0$ if $\mathbf{p} \neq \mathbf{q}$. 
However, if we restrained on the operator subspace spanned by $\mathbf{S}:=\mathrm{span}\{\sigma_{\mathbf{p} }\otimes \sigma_{\mathbf{p}}\}_{\mathbf{p} \in \mathbb{F}_2^{2n}}$, this channel takes itself as the inverse.
Then given an operator $O = \sum_{\mathbf{p}} \alpha_{\mathbf{p}} \sigma_{\mathbf{p}} \otimes \sigma_{\mathbf{p}}$, we use Bell measurment result $\{\mathbf{r}\}$ to construct the shadow estimator $\hat{o}(\mathbf{r}) = \mathrm{tr}\left[\mathcal{M}_{I}^{-1,\dagger}[O] \ket{\Phi_{\mathbf{r}}}\!\bra{{\Phi_{\mathbf{r}}}}\right]= \mathrm{tr}\left[O \ket{\Phi_{\mathbf{r}}}\!\bra{{\Phi_{\mathbf{r}}}}\right]$. This estimator has variance $\mathrm{tr}(\rho O^2) - \mathrm{tr}^2(\rho O) \leq \left[(\lambda_{\max}(O) - \lambda_{\min}(O))/2\right]^2$, as shown by the lemma~\ref{lemma:max_variance}:
\begin{lemma}
    \label{lemma:max_variance}
    Let $O$ be a Hermitian operator with the maximum and minimum eigenvalues $\lambda_{\text{max}}$ and $\lambda_{\text{min}}$. Then its variance $\text{Var}_{\rho}(O) := \langle O^2 \rangle_{\rho} - \langle O \rangle_{\rho}^2$ satisfied $ \sup_{\rho \in \mathcal{D}(\mathcal{H}_n)}\text{Var}_{\rho}(O) = \left( {(\lambda_{\text{max}} - \lambda_{\text{min}})}/{2} \right)^2$.
     This maximum is achieved at pure state such as $|\psi\rangle = \frac{1}{\sqrt{2}}(|\psi_{\text{max}}\rangle + |\psi_{\text{min}}\rangle)$.
\end{lemma}
    \begin{proof}
    The variance of an observable $O$ with respect to a quantum state $\rho$ is defined as
    We first show that the function $f(\rho):= \text{Var}_{\rho}(O)$ is a concave function of $\rho$. Use the fact that $[(1 - \theta)x_1 + \theta x_2]^2 \geq (1- \theta)x_1^2 + \theta x_2^2$ due to the convexity of function $g(x) = x^2$, we have
    \begin{align}
        f((1- \theta)\rho + \theta \sigma) - (1- \theta)f(\rho) - \theta f(\sigma) = -\left[(1- \theta)\text{tr}(O\rho) + \theta \text{tr}(O\sigma)\right]^2 + (1- \theta)\left[\text{tr}(O\rho)\right]^2 + \theta \left[\text{tr}(O\sigma)\right]^2 \leq 0, \notag
    \end{align}
    which showed the concavity of $f(\rho)$.
    The maximum value of a concave function over a convex set is attained at an extreme point. Therefore, the problem reduces to maximizing the variance over all pure states $|\psi\rangle,  \text{Var}_{\psi}(O) = \langle\psi|O^2|\psi\rangle - (\langle\psi|O|\psi\rangle)^2 $.
    Let $\{|\lambda_i\rangle\}$ be the orthonormal eigenbasis of $O$ with corresponding eigenvalues $\{\lambda_i\}$, and express pure state $|\psi\rangle$ in this basis as $|\psi\rangle = \sum_i c_i |\lambda_i\rangle$, where $\{p_i:=|c_i|^2\}_i$ forms a valid probability. Then
    $ \text{Var}(O) = \sum_i p_i \lambda_i^2 - \left(\sum_i p_i \lambda_i\right)^2 = \sum_{i} p_i(\lambda_i - \bar{\lambda})^2 $.
    Keeping $\bar \lambda$ unchanged, we should concentrate $p_i$ on the most extreme values available $\lambda_{\text{max}}$ and $\lambda_{\text{min}}$. Let us consider a distribution where probability $p$ is assigned to $\lambda_{\text{max}}$ and $1-p$ to $\lambda_{\text{min}}$. The variance for this distribution is $ \mathbf{V}(p) = [p\lambda_{\text{max}}^2 + (1-p)\lambda_{\text{min}}^2] - [p\lambda_{\text{max}} + (1-p)\lambda_{\text{min}}]^2 $. 
    Assuming $\lambda_{\text{max}} \neq \lambda_{\text{min}}$, we can find the maximum by differentiate with respect to $p$ and find $p=1/2$. This corresponds to an equal mixture of the minimum and maximum eigenvalues.
    Substituting $p=1/2$ back into the variance expression gives 
    \begin{align}
    \text{Var}_{\text{max}}(O) &= \left[\frac{1}{2}\lambda_{\text{max}}^2 + \frac{1}{2}\lambda_{\text{min}}^2\right] - \left[\frac{1}{2}\lambda_{\text{max}} + \frac{1}{2}\lambda_{\text{min}}\right]^2 = \left( \frac{\lambda_{\text{max}} - \lambda_{\text{min}}}{2} \right)^2.
    \end{align}
    \end{proof}
This means that Bell shadow can effeciently estimates operators $O = \sum_{\mathbf{p}} \alpha_{\mathbf{p}} \sigma_{\mathbf{p}} \otimes \sigma_{\mathbf{p}}$ with bounded $\infty$-norm $\|O\|_{\infty}$.
This include the SWAP operator and single Pauli operator $P \otimes P$, which estimates critical values for $\mathrm{tr}(\rho^2)$ and $|\mathrm{tr}(P \rho)|^2$. Their estimator writes 
\begin{align}
    &\hat{o}_{\mathrm{swap}}(\mathbf{r}) = \mathrm{tr}\left[\mathrm{SWAP}^{\otimes n} \ket{\Phi_{\mathbf{r}}}\!\bra{{\Phi_{\mathbf{r}}}}\right] = 2^{-n}\sum_{\mathbf{p}_x,\mathbf{p}_z} (-1)^{\mathbf{r}_z\mathbf{p}_x + (\mathbf{r}_x + \mathbf{p}_x)\mathbf{p}_z} = (-1)^{\mathbf{r}_z \mathbf{r}_x},\\
&\hat{o}_{\mathbf{q}}(\mathbf{r}) =\mathrm{tr}\left[\sigma_{\mathbf{q}} \otimes \sigma_{\mathbf{q}} \ket{\Phi_{\mathbf{r}}}\!\bra{{\Phi_{\mathbf{r}}}}\right] =(-1)^{[\mathbf{r}, \mathbf{q}] + \mathbf{q}_x\cdot \mathbf{q}_z},
\end{align}
respectively. Both of them are bounded within $[-1,1]$, thus induces two effecient estimation algorithm according to the Hoffding inequality.
\begin{figure}[htbp!]
    \centering
    \hrule height 0.8pt \vspace{-8pt} 
\caption{Estimation of $\text{tr}(\rho^2)$~\cite{hangleiter2024bell}}
    \label{alg:Bell_for_purity}
    \vspace{1pt} \hrule \vspace{2pt}
    % content
    \begin{algorithmic}[1]
        \Require  $n$-qubits density matrix $\rho$
        \Ensure Estimates $\text{tr}(\rho^2)$ within $\varepsilon$ additive error and confidence $1 -\delta$.
        \State Implement $N \sim 4 \ln(2/\delta)/\varepsilon^2$ times Bell measurement and get the result $\mathbf{r}_1, \cdots, \mathbf{r}_N$.
        \State For each result $\mathbf{r}_i$, denotes $s_i =\mathbf{r}_{i,x}\cdot \mathbf{r}_{i,z} $, which indicates the number of Bell measurement result $(11)$.
        \State Return ${M}(\rho) = \sum_{i = 1}^N(-1)^{s_i}/{N}$.
        \end{algorithmic}
    \vspace{5pt} \hrule % 底部的横线
\end{figure}
\begin{figure}[htbp!]    
    \centering
    \hrule height 0.8pt \vspace{-8pt} 
    \caption{Estimation of $|\text{tr}(\rho P)|^2$.}
    \label{alg:est_Pauli_const}
    \vspace{1pt} \hrule \vspace{2pt}
        \begin{algorithmic}[1]
        \Require  $n$-qubit density matrix $\rho$ and a set of target Pauli operator $\{\sigma_{\mathbf{y}^{(1)}},\cdots,\sigma_{\mathbf{y}^{(M)}}\}$.
        \Ensure Collect enough classical information such that  $\forall i, |\text{tr}(\rho \sigma_{\mathbf{y}^{(i)}})|^2$ can be estimated within $\varepsilon$ additive error and $1 - \delta$ success rate.
        \State Implements $N \sim 4 \ln(2M/\delta)/\varepsilon^2$ times Bell measurement and restore the result $\mathbf{r}_1, \cdots, \mathbf{r}_N$ with $\mathbf{r}_i \in \mathbb{F}_2^{2n}$.
        \State To estimate $|\text{tr}(\rho \sigma_{\mathbf{y}})|^2$, take the xz decomposition of $\mathbf y = (\mathbf{y}_x, \mathbf{y}_z)$
        \State Calculate the average and return $R =({1}/{N})\cdot\sum_{i = 1}^N(-1)^{[\mathbf r_i, \mathbf y]+\mathbf y_{x} \cdot \mathbf y_z}$.
        \end{algorithmic}
    \vspace{5pt} \hrule % 底部的横线
    \end{figure}

\subsection{A non-adaptive randomized strategy}
Unfortunately, for operators out of $\mathbf{S}$ such as
\begin{align}
    O_{\mathbf{q}}:=(\sigma_{\mathbf{q}} \otimes \mathbb{I}_n) \cdot \mathrm{SWAP}^{\otimes n}=2^{-n}\sum_{\mathbf{p}}\sigma_{\mathbf{q}}\sigma_{\mathbf{p}} \otimes \sigma_{\mathbf{p}},
\end{align}
which is a critical observable that investigate the value of $\mathrm{tr}(\sigma_{\mathbf{r}}\rho^2)$, Bell sampling lost its capacity.
We consider the case $\mathbf{q}_i \neq 0, \forall i = 1,\cdots,n$, otherwise we reorder the terms in $\sigma_{\mathbf{q}}$ such that $O = (\sigma_{\mathbf{q}'} \otimes \mathbb{I}_{n'}) \cdot \mathrm{SWAP}^{\otimes n'} \otimes \mathrm{SWAP}^{\otimes (n -n')} $, and run algorithms $\mathcal{A}$ on the first $n'$ qubits and Bell measurement strategy on the last $(n-n')$ qubits.
Given $\mathbf{p}$, we defined the index set $ A_{\mathbf{p}}:=\{i\in \{1,\cdots,n\}|\mathbf{p}_i \notin \{\mathbf{q}_i,\mathbf{0}\}\}$. For $i\in A_{\mathbf{p}}$, the Pauli operator $\sigma_{\mathbf{p}_i}$ and $\sigma_{\mathbf{q}_i}$ are anti-commute. 
We can then construct $L_\mathbf{q}:= \bigotimes_{i\in A_{\mathbf{p}}} e^{i\frac{\pi}{4} \sigma_{\mathbf{q}_i}} $ and conclude that $ i\sigma_{\mathbf{q}} \sigma_{\mathbf{p}} =L_\mathbf{q}\sigma_{\mathbf{p}}L_\mathbf{q}^\dagger$ on the index set $A_{\mathbf{p}}$.
Unfortunately, for the rest indexes  $\bar{A}_{\mathbf{p}}:= \{i \in \{1,\cdots,n\}| \mathbf{p}_i=\mathbf{q}_i\text{~or~} \mathbf{0}\}$ ,  there exists no local Clifford that realizes $\sigma_{\mathbf{p}_i} \otimes \sigma_{\mathbf{p}_i} \to \sigma_{\mathbf{q}_i}\sigma_{\mathbf{p}_i} \otimes \sigma_{\mathbf{p}_i}$. 
In this case, we can separate $\bar{A}_{\mathbf{p}}$ into two parts $R_{\mathbf{p}}:=\{i| \mathbf{p}_i = \mathbf{q}_i\}$ and $L_{\mathbf{p}}:=\{i| \mathbf{p}_i = \mathbf{0}\}$ and conclude that
\begin{align}
\label{eq:2rep_decomp_Prho2}
    O = 2^{-n}\sum_{\mathbf{p}}\sigma_{\mathbf{q}}\sigma_{\mathbf{p}} \otimes \sigma_{\mathbf{p}}= 2^{-n}\sum_{\mathbf{p}} \bigotimes_{i \in A_{\mathbf{p}}} (-iL_{\mathbf{q}_i} \sigma_{\mathbf{p}_i}L_{\mathbf{q}_i}^{\dagger} \otimes \sigma_{\mathbf{p}_i}) \bigotimes_{i \in R_{\mathbf{p}}} (\mathbb{I} \otimes \sigma_{\mathbf{q}_i}) \bigotimes_{i \in L_{\mathbf{p}}} (\sigma_{\mathbf{q}_i} \otimes \mathbb{I})
\end{align}
For the first part $A_{\mathbf{p}}$, we conjugate one side of Bell projector by local Clifford operation $L$ and defined the shifted Bell sampling projector
\begin{align}
    (L \otimes \mathbb{I}_n)\ket{\Phi_{\mathbf{r}}}\!\bra{\Phi_{\mathbf{r}}}(L^\dagger \otimes \mathbb{I}_n) = 4^{-n}\sum_{\mathbf{p}} (-1)^{[\mathbf{r}, \mathbf{p}] + \mathbf{p}_x\cdot \mathbf{p}_z}(L\sigma_{\mathbf{p}}L^{\dagger}) \otimes \sigma_{\mathbf{p}},
\end{align}
 This then extends the Bell decoherence channel in Eq.\eqref{eq:Bell_deco} into the shifted one:
\begin{align}
\label{eq:Shif_Bell_deco}
    \mathcal{M}_{L}(\rho^{\otimes 2}) = 4^{-n}\sum_{\mathbf{p}}\mathrm{tr}\left[\left[(L\sigma_{\mathbf{p}}L^{\dagger}) \otimes \sigma_{\mathbf{p}}\right]\cdot\rho^{\otimes 2}\right](L\sigma_{\mathbf{p}}L^{\dagger}) \otimes \sigma_{\mathbf{p}}.
\end{align}
For the rest part $R_{\mathbf{p}},L_{\mathbf{p}}$, we have to extract the information of insymmetric Pauli operator $\mathbb{I} \otimes \sigma_{\mathbf{q}}, \sigma_{\mathbf{q}} \otimes \mathbb{I}$, which is prohibitive for local conjugated Bell measurement.
To overcome this, we implement the $\sigma_{\mathbf{q}_i}$ Pauli measurements on the position $i$.
This measurement constitutes the Pauli decoherence channel
\begin{align}
    \Lambda_{i}(\cdot ) = (1/2)\mathbb{I}\cdot\mathrm{tr}(\cdot) + (1/2)\mathrm{tr}\left[\sigma_{\mathbf{q}_i} (\cdot)\right] \sigma_{\mathbf{q}_i}.
\end{align}
In the observable in Eq.~\eqref{eq:2rep_decomp_Prho2}, we average all possible separation.
Since the $\sigma_{\mathbf{q}_i}$ Pauli measurements is incompatiable to the Bell measurment,  this indicates the naive strategy that uniformly selects a separation $A \cup ([n]/A)$ and then implements single-copy $\sigma_{\mathbf q}$ Pauli measurement on positions $[n]/A$ and $L_{\mathbf q}$ conjugated Bell measurement on positions $A$.

% \begin{figure}[htbp!]
%     \centering
%     \hrule height 0.8pt \vspace{-8pt} 
%     \caption{Two-replica local randomized algorithm for $\mathrm{tr}(P\rho^2)$ } % title
%     \label{alg:my_algo}
%     \vspace{1pt} \hrule \vspace{2pt}
%     % content
%     \begin{algorithmic}[1]
%         \Require Initial state $\rho^{\otimes 2}$, target Pauli operator $\sigma_{\mathbf{q}}$ with $\mathbf{q}_i \neq \mathbf{0}$. 
%         \Ensure An estimator $\hat o$ such that $|\hat{o} - \mathrm{tr}(\sigma_{\mathbf{q}} \rho^2)|<\varepsilon$ with probability $1-\delta$.
%         \State Sample the subsets $A \subset [n]$ with uniform probability.
%             \State Apply local Clifford $L_{\mathbf{q}_i} =  e^{i \frac{\pi}{4}\sigma_{\mathbf{q}_i}}$ on the first replica at position $i \in A$.
%             \State Perform Bell measurement on all qubits pairs at $i \in A$. Get the outcome $\mathbf{r} \in \{0,1\}^{2|A|}$.
%             \State Construct estimator $\hat{o}(\mathbf{r}) \leftarrow 2  (-1)^{\mathbf{r}_x \cdot \mathbf{r}_z} $.
%             \State On subsystem $i\in [n]/A$: Perform single-copy Pauli measurement in basis $\sigma_{\mathbf{q}_i}$ on the first and the second replica. Let outcome be $L_i,R_i$.
%             \State Construct estimator $\hat{o} (\mathbf{r}, L_i,R_i) = 2\cdot (-1)^{\mathbf{r}^{E}_{x} \cdot \mathbf{r}^{E}_{z}} \cdot m_I$.
%         \State \Return $\hat{o}$
%     \end{algorithmic}
%     \vspace{5pt} \hrule % 底部的横线
% \end{figure}
  Notice that the measurement of operator $\sigma_{\mathbf{q}_i}$ with result ${r}_i \in \{0,1\}$ forms the $n$-qubit decoherence channel $\Lambda_{\mathbf{q}}(\rho) = \bigotimes_{i= 1}^{n} \Lambda_{i}(\rho)$, where $\Lambda_{i}(\cdot ) = (1/2)\mathbb{I}\cdot\mathrm{tr}(\cdot) + (1/2)\mathrm{tr}\left[\sigma_{\mathbf{q}_i} (\cdot)\right] \sigma_{\mathbf{q}_i}$. 
  We also denote the channel $\mathcal{B}_i$ as the localized shifted Bell decoherence channel as 
  \begin{align}
    \mathcal{B}_{i}(\cdot) = 4^{-1}\sum_{\mathbf{p}_i \in \{00,01,10,11\}}\mathrm{tr}\left[\left[(L_{\mathbf{q}_i}\sigma_{\mathbf{p}_i}L_{\mathbf{q}_i}^{\dagger}) \otimes \sigma_{\mathbf{p}_i}\right](\cdot) \right](L_{\mathbf{q}_i}\sigma_{\mathbf{p}_i}L_{\mathbf{q}_i}^{\dagger}) \otimes \sigma_{\mathbf{p}_i}.
\end{align}
 Then locallly, the shadow channel $\mathcal{M}$ induced by strategies is the tensor products of $\mathcal{M}_i = ({1}/{2}) \mathcal{B}_i + (1/2)\Lambda^{\otimes 2}_{i}$.
we noted that
  \begin{align}
      \left\{
      \begin{aligned}
      &\mathcal{M}_i(\sigma_{\mathbf{q}_i} \otimes \mathbb{I}) = (1/2)\cdot \sigma_{\mathbf{q}_i} \otimes \mathbb{I},\quad \mathcal{M}_i( \mathbb{I}\otimes\sigma_{\mathbf{q}_i}) =(1/2)\cdot \mathbb{I}\otimes\sigma_{\mathbf{q}_i}, \\
      &\mathcal{M}_i\left[(L_{\mathbf{q}_i}\sigma_{\mathbf{p}_i}L_{\mathbf{q}_i}^{\dagger}) \otimes \sigma_{\mathbf{p}_i}\right]=(1/2)\cdot(L_{\mathbf{q}_i}\sigma_{\mathbf{p}_i}L_{\mathbf{q}_i}^{\dagger}) \otimes \sigma_{\mathbf{p}_i},\forall \mathbf{p}_i \notin \{\mathbf{q}_i,\mathbf{0}\}.\\
      \end{aligned}
      \right.
  \end{align}
    Then we conclude that $\mathcal{M}^{-1}(O) = 2^n O$ .
    Noted that
    \begin{align}
        &\mathrm{tr}\left[(L_{\mathbf{q}_i} \otimes \mathbb{I})\Phi_{\mathbf{r}_i}(L^\dagger_{\mathbf{q}_i} \otimes \mathbb{I}) \cdot (\sigma_{\mathbf{q}_i} \otimes \mathbb{I})\cdot \mathrm{SWAP}\right] \\
        &= 4^{-1}\sum_{\mathbf{u}_i} (-1)^{[\mathbf{r}_i, \mathbf{u}_i] +{u}_x\cdot {u}_z}\mathrm{tr}\left[(L_{\mathbf{q}_i} \otimes \mathbb{I})\sigma_{\mathbf{u}_i} \otimes \sigma_{\mathbf{u}_i} (L^\dagger_{\mathbf{q}_i} \otimes \mathbb{I}) \cdot (\sigma_{\mathbf{q}_i} \otimes \mathbb{I})\cdot \mathrm{SWAP}\right]\\
        &=4^{-1}\sum_{\mathbf{u}_i} (-1)^{[\mathbf{r}_i, \mathbf{u}_i] +{u}_x\cdot {u}_z}\mathrm{tr}\left[L_{\mathbf{q}_i} \sigma_{\mathbf{u}_i} L_{\mathbf{q}_i}^\dagger \sigma_{\mathbf{q}_i}\sigma_{\mathbf{u}_i}\right]=-(i/2)\cdot \sum_{\mathbf{u}_i \notin \{\mathbf{0}, \mathbf{q}_i\}} (-1)^{[\mathbf{r}_i, \mathbf{u}_i] + u_xu_z}.\\
        &\mathrm{tr}\left[\left(\frac{ \mathbb{I} + r_{i,x} \sigma_{\mathbf{q}_i}}{2}\right)\otimes  \left(\frac{ \mathbb{I} + r_{i,z} \sigma_{\mathbf{q}_i}}{2}\right)\cdot (\sigma_{\mathbf{q}_i} \otimes \mathbb{I})\cdot \mathrm{SWAP}\right]=|\mathbf{r}_i|/2.
    \end{align}
    We got the estimator
    \begin{align}
        \hat{o}(\{t_i,\mathbf{r}_i\}) =\prod_{i, t_i= 0}\left[-i\cdot \sum_{\mathbf{u}_i \notin \{\mathbf{0}, \mathbf{q}_i\}} (-1)^{[\mathbf{r}_i, \mathbf{u}_i] + u_xu_z} \right]\cdot \prod_{i, t_i= 1}|\mathbf{r}_i|  .
    \end{align}
The variance of this estimator is exponentially large since both weight $|\mathbf{r}_i|$ and summation $\sum_{\mathbf{u}_i \notin \{\mathbf{0}, \mathbf{q}_i\}} (-1)^{[\mathbf{r}_i, \mathbf{u}_i] + u_xu_z} $ range from $[-2,2]$, which gives exponentially large variance.



\subsection{A adaptive local strategy}
However, in the case that the non-local operations is allowed, 
there already exists efficient $2$-replica estimation strategy of $\mathrm{tr}(P\rho^2)$ for a known fixed Pauli operators $P$. 
The simplest one is the SWAP-test that use single anccilla qubits and non-local operation between ancilla qubits and data qubits to estimate the observable $\mathrm{SWAP}^{\otimes n}\cdot( P\otimes\mathbb{I}_n )$, as illustrated in Figure~\ref{fig:nonlocal_strate} for $n=4$ and $P = XZYX$.
If no-ancilla qubit is allowed,  we can find a Clifford gate $C$ that map $P = XZYX$ to $P' = CPC^\dagger= ZIII$. Then $\mathrm{tr}(P\rho^2) = \mathrm{tr}\left((ZIII) \cdot  \sigma^2\right)$ for $\sigma = C^\dagger\rho C$.
Since $ZIII$ contains non-identity Pauli operator only on the first qubits.
We can thus measure the first qubit on the eigenbasis of $2^{-1}(ZI+IZ)\cdot\mathrm{SWAP} = \ket{00}\!\bra{00} - \ket{11}\!\bra{11} $.
The rest $(n-1)$ qubits are measured on the localized eigenbasis of $\mathrm{SWAP}^{\otimes (n-1)}$, which are the $(n-1)$ tensor product of Bell basis.
\begin{figure}[htbp!]
    \centering
    \includegraphics[width=0.3\linewidth]{Figures/ancci_strat.pdf}
    \includegraphics[width=0.2\linewidth]{Figures/nonlocal_strategy.pdf}
    \includegraphics[width=0.4\linewidth]{Figures/local_fixed.png}
    \caption{Two-copy estimation strategies for $\tr(P\rho^2)$: (left) ancilla-based SWAP test, (center) no-ancilla Clifford rotation, (right) local adaptive measurement.}
    \label{fig:nonlocal_strate}
\end{figure}

In this section, we proposed a non-randomized adaptive strategy that target to estimate $\mathrm{tr}(P \rho^2)$ given a known Pauli operator $P$.
As indicated before, it is sufficient to consider the full weight $n$-qubit Pauli $\sigma_{\mathbf{r}_i}$ such that $\sigma_{\mathbf{r}_i} \neq \mathbb{I},\forall i= 1,\cdots, n$.
Then we noted that full weight Pauli operators can be locally transformed into Pauli $Z$ strings $Z^{\otimes n}$ using local Clifford gate since $HXH = Z$ and $HS^{\dagger}YSH = Z$. 
In this case, this observable $O_{\mathbf{r}}$ is a tensor product of local non-Hermitian operators $O_{\mathbf{r}} = [(Z \otimes \mathbb{I}) \cdot \mathrm{SWAP}]^{\otimes n}:= A^{\otimes n}$. The local operator $A$, although non-Hermitian, is a normal matrix that satisfied $[A,A^\dagger] = 0$, thus can still be diagonalized under unitary transformation. Actually, it satisfied the decomposition
\begin{align}
    A &= \ket{00}\!\bra{00} - \ket{11}\!\bra{11} + i\ket{\Phi^+}\!\bra{\Phi^+} - i \ket{\Phi^-}\!\bra{\Phi^-}, \text{~with~} \ket{\Phi^+} = \frac{\ket{01} + i \ket{10}}{\sqrt{2}}, \ket{\Phi^-} = \frac{\ket{01} - i \ket{10}}{\sqrt{2}}. 
\end{align}
Thus we measure locally on the basis $\{\ket{00}\!\bra{00}, \ket{11}\!\bra{11}, \ket{\Phi^+}\!\bra{\Phi^+},\ket{\Phi^-}\!\bra{\Phi^-}\}$ with result $\hat{o}_i \in \{+1,-1,+i,-i\}$, correspondingly. 
Since both $\mathrm{Im}(\prod_{i} \hat{o}_i)$ and $\mathrm{Re}(\prod_{i} \hat{o}_i)$ are bounded, based on the Hoeffding inequality, the average $\mathrm{tr}(\rho A^{\otimes n})$ can be estimated by $\hat{o}:=\prod_{i}\hat{o}_i$ efficiently.

In general, every two-qubit projective measurements that unchanged under qubits exchange can be written in the form
\begin{align}
    \mathcal{B}&:=\left\{\sin \theta \ket{\psi_0}\ket{\psi_0} + \cos \theta \ket{\psi_1}\ket{\psi_1}, \cos \theta \ket{\psi_0}\ket{\psi_0} - \sin \theta \ket{\psi_1}\ket{\psi_1},\right.\\
    &\left.\frac{1}{\sqrt{2}}\ket{\psi_0}\ket{\psi_1} + \frac{i}{\sqrt{2}} \ket{\psi_1}\ket{\psi_0}, \frac{1}{\sqrt{2}}\ket{\psi_0}\ket{\psi_1} - \frac{i}{\sqrt{2}} \ket{\psi_1}\ket{\psi_0}\right\}.
\end{align}
Thus, in general the two-copy disjoint measurement strategy has the form
\begin{align}
   \mathcal{B}(\theta, U)&:=U \otimes U  \{\ket{\Phi_{00}^{\theta}},\ket{\Phi_{01}^{\theta}},\ket{\Phi^+},\ket{\Phi^-}\}\\
   &:= U \otimes U  \{ \sin \theta \ket{00} + \cos \theta \ket{11},  \cos \theta \ket{00} - \sin \theta \ket{11}, \frac{(\ket{01}+i\ket{10})}{\sqrt{2}},  \frac{(\ket{01}-i\ket{10})}{\sqrt{2}}\}.
\end{align}
The basis with parameter $\theta$ decomposed into Pauli operators as
\begin{align}
    &\ket{\Phi_{00}^{\theta}}\!\bra{\Phi_{00}^{\theta}} = \frac{II + ZZ + \sin(2\theta)(XX-YY) - \cos(2\theta)(IZ+ZI)}{4},\\
    &\ket{\Phi_{01}^{\theta}}\!\bra{\Phi_{01}^{\theta}} = \frac{II + ZZ - \sin(2\theta)(XX-YY) + \cos(2\theta)(IZ+ZI)}{4}.
\end{align}
Then given the localized two-copy measurement basis $\mathcal{B}(\theta_i, U_i)$ at site $i$,
%%%%%%%%%%%%%%%%%%%%%%%%%%%%%%%%%%%%%%%%%%%%%%%%%%%%%%%%%%%%%%%%%%%%%%%%%%%%%%%%%%%%%%%%%%%%%%%%%%%
\section{The  multi-copy local shadows}\label{app:local_shadow}
\begin{figure}[htbp!]
    \centering
    \includegraphics[width=1\linewidth]{Figures/shadows.png}
    \caption{Overview of multi-copy shadow tomography protocols.}
    \label{fig:shadows}
\end{figure}
Multi-copy shadow tomography restores $k$-copies of $\rho$, with each copy composed by $n$ qudits. 
Then it repeats randomized local measurements that supports on $k$-qudits,
and try to use these measurement results to estimate \textbf{MANY} properties of $\rho$,
which includes linear property $\text{tr}(O_i\rho)$, non-linear properties $\text{tr}(O_i \rho^{t_i})$, and even their products $\prod_{i}\text{tr}(O_i \rho^{t_i})$.

The simplest non-adaptive case of multi-copy local shadow tomography is shown below:
%%%%%%%%Algorithm%%%%%%%%Algorithm%%%%%%%%Algorithm%%%%%%%%Algorithm
\begin{itemize}
    \item Set $U = \bigotimes_{i=1}^n U_i$, where $(U_1,\cdots,U_n)$ sampled from $n$ independent $d^k$-dimensional unitary ensemble $(\mathcal{U}_1,\cdots,\mathcal{U}_n)$.
    \item Measure the state $U \rho^{\otimes k} U^\dagger$ on the computational basis with result $\ket{\mathbf{b}}\!\bra{\mathbf{b}}$. Restore the classical description of circuit $U$ and length-n bit string $\mathbf b$.
    \item \textit{Virtually} computes the shadow matrix $\hat \rho^{\otimes k}:= \mathcal{M}^{-1}(U^\dagger \ket{\mathbf{b}}\!\bra{\mathbf{b}} U)$, where $\mathcal{M}$ is the entanglement-breaking linear map $\rho^{\otimes k} \to \mathbb{E}_{U \sim \mathcal{U}} \sum_{\mathbf b}\bra{\mathbf b}U \rho^{\otimes k} U^\dagger \ket{\mathbf{b}} \cdot  U^\dagger\ket{\mathbf b}\!\bra{\mathbf b} U$.
\end{itemize}
%%%%%%%%%%%%%%%%%%%%%%%%%%%%%%%%%%%%%%
 As shown in Lemma~\ref{lemma:single_shadow_unbiased}, to estimate average $\text{tr}(O \hat{\rho}^{\otimes k})$, ones classically post-processing those sampled classical shadow. 
\begin{lemma}[Proof of unbiasedness]\label{lemma:single_shadow_unbiased}
    The estimator $\hat{o}^{(k)} := \text{tr}(O \hat{\rho}^{\otimes k})$ is an unbiased estimator for the expectation $\text{tr}(O {\rho}^{\otimes k})$, where $\hat{\rho}^{\otimes k}$ is sampled according to multi-copy local shadows above.
\end{lemma}
\begin{proof}
    \begin{align}
        \mathbb{E}_{\hat \rho} \text{tr}(O \hat{\rho}^{\otimes k}) &= \mathbb{E}_{U \sim \mathcal{U}}\mathbb{E}_{\mathbf{b} \sim U \rho^{\otimes k} U^\dagger} \text{tr}\left[O \mathcal{M}^{-1}(U^\dagger \ket{\mathbf{b}}\!\bra{\mathbf{b}} U)\right]\\
        &=  \text{tr}\left[O \mathcal{M}^{-1}(\mathbb{E}_{U\sim \mathcal{U}}\mathbb{E}_{\mathbf{b} \sim U \rho^{\otimes k}  U^\dagger}U^\dagger \ket{\mathbf{b}}\!\bra{\mathbf{b}} U)\right]=\text{tr}(O \rho^{\otimes k}).
    \end{align}
\end{proof}
The explicit form of channel $\mathcal{M}$ rely on the ensembles $\mathcal{U}_i$.

\subsection{Independent local 3-design}
\begin{definition}[Unitary $t$-design]\label{def:t_design}
    A unitary $t$-design $\mathbf{D}_t(d)$ on a $d$-dimensional Hilbert space is a distribution over the unitary group $U(d)$ such that the $t$-th moment matches that of the Haar measure: $\mathbb{E}_{U \sim \mathbf{D}_t(d)} U^{\otimes t} (\cdot) U^{\dagger \otimes t} = \mathbb{E}_{U \sim \mathrm{Haar}(d)} U^{\otimes t} (\cdot) U^{\dagger \otimes t}$.
\end{definition}
\begin{lemma}[State $t$-design averaging]\label{lemma:state_trajectory}
    If $\ket{\psi}$ is drawn from a state $t$-design on $\mathbb{C}^d$ (equivalently, $\ket{\psi} = U\ket{0}$ with $U$ from a unitary $t$-design), then
    $\mathbb{E}\left[\left(\ket{\psi}\!\bra{\psi}\right)^{\otimes t}\right] = \frac{1}{d^{\uparrow t}} \sum_{\pi \in S_t} \mathbb{P}_{\pi}^{(d)}$,
    where $d^{\uparrow t} = d(d+1)\cdots(d+t-1)$ is the rising factorial and $\mathbb{P}_{\pi}^{(d)}$ denotes the permutation operator on $(\mathbb{C}^d)^{\otimes t}$.
\end{lemma}
In this section we consider the case that $(\mathcal{U}_1,\cdots,\mathcal{U}_n)$ being $n$ independent unitary 3-design.
This i.i.d. unitary 3-design setting includes the Clifford randomized measurement and the randomized Pauli basis measurement proposed in~\cite{huang2020predicting}.
The result of $\mathcal{M}$ is firstly summarized in Lemma~\ref{lemma:onsite_2design}.
%%%%%%%%%%%Lemma%%%%%%%%%%%Lemma%%%%%%%%%%%Lemma%%%%%%%%%%%Lemma
\begin{lemma}[On site 2-design unitary ensemble]\label{lemma:onsite_2design}
    A $k$-site 3-design unitary ensembles $\mathbf{D}_{3}(d^k)^{\otimes n}$ separates $n\cdot k$ numbers of qudits in to $n$ parts, and independently implements $k$-qudit 2-design unitary, as defined in Definition~\ref{def:t_design}, within each part. In other words, $\mathbf{D}_{3}(d^k)^{\otimes n}:=\{\bigotimes_{i  =1}^n U_i| U_i \sim \mathbf{D}_{2}(d^k)\}$.
    This ensemble induces the channel
    \begin{align}
        \mathcal{M}(\rho) = (\mathcal{D}_{1/(2^k + 1)})^{\otimes n}(\rho), \text{with~} \mathcal{D}_{1/(2^k + 1)}(\cdot) = \frac{1}{2^k +1}(\cdot) + \frac{\text{tr}(\cdot)}{2^k + 1}{\mathbb{I}(d^k)}.
    \end{align}
    Specifically, the single-copy global 2-design unitary ensemble takes the case $n = 1$, $U \sim \mathbf{D}_{3}(d^k)$,and the channel $     \mathcal{M}(\rho) ={(d^k + 1)^{-1}} \rho + {d^k{(d^k + 1)^{-1}}  \text{tr}\rho \cdot }({\mathbb{I}(d^k)}/{d^k})$.
    The single-copy local 2-design unitary ensemble takes the case $k = 1$, $U \sim \mathbf{D}_{3}(d)^{\otimes n}$,and the channel $ \mathcal{M}(\rho) = {(d + 1)}^{-1} \rho + {d({d + 1})^{-1} \text{tr}\rho} \cdot ({\mathbb{I}(d)}/{d})$.
\end{lemma}
\begin{proof}
    For tensor product inputs $\otimes_{i = 1}^n \rho_i, \rho_i \in \mathcal{D}(\mathcal{H}(d^k))$, the channel $\mathcal{M}$ in this case became
\begin{align}
    \mathcal{M}(\otimes_{i = 1}^n \rho_i) 
    &=\mathbb{E}_{\substack{U = \bigotimes_{i=1}^n\text{U}_i \\ \text{U}_i \sim \mathbf{D}_{3}(d^k)}} \sum_{\mathbf{r} \in \mathbb{F}_d^{kn}} \bra{{\mathbf r}} \text{U} (\otimes_{i = 1}^n \rho_i) \text{U}^{\dagger} \ket{{\mathbf r}} \cdot \text{U}^{\dagger} \ket{{\mathbf r}}\!\bra{{\mathbf r}} \text{U} \notag\\
    &=\mathbb{E}_{\substack{U = \bigotimes_{i=1}^n\text{U}_i \\ \text{U}_i \sim \mathbf{D}_{3}(d^k)}} \sum_{\mathbf{r} \in \mathbb{F}_d^{kn}} \text{tr}_{\bar A} \left( \left(\otimes_{i = 1}^n \rho_i^{(\bar A)} \otimes \mathbb{I}^{(A)}(d^{kn})\right)\cdot  (\text{U}^{\dagger} \ket{{\mathbf r}}\! \bra{{\mathbf r}} \text{U})^{(\bar A)} \otimes (\text{U}^{\dagger} \ket{{\mathbf r}}\!\bra{{\mathbf r}} \text{U})^{(A)} \right)\notag\\
    &= \text{tr}_{\bar A}\left( \left(\otimes_{i = 1}^n \rho_i^{(\bar A)}\otimes \mathbb{I}^{(A)}(d^{kn})\right) \cdot \bigotimes_{i = 1}^n \sum_{\mathbf{r}_i \in \mathbb{F}_d^k} \mathbb{E}_{\substack{\text{U}_i \sim \mathbf{D}_{3}(d^k)}}  (\text{U}_i^{\dagger} \ket{\mathbf{r}_{i}}\bra{\mathbf{r}_{i}} \text{U}_i)^{(\bar A)} \otimes (\text{U}_i^{\dagger} \ket{\mathbf{r}_{i}}\!\bra{\mathbf{r}_{i}} \text{U}_i)^{(A)} \right)\notag\\
    &=  \text{tr}_{\bar A}\left( \left(\otimes_{i = 1}^n \rho_i^{(\bar A)}\otimes \mathbb{I}^{(A)}(d^{kn})\right) \cdot \bigotimes_{i = 1}^n \sum_{\mathbf{r}_{i}\in \mathbb{F}_d^{k}} \Phi_2(\ket{\mathbf{r}_{i}}\!\bra{\mathbf{r}_{i}}^{\otimes 2})^{(\bar A A)} \right)\notag \\
    &= \text{tr}_{\bar A}\left( \left(\otimes_{i = 1}^n \rho_i^{(\bar A)}\otimes \mathbb{I}^{(A)}(d^{kn})\right) \cdot \bigotimes_{i = 1}^n d^{k} \frac{\mathbb{I}^{(\bar A)}(d^k)\otimes \mathbb{I}^{(A)}(d^k)+ \mathbb{P}_{(12)}^{d,(\bar A A)}}{(d^k+1) \cdot d^k}\right).\label{eq:k_copy_shadow_channel}
\end{align}
Use the fact that  $\text{tr}_{\bar A}(\mathbb{P}_{(12)}^{d,(\bar A A)}\cdot \rho^{(\bar A)}\otimes \mathbb{I}^{(A)}) = \rho^{(A)}$, we get the channel output
\begin{align}
    \mathcal{M}(\otimes_{i = 1}^n \rho_i) 
    &= \bigotimes_{i = 1}^n \text{tr}(\rho_i) \frac{\mathbb{I}(d^k)}{d^k + 1} + \frac{1}{d^k +1}\rho_i = \mathcal{D}^{\otimes n}_{{1}/{(d^k +1)}}(\otimes_{i = 1}^n \rho_i).
\end{align}
By the linearity of $\mathcal{M}$, we  conclude that its inverse $\mathcal{D}_{{1}/{(d^k +1)}}^{-1}(\rho) = (d^k + 1)\rho - \text{tr}\rho \cdot \mathbb{I}(d^k) $ constitute an unbiased estimator of $\rho \in \mathcal{D}(\mathcal{H}(d^{nk}))$ through the transformation
\begin{align}
     \hat \rho_{U,\mathbf b} =\mathcal{M}^{-1}(\text{U}^\dagger \ket{\mathbf b}\!\bra{\mathbf b} \text{U})=\bigotimes_{i =1}^n\left((d^k +1)\text{U}_i^{\dagger} \ket{\mathbf{b}_i}\!\bra{\mathbf{b}_i} \text{U}_i - \mathbb{I}(d^k)\right).
\end{align}
\end{proof}
%%%%%%%%%%%%%%%%%%%%%%%%%%%%%%
Following this, we calculate the estimation variance for i.i.d. 3-design unitary ensemble.
In this case, we put $\rho^{\otimes k}$ in quantum memory and then implement randomized k-qubits Clifford gate $U = {U_i}^{\otimes n}$ with $U_i \in \mathbf{D}_{3}(2^k)$ the $k$-qubits gate and forms $3$-design. After this, we measured each qubit in the computational basis. 
Here below, we consider the qubit case $d = 2$ only.
%%%%%%%%Lemma%%%%%%%%Lemma%%%%%%%%Lemma%%%%%%%%Lemma
\begin{lemma}[Local k-copy 3-design shadow]\label{lemma:LC_shadow_kcopy}
    In the case of local k-copy randomized measurements on $\rho^{\otimes k}$, let the observable $O^{(k)}$ acting on $n\cdot k$ qubits have decomposition $O^{(k)} = \sum_{W} \alpha_{W} W$, where $W = \bigotimes_{i=1}^n W_i$ and $W_i \in \{I,X,Y,Z\}^{\otimes k}$. Denoting $\mathcal{C}(W)$ as the set of $k$-qubit Pauli operator that commute with $W$ and $\delta_{e,S} = 1$ iff $e\in S$, we define a function $R:\mathcal{P}_k^{\otimes n} \to \mathbb{R}$ as
    \begin{align}\label{eq:R_expression}
        R(W,W') = \prod_{i = 1}^n r(W_i,W'_i), r(W_i,W'_i) := \begin{cases}
            1 &\text{if~}W_i = \mathbb{I}_k \text{~or~}W'_i = \mathbb{I}_k ,\\
            (2^k+1) \cdot \frac{2^{k-1}\delta_{W_i,W'_i} +\delta_{W_i, \mathcal{C}(W'_i)\backslash \mathbb{I}_k } }{2^{k-1}+1} &\text{if~}W_i \neq \mathbb{I}_k \text{~and~} W'_i \neq \mathbb{I}_k .
        \end{cases}
    \end{align}
    we have $\langle (\hat{o}^{(k)})^2 \rangle_{\rho} =\text{tr} \left(\rho^{\otimes k} \cdot \mathbf{V}(O^{(k)}) \right)$, where $\mathbf{V}(O^{(k)}) := \sum_{W,W'}\alpha_W \alpha_{W'} R(W,W')W W'$. Specifically, if we take a local operator $O^{(k)} = \bigotimes_{i = 1}^n \left(\sum_{W_i \in \mathcal{P}_k}\alpha_{W_i} W_i\right)$, we have $\mathbf{V}(O^{(k)}) :=\bigotimes_{i = 1}^n  \sum_{W_i,W_i'}\alpha_{W_i} \alpha_{W_i'} r(W_i,W_i')W_i W_i'$.
\end{lemma}
\begin{proof}
By Lemma~\ref{lemma:onsite_2design}, we conclude that its inverse is $\mathcal{D}_{1/(2^k + 1)}^{-1}(\cdot) = (2^k + 1)(\cdot) - \text{tr}(\cdot) \mathbb{I}_k $ constitute an unbiased estimator of $\rho^{\otimes k}$ through the transformation $\hat{\rho}^{\otimes k} =\mathcal{M}^{-1}(U^\dagger \ket{\mathbf r}\!\bra{\mathbf r} U)=\bigotimes_{i =1}^n\left((2^k + 1)U_i^{\dagger} \ket{\mathbf{r}_{i}}\bra{\mathbf{r}_{i}} U_i - \mathbb{I}_k\right).$
The second moment is given by $  \langle (\hat{o}^{(2)})^2 \rangle_{\rho} =\mathbb{E}_{U \sim \mathbf{D}_3(2^k)^{\otimes n}}\sum_{\mathbf r}\bra{\mathbf r}U \rho^{\otimes k} U^\dagger \ket{\mathbf{r}} \bra{\mathbf r}U\mathcal{M}^{-1}(O^{(k)}) U^\dagger\ket{\mathbf r}^2$.
    Similar to the proof in Lemma~\ref{lemma:LC_shadow_kcopy}, we decomposed the operator $O^{(k)}$ on the Pauli product basis $W_i \in \{I,X,Y,Z\}^{\otimes k}$ that supports on those k qubits located at site $i$. By denoting $W = W_1\cdots W_n$, we have $O^{(k)} = \sum_{W} \alpha_{W} W$. According to Lemma~\ref{lemma:onsite_2design}, the inverse channel $\mathcal{M}^{-1}$ acts locally as $\mathcal{M}_i^{-1}(W_i) = g(W_i) W_i$, where the scaling factor $g(W_i)$ is determined by $g(W_i) = \begin{cases} (2^k+1) & \text{if~} W_i \neq \mathbb{I}_k \\ 1 &\text{if~} W_i = \mathbb{I}_k \end{cases}.$
    Use this decomposition and again supposed $\rho^{\otimes k} = \sum_{s} \beta_s \bigotimes_{i=1}^n \rho_i^{(s)}$, where $ \rho_i^{(s)} \in \mathcal{D}(\mathcal{H}_k)$, the result involves averaging over the k-qubit 3-design $U_i \sim \mathbf{D}_3(2^k)$ as:
    \begin{align}
        \langle (\hat{o}^{(k)})^2 \rangle_{\rho} &= \sum_{W,W'}\alpha_{W}\alpha_{W'}\sum_{s}\beta_s\prod_{i = 1}^n g(W_i)g(W'_i)\sum_{\mathbf{r}_i \in \{0,1\}^k} \left[({\rho}^{(s)}_i \otimes W_i \otimes W'_i) \cdot \mathbb{E}_{U_i} U_i^{\otimes 3}\ket{\mathbf{r}_i}\!\bra{\mathbf{r}_i}^{\otimes 3}U_i^{\otimes 3,\dagger}\right] \\
        &= \sum_{W,W'}\alpha_{W}\alpha_{W'}\sum_s \beta_s\prod_{i = 1}^n 2^k p \cdot g(W_i)g({W'}_i) \Bigl[ \text{tr}(W_i)\text{tr}(W'_i)\text{tr}({\rho}^{(s)}_i) + \text{tr}({\rho}^{(s)}_i W_i)\text{tr}(W'_i) \notag\\
        & \qquad\qquad\qquad + \text{tr}(W_i)\text{tr}({\rho}^{(s)}_i W'_i)+\text{tr}({\rho}^{(s)}_i)\text{tr}(W_iW'_i) + \text{tr}(W_i {\rho}^{(s)}_i W'_i) + \text{tr}(W'_i {\rho}^{(s)}_i W_i) \Bigr]. \label{eq:mid_sqr_ave_2copy}
    \end{align}
    
    Here we use the fact that the term $\ket{E} :=U_i\ket{\mathbf{r}_i}$ is average over the 3-design states, yielding terms based on the formula $\mathbb{E} [E^{\otimes 3}] ={d^{-1}(d+1)^{-1}(d+2)^{-1}} \sum_{\pi \in S_3} \mathbb{P}_{\pi}^{(d)}$ according to Lemma~\ref{lemma:state_trajectory}. With $d=2^k$, the prefactor per site is $p = 1/(2^k \cdot (2^k+1) \cdot (2^k+2))$.  We evaluate the expression by considering terms for each site $i$, analogous to the single-copy proof of Lemma~\ref{lemma:LC_shadow_kcopy}. Depending on three cases, we calculate:
    \begin{align*}
        \langle (\hat{o}^{(2)})^2 \rangle_{\rho} &= 
        \sum_{W,W'}\alpha_{W}\alpha_{W'}\sum_{s}\beta_s\prod_{i = 1}^n 2^k p \cdot \begin{cases}
         1\cdot 1\cdot [2^{2k}+3\cdot 2^k+2]\text{tr}{\rho}^{(s)}_i  &\text{if~} W_i=W_i' = \mathbb I_k,\\
         1 \cdot (2^k+1) \cdot[2^k + 2]\cdot \text{tr}({\rho}^{(s)}_i W'_i)&\text{if~} W_i = \mathbb I_k \neq W'_i,\\
         (2^k+1) \cdot 1 \cdot[2^k + 2]\cdot \text{tr}({\rho}^{(s)}_i W_i)&\text{if~} W'_i = \mathbb I_k \neq W_i,\\
         (2^k+1)^2 \cdot[2^k\delta_{W_i, W'_i}\text{tr}{\rho}^{(s)}_i  + \text{tr}({\rho}^{(s)}_i\{W_i,W_i'\})]&\text{if~} W_i, W'_i \neq \mathbb I_k.
     \end{cases}\\
    %  &=\sum_{W,W'}\alpha_{W}\alpha_{W'}\sum_{s}\beta_s\prod_{i = 1}^n \begin{cases}
    %     \text{tr}{\rho}^{(s)}_i  &\text{if~} W_i=W_i' = \mathbb I_k,\\
    %      \text{tr}({\rho}^{(s)}_i W'_i)&\text{if~} W_i = \mathbb I_k\neq W'_i,\\
    %      \text{tr}({\rho}^{(s)}_i W_i)&\text{if~} W'_i = \mathbb I_k\neq W_i,\\
    %     \frac{2^k+1}{2^k+2} \cdot\text{tr}\left[(2^k\delta_{W_i, W'_i}\mathbb{I}_k  +2\delta_{W_i, \mathcal{C}(W'_i)\backslash \mathbb{I}_k}  W_iW_i'){\rho}^{(s)}_i\right]&\text{if~} W_i, W'_i \neq \mathbb I_k.
    %  \end{cases}\\
     &= \sum_{W,W'}\alpha_{W}\alpha_{W'}\cdot \sum_{s}\beta_s\prod_{i = 1}^n \begin{cases}
        \text{tr}({\rho}^{(s)}_i W_iW'_i)  &\text{if~} W_i=W_i' = \mathbb I_k,\\
         \text{tr}({\rho}^{(s)}_i W_i W'_i)&\text{if~} W_i = \mathbb I_k\neq W'_i,\\
         \text{tr}({\rho}^{(s)}_i W_i W'_i)&\text{if~} W'_i = \mathbb I_k\neq W_i,\\
         \frac{2^k+1}{2^k+2} \cdot\text{tr}\left[{\rho}^{(s)}_i(2^k\delta_{W_i, W'_i}+2\delta_{W_i, \mathcal{C}(W'_i)\backslash \mathbb{I}_k})W_iW_i'\right]&\text{if~} W_i, W'_i \neq \mathbb I_k.
     \end{cases}\\
     &=\text{tr}\left( \rho^{\otimes k}\cdot \sum_{W,W' \in \mathcal{P}_k^{\otimes n}} \alpha_{W}\alpha_{W'} R(W,W') W W'\right).
    \end{align*}
    To prove the second equality, we first defined $\mathcal{C}(W)$ as the set of k-qubit Pauli operator that commute with $W$ and $\delta_{e,S} = 1$ iff $e\in S$, otherwise $\delta(e,S) = 0$. Then we use the fact that
    \begin{align}
        \forall W_i, W'_i \neq \mathbb I_k, \text{tr}({\rho}^{(s)}_i\{W_i,W_i'\}) = \begin{cases}0 \!\!\!&\text{if~} W'_i \notin \mathcal{C}(W_i)\backslash \mathbb{I}_k\\2\text{tr}(\tilde{\rho}^{(s)}_i W_iW'_i)  \!\!\!&\text{if~} W_i \in \mathcal{C}(W'_i)\backslash \mathbb{I}_k \end{cases} = 2\delta_{W_i, \mathcal{C}(W'_i)\backslash \mathbb{I}_k} \text{tr}(\tilde{\rho}^{(s)}_i W_iW'_i),
    \end{align}
In the last equality, we use the definition
\begin{align}
    R(W,W') = \prod_{i = 1}^n r(W_i,W'_i),  r(W_i,W'_i) := \begin{cases}
        1 &\text{if~}W_i = \mathbb{I}_k \text{~or~}W'_i = \mathbb{I}_k ,\\
        (2^k+1) \cdot \frac{2^{k-1}\delta_{W_i,W'_i} +\delta_{W_i, \mathcal{C}(W'_i)\backslash \mathbb{I}_k } }{2^{k-1}+1} &\text{if~}W_i \neq \mathbb{I}_k \text{~and~} W'_i \neq \mathbb{I}_k .
    \end{cases}
\end{align}
\end{proof}
%%%%%%%%%%%%%%%%%%%%%%%%%%%%%%
The inherient difference between the single-copy local Clifford strategy and the multi-copy one reflected on the set of $k$-qubit Pauli operators that commute with a given $k$-qubit Pauli operator $W$, which we denoted as $ \mathcal{C}(W)$.
In the single-copy case, the center of $P_i \in \{X,Y,Z\}$ satisfied $\mathcal{C}(P_i) = \{P_i, \mathbb{I}\}$.
     However, for example, for the two-qubit Pauli operators $XX$, extra terms such as $ZZ$ involved.
     We denote the $n$-qubit Pauli operators as $\sigma_{\mathbf{r}} = \bigotimes_{i=1}^n \sigma_{{r}_{x,i},{r}_{z,i}} \in \mathcal{P}_n$ with $\mathbf{r} = (\mathbf{r}_x = ({r}_{x,1},\cdots,{r}_{x,n}),\mathbf{r}_z = ({r}_{z,1},\cdots,{r}_{z,n})) \in (\mathbb{F}_2^n)^2$ and $\{\sigma_{00},\sigma_{01},\sigma_{10},\sigma_{11}\} = \{\mathbb{I}, Z,X,Y = iXZ\}$. We denote $a \oplus b = a + b \mathrm{~mod~}2 = a + b - 2a\cdot b$ and $\mathbf{a} \odot \mathbf{b} = (a_1\cdot b_1,\cdots,a_n\cdot b_n)$ with $a,b \in \{0,1\}$.
In the following, we follow the Pauli multiplication rule as proved below 
     \begin{align}
         \sigma_{\mathbf{s}}\sigma_{\mathbf{s}'} &= \bigotimes_{j = 1}^n i^{s_{x,j} s_{z,j}} X^{s_{x,j}}Z^{s_{z,j}} \cdot i^{s'_{x,j} s'_{z,j}} X^{s'_{x,j}}Z^{s'_{z,j}} =  \bigotimes_{j = 1}^n X^{s_{x,j}+s_{x,j}'}Z^{s_{z,j}+s_{z,j}'} (-1)^{s_{z,j}\cdot s_{x,j}'} \cdot i^{s_{x,j}s_{z,j}+s_{x,j}'s_{z,j}'} \notag \\
         &= \sigma_{\mathbf{s}\oplus \mathbf{s}'} (-i)^{(\mathbf{s}_x+ \mathbf{s}_x'-2\mathbf{s}_x\odot\mathbf{s}_x')(\mathbf{s}_z+\mathbf{s}_z' - 2\mathbf{s}_z\odot\mathbf{s}_z')} i^{\mathbf{s}_x\mathbf{s}_z + \mathbf{s}_x'\mathbf{s}_z'}(i)^{2\mathbf{s}_z\mathbf{s}_x'} \\
         &=i^{\mathbf{s}_x'\mathbf{s}_z - \mathbf{s}_z'\mathbf{s}_x} (-1)^{(\mathbf{s}_x + \mathbf{s}_x') \cdot (\mathbf{s}_z \odot \mathbf{s}_z') + (\mathbf{s}_z + \mathbf{s}_z') \cdot (\mathbf{s}_x \odot \mathbf{s}_x')}\sigma_{\mathbf{s}+ \mathbf{s}'} := i^{\mathbf{s}' J \mathbf{s}} (-1)^{d(\mathbf{s},\mathbf{s}')}\sigma_{\mathbf{s}+ \mathbf{s}'} .
     \end{align}

Defined $[\mathbf{s}', \mathbf{s}]:= \mathbf{s}_x'\cdot \mathbf{s}_z - \mathbf{s}_z'\cdot \mathbf{s}_x \text{~mod~} 2$ as the binary symplectic product. Use this rule, we calculate  $ \sigma_{\mathbf{s}}\cdot \sigma_{\mathbf{s}'}=i^{2\mathbf{s}_x'\cdot \mathbf{s}_z - 2\mathbf{s}_z'\cdot \mathbf{s}_x }  (\sigma_{\mathbf{s}'} \cdot \sigma_{\mathbf{s}}) \Rightarrow [\mathbf{s}, \mathbf{s}'] = 0,$ and show that $\mathcal{C}(\sigma_{\mathbf{u}} )\backslash \mathbb{I}_k= \{\sigma_{\mathbf{v}}|\mathbf{v}\in \mathbb{F}_2^k,\mathbf{v}\neq \mathbf{0},[\mathbf{v},\mathbf{u}] = 0\}.$
This means that $\delta_{\sigma_{\mathbf{s},\mathcal{C}(\sigma_{\mathbf{s}'} )\backslash \mathbb{I}_k}} = \delta_{\mathbf{s}, \{\mathbf{v}\in \mathbb{F}_2^k|\mathbf{v}\neq \mathbf{0},[\mathbf{v},\mathbf{s}'] = 0\}} = {\left(1 + (-1)^{[\mathbf{s},\mathbf{s}']}\right)}/{2}$ for $\mathbf{s},\mathbf{s}' \neq \mathbf{0}$ and 
\begin{align}\label{eq:lc_Pauli_rescale}
r(\sigma_{\mathbf{s}},\sigma_{\mathbf{s}'})& := \begin{cases}
            1 &\text{if~}\sigma_{\mathbf{s}} = \mathbb{I}_k \text{~or~}\sigma_{\mathbf{s}'} = \mathbb{I}_k ,\\
            (2^k+1) \cdot \frac{2^{k-1}\delta_{\sigma_{\mathbf{s}},\sigma_{\mathbf{s}'}} +\delta_{\sigma_{\mathbf{s}}, \mathcal{C}(\sigma_{\mathbf{s}'})\backslash \mathbb{I}_k } }{2^{k-1}+1} &\text{if~}\sigma_{\mathbf{s}} \neq \mathbb{I}_k \text{~and~} \sigma_{\mathbf{s}'}\neq \mathbb{I}_k .
        \end{cases}\\
&= \begin{cases}
            1 &\text{if~}\mathbf{s} = \mathbf{0}\text{~or~}\mathbf{s}' = \mathbf{0} ,\\
            \frac{2^k+1}{2^{k}+2} \cdot \left[2^{k}\delta_{\mathbf{s},\mathbf{s}'} +1+ (-1)^{[\mathbf{s}',\mathbf{s}]}\right]  &\text{if~}\mathbf{s} \neq \mathbf{0}\text{~and~}\mathbf{s}' \neq \mathbf{0} .
        \end{cases}
\end{align}

A generalization of the variance operator $\mathbf{V}(O^{(k)})$ is to applied it to two operators $O_\alpha,O_\beta$, as defined below:
\begin{definition}[Covariance operator]
    \label{def:cov_oper}
    Given two operators with decomposition $O_\alpha = \sum_{W} \alpha_{W} W, O_\beta = \sum_{W} \beta_{W} W$ with $W \in \{\mathbb{I},X,Y,Z\}^{\otimes k}$, the we defined the covariance operator as 
    \begin{align}
        \mathbf{V}(O_\alpha,O_\beta) = \sum_{W,W'} \alpha_{W} \beta_{W'} R(W,W') WW',
    \end{align}
    where $R(W,W')$ is defined in Lemma~\ref{lemma:LC_shadow_kcopy} before.
\end{definition}
The covariance operator is important due to the Lemma below
\begin{lemma}
    \label{lemma:covariance_decomp}
    Given an k-copy operator $O = \lambda_{\alpha} O_{\alpha} + \lambda_{\beta} O_{\beta}$, we have \begin{align}
        \mathbf{V}(O,O) = \lambda_{\alpha}^2 \mathbf{V}(O_{\alpha},O_{\alpha}) + 2 \lambda_{\alpha} \lambda_{\beta} \mathbf{V}(O_{\alpha},O_{\beta}) + \lambda_{\beta}^2 \mathbf{V}(O_{\beta},O_{\beta}).
    \end{align}
\end{lemma}
\begin{proof}
    Supposed the decomposition $O_\alpha = \sum_{W} \alpha_{W} W, O_\beta = \sum_{W} \beta_{W} W$, we noted that 
    \begin{align}
        \mathbf{V}(O,O) &=  \sum_{W,W'} (\lambda_{\alpha}\alpha_{W} + \lambda_{\beta} \beta_{W})(\lambda_{\alpha}\alpha_{W'} + \lambda_{\beta} \beta_{W'}) R(W,W') WW',\\
        &=  \sum_{W,W'} \lambda^2_{\alpha}\alpha_{W}\alpha_{W'}R(W,W') WW' +\sum_{W,W'}  \lambda_{\alpha}\lambda_{\beta} \alpha_{W}\beta_{W'}R(W,W') WW' \notag\\
        & +\sum_{W,W'}  \lambda_{\beta} \lambda_{\alpha}\beta_{W}\alpha_{W'}R(W,W') WW' + \sum_{W,W'} \lambda^2_{\beta}\beta_{W}\beta_{W'}R(W,W') WW'\\
        &=\lambda_{\alpha}^2 \mathbf{V}(O_{\alpha},O_{\alpha}) + 2 \lambda_{\alpha} \lambda_{\beta} \mathbf{V}(O_{\alpha},O_{\beta}) + \lambda_{\beta}^2 \mathbf{V}(O_{\beta},O_{\beta}).
    \end{align}
    Here at the last line, we use the definition~\ref{def:cov_oper} and the fact that $R(W,W') = R(W',W)$. 
\end{proof}

\paragraph{Application to $\text{tr}(O \rho^2)$.}
A critical application of this local multi-copy shadow is to use it to estimate the non-linear term $\text{tr}(O \rho^k)$ for arbitrary operator $O \in \text{Herm}(\mathcal{H}_n)$.
In the case of $k  =2$, $\text{tr}((O \otimes \mathbb{I}_n) \text{SWAP}^{\otimes n}\cdot (\rho \otimes \rho))$ is an unbiased estimator for $\text{tr}(O \rho^2)$. If we decomposed the operator $O$ on the Pauli basis, it turns out $(O \otimes \mathbb{I}_n)\text{SWAP}^{\otimes n} = \sum_{\mathbf{r} \in \mathbb{F}_2^{2n}} (\sigma_{\mathbf{r}} \otimes  \mathbb{I}_n)\text{SWAP}^{\otimes n}:= \sum_{\mathbf{r} \in \mathbb{F}_2^{2n}} O_{\mathbf{r}} $.
Here below we use Lemma~\ref{lemma:LC_shadow_kcopy} to evaluate the covariance operator of $O_{\mathbf{r}}$ explicitly. 
%%%%%%%%%%%Lemma%%%%%%%%%%%Lemma%%%%%%%%%%%Lemma%%%%%%%%%%%Lemma
\begin{lemma}[Variance of two-copy Pauli]\label{lemma:covar_Pauli}
    If we take $O_{\mathbf{r}}= (\sigma_{\mathbf{r}} \otimes \mathbb{I}_n) \cdot \text{SWAP}^{\otimes n}$ for Pauli operator $\sigma_{\mathbf{r}}$ and denote $\text{wt}(\sigma_{\mathbf{r}})$ as the weight of this Pauli operator, then we have $\mathbf{V}(O_{\mathbf{r}}, O_{\mathbf{r}'}) = \bigotimes_{j = 1}^n v(\mathbf{r}_j, \mathbf{r}_j') $ with
    \begin{align}\label{eq:v_full_operator}
        v(\mathbf{r}_j, \mathbf{r}_j') =
        \left\{\begin{aligned}
            &\frac{5}{6}\left(\sigma_{\mathbf{r}_j}\otimes \sigma_{\mathbf{r}'_j} +\sigma_{\mathbf{r}'_j}\otimes \sigma_{\mathbf{r}_j} \right)&\text{~if~} \mathbf{r}_j, \mathbf{r}_j'  \neq \mathbf{0},\\
            &\frac{11\mathbb{I}_4 + 4 \Pi_{\mathrm{asym}}}{3} &\text{~if~} \mathbf{r}_j=\mathbf{r}_j'= \mathbf{0}, \\
            &\left(\frac{1}{2} \Pi_{\mathrm{sym}} + \frac{7}{6} \Pi_{\mathrm{asym}}\right)(\mathbb{I} \otimes \sigma_{\mathbf{r}_j})+ \frac{5}{6}\sigma_{\mathbf{r}_j}\otimes\mathbb{I} &\text{~if~} \mathbf{r}_j' = \mathbf{0},  \mathbf{r}_j\neq \mathbf{0},\\
            &\left(\frac{1}{2} \Pi_{\mathrm{sym}} + \frac{7}{6} \Pi_{\mathrm{asym}}\right)(\mathbb{I} \otimes \sigma_{\mathbf{r}'_j}) + \frac{5}{6}\sigma_{\mathbf{r}'_j}\otimes\mathbb{I}&\text{~if~} \mathbf{r}_j = \mathbf{0},  \mathbf{r}'_j\neq \mathbf{0}.
        \end{aligned}\right.
    \end{align}
    In particular, the diagonal element gives $v(\mathbf{r}_j,\mathbf{r}_j) = ({5}/{3})\sigma_{\mathbf{r}_j}^{\otimes 2}$ when $\mathbf{r}_j \neq \mathbf{0}$ and $v(\mathbf{0},\mathbf{0}) = ({11\mathbb{I}_4 + 4 \Pi_{\mathrm{asym}}})/{3}$ at identity sites. Note that the $\Pi_{\mathrm{asym}}$ terms do \emph{not} vanish on $\rho^{\otimes 2}$ for entangled $\rho$: while $\rho^{\otimes 2}$ is globally symmetric under copy exchange, it is not contained in $\bigotimes_j \mathrm{Sym}^2_j$ when $\rho$ has entanglement across qubit sites. For product states $\rho = \bigotimes_l \rho_l$, the per-site state $\rho_l^{\otimes 2} \in \mathrm{Sym}^2(\mathbb{C}^2)$ so $\Pi_{\mathrm{asym}}$ terms vanish and the second moment simplifies to $\mathrm{tr}[\mathbf{V}(O_{\mathbf{r}},O_{\mathbf{r}})\rho^{\otimes 2}] = ({11}/{3})^{n}({5}/{11})^{\mathrm{wt}(\sigma_{\mathbf{r}})}$. For entangled states the second moment can be larger: since $v(\mathbf{0},\mathbf{0})$ has eigenvalues $11/3$ on $\mathrm{Sym}^2$ and $5$ on the antisymmetric subspace, the operator norm is $\|\mathbf{V}(O_\mathbf{0},O_\mathbf{0})\| = 5^n$.
\end{lemma}
\begin{proof}
    We consider the observable $O_{\mathbf{r}}= (\sigma_{\mathbf{r}} \otimes \mathbb{I}_n) \cdot \text{SWAP}^{\otimes n}$ with $\mathbf{r} = (\mathbf{r}_1, \cdots, \mathbf{r}_n) \in (\mathbb{F}_2^2)^n$ and rewrite the observable according to the fact $\text{SWAP}^{\otimes n} = 2^{-n}\sum_{\mathbf{p} \in \mathbb{F}_2^{2n}} \sigma_{\mathbf{p}} \otimes \sigma_{\mathbf{p}}$ as:
    \begin{align}
        O_{\mathbf{r}}=  2^{-n}\sum_{\mathbf{p} \in (\mathbb{F}_2^{2})^n} \sigma_{\mathbf{r}}\sigma_{\mathbf{p}}\otimes \sigma_{\mathbf{p}},\quad O_{\mathbf{r}'}=  2^{-n}\sum_{\mathbf{p}' \in (\mathbb{F}_2^{2})^n} \sigma_{\mathbf{r}'}\sigma_{\mathbf{p}'}\otimes \sigma_{\mathbf{p}'} 
    \end{align}
    This means that the operator $O_{\mathbf{r}},O_{\mathbf{r}'}$ satisfied the local decomposition $\left\{\begin{aligned}W_j \propto \sigma_{\mathbf{r}_j + \mathbf{p}_j} \otimes \sigma_{\mathbf{p}_j}\\ W'_j \propto \sigma_{\mathbf{r}'_j + \mathbf{p}'_j} \otimes \sigma_{\mathbf{p}'_j}\end{aligned}\right.$
    Here below we abbreviate the foot label $j$ for simplicity.
    Using Definition~\ref{def:cov_oper}, we calculate the covariance operator as
    \begin{align}
        &\mathbf{V}(O_{\mathbf{r}}, O_{\mathbf{r}'}) = 4^{-n} \bigotimes_{j=1}^n \left(\sum_{\mathbf{p}, \mathbf{p}'} (-1)^{[\mathbf{p},\mathbf{r}']}r\left((\mathbf{r}+\mathbf{p})\mathbf{p},(\mathbf{r}'+\mathbf{p}')\mathbf{p}'\right)\sigma_{\mathbf{r}} \sigma_{\mathbf{r}'}\sigma_{\mathbf{p}} \sigma_{\mathbf{p}'} \otimes \sigma_{\mathbf{p}}\sigma_{\mathbf{p}'} \right) \notag \\
        &=4^{-n} \bigotimes_{j=1}^n \left(\sum_{\mathbf{p}, \mathbf{p}'} (-1)^{[\mathbf{p},\mathbf{r}']+[\mathbf{p},\mathbf{p}']}r\left((\mathbf{r}+\mathbf{p})\mathbf{p},(\mathbf{r}'+\mathbf{p}')\mathbf{p}'\right)\sigma_{\mathbf{r}} \sigma_{\mathbf{r}'}\sigma_{\mathbf{p}+\mathbf{p}'}\otimes \sigma_{\mathbf{p}+\mathbf{p}'} \right) \notag \\
        &= \bigotimes_{j=1}^n \left( \left\{\begin{aligned}
            &\frac{5}{24}\sum_{\mathbf{p}, \mathbf{p}'}(-1)^{[\mathbf{p},\mathbf{r}']+[\mathbf{p},\mathbf{p}']} \left[4\delta_{\mathbf{r},\mathbf{r}' }\delta_{\mathbf{p},\mathbf{p}' }+ 1 + (-1)^{[\mathbf{r}+\mathbf{p},\mathbf{r}'+\mathbf{p}']+[\mathbf{p},\mathbf{p}']}\right]\sigma_{\mathbf{r}} \sigma_{\mathbf{r}'}\sigma_{\mathbf{p}+\mathbf{p}'} \otimes \sigma_{\mathbf{p}+\mathbf{p}'} \text{~if~} \mathbf{r},\mathbf{r}' \neq 0\\
            &4^{-1}\sum_{\mathbf{p}, \mathbf{p}'}(-1)^{[\mathbf{p},\mathbf{p}']} r((\mathbf{r}+\mathbf{p})\mathbf{p},\mathbf{p}'\mathbf{p}')\sigma_{\mathbf{r}} \sigma_{\mathbf{p}+\mathbf{p}'} \otimes \sigma_{\mathbf{p}+\mathbf{p}'}\text{~if~} \mathbf{r}' =0\neq \mathbf{r}\\
            &4^{-1}\sum_{\mathbf{p}, \mathbf{p}'}(-1)^{[\mathbf{p},\mathbf{r}']+[\mathbf{p},\mathbf{p}']} r(\mathbf{p}\mathbf{p},(\mathbf{r}' + \mathbf{p}')\mathbf{p}') \sigma_{\mathbf{r}'}\sigma_{\mathbf{p}+\mathbf{p}'}\otimes \sigma_{\mathbf{p}+\mathbf{p}'}\text{~if~} \mathbf{r} =0\neq \mathbf{r}'\\
            &4^{-1}\sum_{\mathbf{p}, \mathbf{p}'}(-1)^{[\mathbf{p}', \mathbf{p}]} r(\mathbf{p}\mathbf{p},\mathbf{p}'\mathbf{p}')\sigma_{\mathbf{p} + \mathbf{p}'} \otimes \sigma_{\mathbf{p}+\mathbf{p}'}\text{~if~} \mathbf{r} = \mathbf{r}'  = 0
        \end{aligned}\right. \right) \notag \\
        &:= \bigotimes_{j=1}^n v(\mathbf{r}_j, \mathbf{r}_j')
        \label{eq:mid_V_Orho2}
        % = \frac{5}{3} \sigma_{\mathbf{r}_j}^{\otimes 2}
    \end{align}
    Here in the third equality, we use the fact that if $\mathbf{r},\mathbf{r}' \neq  \mathbf{0}$, we cannot have $\mathbf{r} + \mathbf{p} = \mathbf{p} = \mathbf 0$ or $\mathbf{r}' + \mathbf{p}' = \mathbf{p}' = \mathbf 0$. Thus, we plug in the second line in Eq.\eqref{eq:lc_Pauli_rescale}.

    In the case of $\mathbf{r}_j,\mathbf{r}'_j \neq 0$, we simplified term by term as below
    \begin{align}
        \text{the 1st term} &= \sum_{\mathbf{p}} (-1)^{[\mathbf{p},\mathbf{r}']}\cdot \frac{5}{6} \delta_{\mathbf{r},\mathbf{r}'}\sigma_{\mathbf{r}}\sigma_{\mathbf{r}'} \otimes \mathbb{I} =0, \\
        \text{the 2nd term} &= \sum_{\mathbf{p}, \mathbf{p}'}(-1)^{[\mathbf{p},\mathbf{r}']+[\mathbf{p},\mathbf{p}']} \frac{5}{24} \sigma_{\mathbf{r}} \sigma_{ \mathbf{r}'}\sigma_{\mathbf{p} + \mathbf{p}'} \otimes \sigma_{\mathbf{p}+\mathbf{p}'}= \sum_{\mathbf{p}, \mathbf{u}}(-1)^{[\mathbf{p}, \mathbf{r}']+[\mathbf{p},\mathbf{u}]} \frac{5}{24} \sigma_{\mathbf{r}} \sigma_{ \mathbf{r}'}\sigma_{\mathbf{u}} \otimes \sigma_{\mathbf{u}}\\ 
        &= \sum_{\mathbf{u}}\delta_{\mathbf{r}',\mathbf{u}} \frac{5}{6} \sigma_{\mathbf{r}} \sigma_{ \mathbf{r}'}\sigma_{\mathbf{u}} \otimes \sigma_{\mathbf{u}} = \frac{5}{6} \sigma_{\mathbf{r}} \otimes \sigma_{\mathbf{r}'},\\
        %%%%%%%%%%%%%%%%
        \text{the 3rd term} &= \sum_{\mathbf{p}, \mathbf{u}}(-1)^{[\mathbf{p}, \mathbf{r}']+[\mathbf{p},\mathbf{u}]} \frac{5}{24}\left[(-1)^{[\mathbf{r}+\mathbf{p},\mathbf{r}'+\mathbf{u}-\mathbf{p}]+[\mathbf{p},\mathbf{u}]}\right] \sigma_{\mathbf{r}}\sigma_{ \mathbf{r}'}\sigma_{\mathbf{u}} \otimes \sigma_{\mathbf{u}}\\
        &=\sum_{\mathbf{p}, \mathbf{u}}\frac{5}{24}\left[(-1)^{[\mathbf{r},\mathbf{r}']+[\mathbf{r},\mathbf{u}]+[\mathbf{r},\mathbf{p}]+[\mathbf{p},\mathbf{u}]}\right]\sigma_{\mathbf{r}}\sigma_{ \mathbf{r}'}\sigma_{\mathbf{u}} \otimes \sigma_{\mathbf{u}}\\
        &=\sum_{\mathbf{u}}\frac{5}{6}\delta_{\mathbf{r},\mathbf{u}}\left[(-1)^{[\mathbf{r},\mathbf{r}']+[\mathbf{r},\mathbf{u}]}\right]\sigma_{\mathbf{r}}\sigma_{ \mathbf{r}'}\sigma_{\mathbf{u}} \otimes \sigma_{\mathbf{u}} =\frac{5}{6}\sigma_{\mathbf{r}'} \otimes \sigma_{\mathbf{r}} .
    \end{align}
    Combining all these gives the result of the first line of Eq.\eqref{eq:mid_V_Orho2} as
    \begin{align}\label{eq:V_eq1}
       v({\mathbf{r}_j}, {\mathbf{r}'}_j) =\frac{5}{6}\left(\sigma_{\mathbf{r}_j}\otimes \sigma_{\mathbf{r}'_j} +\sigma_{\mathbf{r}'_j}\otimes \sigma_{\mathbf{r}_j} \right)~ \text{if~} \mathbf{r}_j, \mathbf{r}_j'  \neq \mathbf{0}
    \end{align} 
    
    In the case of $\mathbf{r}_j =\mathbf{r}'_j= 0$, we calculate by list 16 terms one-by-one as
    \begin{align}
        \text{second line} 
        &= 4^{-1}\left[(-1)^{0} \cdot 1\cdot  \sigma_{00}^{\otimes 2} +  (-1)^{0} \cdot 1\cdot  \sigma_{01}^{\otimes 2} +(-1)^{0} \cdot 1\cdot   \sigma_{10}^{\otimes 2} + (-1)^{0} \cdot 1\cdot   \sigma_{11}^{\otimes 2}\right. \notag\\
        &\hspace{0.8cm} + (-1)^{0} \cdot 1\cdot \sigma_{01}^{\otimes 2} + (-1)^{0} \cdot 5\cdot \sigma_{00}^{\otimes 2} + (-1)^{1} \cdot \frac{5}{3}\cdot\sigma_{11}^{\otimes 2} +  (-1)^{1} \cdot \frac{5}{3}\cdot \sigma_{10}^{\otimes 2} \notag\\
        &\hspace{0.8cm}+ (-1)^{0} \cdot 1\cdot \sigma_{10}^{\otimes 2} +  (-1)^{1} \cdot \frac{5}{3}\cdot \sigma_{11}^{\otimes 2} + (-1)^{0} \cdot 5\cdot \sigma_{00}^{\otimes 2} + (-1)^{1} \cdot \frac{5}{3}\cdot \sigma_{01}^{\otimes 2} \notag\\
        &\left.\hspace{0.8cm}+ (-1)^{0} \cdot 1\cdot \sigma_{11}^{\otimes 2} + (-1)^{1} \cdot \frac{5}{3}\cdot \sigma_{10}^{\otimes 2} + (-1)^{1} \cdot \frac{5}{3}\cdot \sigma_{01}^{\otimes 2} + (-1)^{0} \cdot 5\cdot \sigma_{00}^{\otimes 2} \right]\\
        &= 4 \sigma_{00}^{\otimes 2} - \frac{1}{3}(\sigma_{01}^{\otimes 2}+\sigma_{10}^{\otimes 2}+\sigma_{11}^{\otimes 2}) = \frac{11\mathbb{I}_2 + 4 \Pi_{\mathrm{asym}}}{3}.
    \end{align}
    Here in the last line we use the definition that $\Pi_{\mathrm{asym}} = (\mathbb{I}_2 - \text{SWAP})/2 = {(\mathbb{I}_2 -\sigma_{01}^{\otimes 2}-\sigma_{10}^{\otimes 2}-\sigma_{11}^{\otimes 2})}/{4}$.

    In the case of $\mathbf{r}'= 0 \neq \mathbf{r}$, if we simply adapt the result in Eq.\eqref{eq:V_eq1}, we will get the result $v({\mathbf{r}_j}, {\mathbf{r}_j'}) =({5}/{6})\left(\sigma_{\mathbf{r}}\otimes \mathbb{I} +\mathbb{I}\otimes \sigma_{\mathbf{r}} \right)$. However, in the case of $\mathbf{p}' = 0$ during the summation, we mistakenly counted the constant $r((\mathbf{r}+ \mathbf{p})\mathbf{p}, \mathbf{0}\mathbf{0}) = 1$ into $r((\mathbf{r}+ \mathbf{p})\mathbf{p}, \mathbf{0}\mathbf{0}) = (5/6)[4\delta_{\mathbf{r},\mathbf{0}}\delta_{\mathbf{p},\mathbf{0}} +2]$. To correct this term, we sum in the case $\mathbf{p}' = 0$ and get the result
    \begin{align}
        \Delta = -4^{-1}(5/6)\sum_{\mathbf{p}} [4\delta_{\mathbf{p},0}\delta_{\mathbf{r},0} + 2]\sigma_{\mathbf{r}}\sigma_{\mathbf{p}} \otimes \sigma_{\mathbf{p}} + 4^{-1} \sum_{\mathbf{p}} \sigma_{\mathbf{r}}\sigma_{\mathbf{p}} \otimes \sigma_{\mathbf{p}} = - \frac{1}{3}(\sigma_{\mathbf{r}} \otimes \mathbb{I})\text{SWAP}.
    \end{align} 
    This correction then gives the correct result
    \begin{align}
        v({\mathbf{r}}, {\mathbf{r}'}) &=\frac{5}{6}\left(\mathbb{I}\otimes \sigma_{\mathbf{r}} +  \sigma_{\mathbf{r}}\otimes\mathbb{I} \right) -  \frac{1}{3}(\sigma_{\mathbf{r}} \otimes \mathbb{I}) \cdot \text{SWAP} = \left(\frac{5}{6} \mathbb{I} - \frac{1}{3}\text{SWAP}\right)(\mathbb{I} \otimes \sigma_{\mathbf{r}}) + \frac{5}{6}\sigma_{\mathbf{r}}\otimes\mathbb{I} \notag\\
        &= \left(\frac{1}{2} \Pi_{\mathrm{sym}} + \frac{7}{6} \Pi_{\mathrm{asym}}\right)(\mathbb{I} \otimes \sigma_{\mathbf{r}}) + \frac{5}{6}\sigma_{\mathbf{r}}\otimes\mathbb{I} \text{~if~} \mathbf{r}' = \mathbf{0},  \mathbf{r}\neq \mathbf{0}.
    \end{align}
    Similarly, in the case $\mathbf{r}= 0 \neq \mathbf{r}'$, we calculate the correction in the case $\mathbf{p} = 0$ as
    \begin{align}
        \Delta &= -4^{-1}(5/6)\sum_{\mathbf{p}'}[4\delta_{\mathbf{p}',0}\delta_{\mathbf{r}',0} + 2]\sigma_{\mathbf{r}'}\sigma_{\mathbf{p}'} \otimes \sigma_{\mathbf{p}'} + 4^{-1} \sum_{\mathbf{p}'} \sigma_{\mathbf{r}'}\sigma_{\mathbf{p}'} \otimes \sigma_{\mathbf{p}'} = - \frac{1}{3}(\sigma_{\mathbf{r}'} \otimes \mathbb{I})\text{SWAP},\\
        v({\mathbf{r}}, {\mathbf{r}'}) &= \left(\frac{1}{2} \Pi_{\mathrm{sym}} + \frac{7}{6} \Pi_{\mathrm{asym}}\right)(\mathbb{I} \otimes \sigma_{\mathbf{r}'}) + \frac{5}{6}\sigma_{\mathbf{r}'}\otimes\mathbb{I} \text{~if~} \mathbf{r} = \mathbf{0},  \mathbf{r}'\neq \mathbf{0}.
    \end{align} 
    Now we conclude the results
    \begin{align}
        v(\mathbf{r}_j, \mathbf{r}_j') = 
        \left\{\begin{aligned}
            &\frac{5}{6}\left(\sigma_{\mathbf{r}_j}\otimes \sigma_{\mathbf{r}'_j} +\sigma_{\mathbf{r}'_j}\otimes \sigma_{\mathbf{r}_j} \right)&\text{~if~} \mathbf{r}_j, \mathbf{r}_j'  \neq \mathbf{0},\\
            &\frac{11\mathbb{I}_2 + 4 \Pi_{\mathrm{asym}}}{3} &\text{~if~} \mathbf{r}_j=\mathbf{r}_j'= \mathbf{0}, \\
            &\left(\frac{1}{2} \Pi_{\mathrm{sym}} + \frac{7}{6} \Pi_{\mathrm{asym}}\right)(\mathbb{I} \otimes \sigma_{\mathbf{r}_j})+ \frac{5}{6}\sigma_{\mathbf{r}}\otimes\mathbb{I} &\text{~if~} \mathbf{r}_j' = \mathbf{0},  \mathbf{r}_j\neq \mathbf{0},\\
            &\left(\frac{1}{2} \Pi_{\mathrm{sym}} + \frac{7}{6} \Pi_{\mathrm{asym}}\right)(\mathbb{I} \otimes \sigma_{\mathbf{r}'_j}) + \frac{5}{6}\sigma_{\mathbf{r}'}\otimes\mathbb{I}&\text{~if~} \mathbf{r}_j = \mathbf{0},  \mathbf{r}'_j\neq \mathbf{0}.
        \end{aligned}\right.
    \end{align}
    This completes the derivation of the full operator identity~\eqref{eq:v_full_operator}. For the special case $\mathbf{r} = \mathbf{r}'$, the non-identity sites give $v(\mathbf{r}_j,\mathbf{r}_j) = (5/3)\sigma_{\mathbf{r}_j}^{\otimes 2}$ while the identity sites give $v(\mathbf{0},\mathbf{0}) = (11\mathbb{I}_4 + 4\Pi_{\mathrm{asym}})/3$.

    \emph{State dependence of the second moment.} While $\rho^{\otimes 2}$ is globally symmetric under copy exchange, it is \emph{not} contained in $\bigotimes_j \mathrm{Sym}^2_j$ for entangled $\rho$. At each identity site, $v(\mathbf{0},\mathbf{0}) = (13\mathbb{I}_4 - 2\,\mathrm{SWAP})/3$ has eigenvalue $11/3$ on $\mathrm{Sym}^2$ and $5$ on the antisymmetric subspace. For product states $\rho = \bigotimes_l \rho_l$, the per-site state $\rho_l^{\otimes 2} \in \mathrm{Sym}^2(\mathbb{C}^2)$ and the $\Pi_{\mathrm{asym}}$ terms vanish, giving $\mathrm{tr}[\mathbf{V}(O_\mathbf{r},O_\mathbf{r})\rho^{\otimes 2}] = (11/3)^n(5/11)^{\mathrm{wt}(\sigma_\mathbf{r})}$. For entangled states the second moment can be larger; the operator norm bound gives $\|\mathbf{V}(O_\mathbf{0},O_\mathbf{0})\| = 5^n$.
\end{proof}
%%%%%%%%%%%%%%%%%%%%%%%%%%%%%%%%%%%%%%%%%%%%%%%

We consider the estimation of the simplest operator $O^{(2)} = \text{SWAP}^{\otimes n} = 2^{-n}(\mathbb{I}\mathbb{I}+XX+YY+ZZ)^{\otimes n}$.
By Lemma~\ref{lemma:covar_Pauli}, the second moment operator for purity estimation is $\mathbf{V}(O_\mathbf{0},O_\mathbf{0}) = \left({(11\mathbb{I}_4 + 4\Pi_{\mathrm{asym}})}/{3}\right)^{\otimes n}$. For product states the second moment is $(11/3)^n$; for entangled states it can be larger (operator norm $5^n$). In all cases the cost is exponential.
For the single copy local Clifford shadow, the variance is calculated through:
\begin{align}
    &\hspace{-1.5cm}\mathbf{V}:= \sum_{W,W' \in \mathcal{P}_2^{\otimes n}} \alpha_{W}\alpha_{W'} R(W,W') W W'= 4^{-n}\bigotimes_{i =1}^n \left(\sum_{W_i,W_i' \in \{\mathbb{I}\mathbb{I},XX,YY,ZZ\}} R(W_i,W'_i) W_i W_i'\right).\notag\\
   = 4^{-n} &\left( \mathbb{I}\mathbb{I}+XX+YY+ZZ+ XX+9\mathbb{I}\mathbb{I} + YY + 9\mathbb{I}\mathbb{I} +ZZ+ 9\mathbb{I}\mathbb{I}\right)^{\otimes n}\notag\\
  & \hspace{-1.15cm}= \left(7 \mathbb{I}\mathbb{I} + \frac{1}{2}(XX+YY+ZZ)\right)^{\otimes n}= \left(\frac{13\mathbb{I}\mathbb{I}+2\text{SWAP}}{2}\right)^{\otimes n} = \left(\frac{11\mathbb{I}\mathbb{I}+ 4\Pi_{\mathrm{sym}}}{2}\right)^{\otimes n}. \label{eq:prity_1rep_shadow_variance}
\end{align}
For the single-copy shadow, this gives a worst-case second moment of $7.5^n$ (the operator norm, attained on product states since the symmetric eigenvalue $15/2$ exceeds the antisymmetric eigenvalue $11/2$). The 2-replica local shadow improves the product-state second moment to $(11/3)^n \approx 3.67^n$, but the antisymmetric eigenvalue $5 > 11/3$ means entangled states can have larger variance (operator norm $5^n$). In all cases, the cost remains exponential; using higher replica numbers $k > 2$ under local measurement constraints only yields marginal improvements in the constant factor, but does not alter the complexity class.


\subsection{A 4-replica non-local strategy for the oblivious case}\label{app:UB_nonlin_P_shad}
We have proved the inefficiency of single-copy strategy while an efficient two-copy Bell strategy exists. 
In analogy, we try to break the lower bound in the last section by extending the single-copy shadow tomography method into the multi-copy case. 
\begin{theorem}[A $2k$-copy shadow protocol for $|\mathrm{tr}(P\rho^k)|$]
    Let $\mathcal{P}_n$ be the $n$-qubit Pauli group. Let $\mathbb{P}^{d}_{\pi}$ be the permutation operator on $2k$ replica of $n$-qubits ($d = 2^n$) and $\pi =(\mathrm{odds})(\mathrm{even})$ is 2-cycle permutation on odd labels $(1,3,5,\cdots,2k-1)$ and even labels $(2,4,\cdots,2k)$ respectively.
    For any $P \in \mathcal{P}_n$, define the operator: $ A_P =\mathrm{sym}_{\pi}\left[P^{\otimes 2} \otimes (I^{\otimes 2})^{\otimes (k-1)} \right] \cdot \mathbb{P}^{d}_{\pi}$.
    The family of observables $\mathcal{A} = \{A_P\}_{P \in \mathcal{P}_n}$ is a commuting family, i.e., $[A_P, A_Q] = 0$ for all $P, Q \in \mathcal{P}_n$.
    Furthermore, its expectation value on $\rho^{\otimes 2k}$ is
    $\mathrm{tr}(A_P \rho^{\otimes 2k}) = [\mathrm{tr}(P\rho^k)]^2 $.
    Therefore, a single projective measurement in the common eigenbasis of $\mathcal{A}$ allows for the simultaneous estimation of $|\mathrm{tr}(P\rho^k)|$ for all $P \in \mathcal{P}_n$.
\end{theorem}
    \begin{proof}
    We first prove the commutativity of the family $\mathcal{A}$.
    Let $S_P = \mathrm{sym}_{\pi}\left[P^{\otimes 2} \otimes (I^{\otimes 2})^{\otimes (k-1)} \right]$ and abbreviate $\pi = \mathbb{P}^d_{\pi}$. The operators are $A_P = S_P \pi$ and $A_Q = S_Q \pi$. First, $\pi$ commutes with $S_P, \forall P \in \mathcal{P}$ since $S_P$ is symmetrized under repeat operation of $\pi$.
    Now, the commutator $[A_P, A_Q]$ became:
    \begin{align}
        [A_P, A_Q] &= (S_P \pi) (S_Q \pi) - (S_Q \pi) (S_P \pi) \\
        &= S_P (S_Q \pi) \pi - S_Q (S_P \pi) \pi \quad \text{(since $\pi$ commutes with $S_P$ and $S_Q$)} \\
        &= [S_P, S_Q]\pi^2.
    \end{align}
    Finally, we show that $[S_P, S_Q] = 0$. This is because each term within $S_P = \mathrm{sym}_{\pi}\left[P^{\otimes 2} \otimes (I^{\otimes 2})^{\otimes (k-1)} \right]$ and $S_Q = \mathrm{sym}_{\pi}\left[Q^{\otimes 2} \otimes (I^{\otimes 2})^{\otimes (k-1)} \right]$
    either have commutators of operators on disjoint subsystems or have two-replica terms with $[P^{\otimes 2}, Q^{\otimes 2}] = [P \otimes P, Q \otimes Q] = 0$.
    Thus we have $[S_P, S_Q] = 0$ for all Pauli operators $P, Q$ .
    Thus, $[S_P, S_Q] = \mathbf{0}$, which implies $[A_P, A_Q] = \mathbf{0}$.
    We now calculate 
    \begin{align}
        \mathrm{tr}(A_P \rho^{\otimes 2k}) &= \mathrm{tr}[\mathrm{sym}_{\pi}(P^{\otimes 2} \otimes \mathbb{I}^{\otimes 2k-2}) \pi \rho^{\otimes 2k}] = \mathrm{tr}[(P^{\otimes 2} \otimes \mathbb{I}^{\otimes 2k-2}) \pi \rho^{\otimes 2k}]\\
        &=\mathrm{tr}[(P\rho \otimes \rho^{\otimes k-1}) \mathbb{P}^{d}_{(\mathrm{odds})} \otimes (P\rho \otimes \rho^{\otimes k-1}) \mathbb{P}^{d}_{(\mathrm{even})}]\\
        &=[\mathrm{tr}(P\rho^k)]^2
    \end{align}
    Thus, measuring the commuting family $\{A_P\}$ simultaneously yields estimates for $|\mathrm{tr}(P\rho^k)|$ for all $P \in \mathcal{P}_n$.
    \end{proof}
\begin{lemma}
    \label{lemma:signal_estim}
    If we can estimate $|\mathrm{tr}(P \rho^k)|, \forall P \in \mathcal{O} \subset \mathcal{P}_n$ within constant additive error $\varepsilon$, then using a $\mathcal{O}(\varepsilon^{-4}\log(|\mathcal{O}|))$ times $k$-replica strategy with $k \cdot n$ ancilla qubits, we can construct an estimator $E_P$ such that $|E_P - \mathrm{tr}(P \rho^{k})| \leq 3\varepsilon$. 
\end{lemma}
\begin{proof}[Algorithmic proof of Lemma~\ref{lemma:signal_estim}]
    \color{red}Let $\hat{f}_P$ denote the estimate of $|\mathrm{tr}(P \rho^k)|$ from the hypothesis, satisfying $|\hat{f}_P - |\mathrm{tr}(P \rho^k)|| \leq \varepsilon$, and write $f_P = |\mathrm{tr}(P \rho^k)|$ for the true value.\color{black} We then construct a quantum state $\sigma$ such that $\left||\mathrm{tr}(P \sigma^k)|- \hat{f}_P\right| \leq \varepsilon$, which implies $||\mathrm{tr}(P\sigma^k)| - f_P| \leq 2\varepsilon$ by the triangle inequality. Such a state $\sigma$ always exists and can be constructed given the estimates $\hat{f}_P$, but may not be found computationally efficiently.
    Then we estimate $\mathrm{tr}(\rho^k P)\cdot \mathrm{tr}(\sigma^k P)$ by measuring the commuting family $\{A_P\}_{P\in \mathcal{O}}$ on state $\rho^{\otimes k} \otimes \sigma^{\otimes k}$ for $\mathcal{O}(\varepsilon^{-4}\log(|\mathcal{O}|))$ times. This gives an estimator $g_P$ such that $|g_P - \mathrm{tr}(\rho^k P)\cdot \mathrm{tr}(\sigma^k P)|\leq \varepsilon^2$.
    \color{red}We set $E_P = 0$ if $\hat{f}_P \leq 2\varepsilon$, and $E_P = g_P/\mathrm{tr}(P \sigma^k)$ otherwise.
    We verify the error bound in two cases:
    \begin{align}
        \text{If~} \hat{f}_P \leq 2\varepsilon:\quad  &f_P \leq \hat{f}_P + \varepsilon \leq 3\varepsilon,\text{~so~} |E_P - \mathrm{tr}(P \rho^k)| = |\mathrm{tr}(P \rho^k)| = f_P \leq 3\varepsilon. \\
        \text{If~} \hat{f}_P > 2\varepsilon:\quad &|\mathrm{tr}(P \sigma^k)| \geq \hat{f}_P - \varepsilon > \varepsilon > 0,\text{~so~} |E_P - \mathrm{tr}(P \rho^k)| \notag\\
            &= \frac{|g_P - \mathrm{tr}(P \sigma^k) \cdot \mathrm{tr}(P \rho^k)|}{|\mathrm{tr}(P \sigma^k)|}
            \leq  \frac{\varepsilon^2}{\hat{f}_P - \varepsilon}  \leq \frac{\varepsilon^2}{\varepsilon} = \varepsilon \leq 3\varepsilon.
    \end{align}\color{black}
\end{proof}
\paragraph{Examples.}
 In the case $k = 2$, the commutants we choose writes $\{A_{\mathbf{r}} := {2}^{-1}{(\sigma_{\mathbf{r}}\sigma_{\mathbf{r}}\mathbb{I}_n\mathbb{I}_n + \mathbb{I}_n\mathbb{I}_n\sigma_{\mathbf{r}}\sigma_{\mathbf{r}})} \mathrm{S}_{13}\mathrm{S}_{24}\}_{\mathbf{r} \in \mathbb{F}_2^{2n}}$. Notice that $\sigma_{\mathbf{r}} \otimes \sigma_{\mathbf{r}} \ket{\Phi_{\mathbf{a}}} = (-1)^{[\mathbf{r},\mathbf{a}]}\ket{\Phi_{\mathbf{a}}}$,
the shared eigenbasis of $\mathcal{A}$ have the form $\ket{\Psi^{s}_{(\mathbf{a}, \mathbf{b})}}$ below with eigenvalues $\lambda_{(s,\mathbf{a},\mathbf{b})} :=s \cdot 2^{-1}[(-1)^{[\mathbf{r}, \mathbf{a}]} + (-1)^{[\mathbf{r}, \mathbf{b}]}] $ and $s = \pm 1$.
\begin{align}
   \ket{\Psi^{s}_{(\mathbf{a}, \mathbf{b})}}&= 
   \begin{cases}   
    {2^{-1/2}}\left[{\ket{\Phi_{\mathbf{a}_x \mathbf{a}_z}} \otimes \ket{\Phi_{\mathbf{b}_x\mathbf{b}_z} }+s \ket{\Phi_{\mathbf{b}_x \mathbf{b}_z}} \otimes \ket{\Phi_{\mathbf{a}_x\mathbf{a}_z}}}\right] & \text{if~} \mathbf{a} < \mathbf{b}, \\ 
    \ket{\Phi_{\mathbf{a}_x \mathbf{a}_z}}\ket{\Phi_{\mathbf{a}_x \mathbf{a}_z}}& \text{if~} \mathbf{a} = \mathbf{b}.
\end{cases}
\end{align}
After implementing the measurement $\left\{\ket{\Psi^{s}_{(\mathbf{a}, \mathbf{b})}} \bra{\Psi^{s}_{(\mathbf{a}, \mathbf{b})}}\right\}_{s \in \{+,-\}, \mathbf{a},\mathbf{b} \in \mathbb{F}_2^{2n}}$ for $T$ times and recording the results $\{(s_t,\mathbf{a}_t, \mathbf{b}_t)\}_{t = 1}^{T}$, the average of $A_{\mathbf{r}} = \sum \lambda_{(\pm,\mathbf{a},\mathbf{b})} \ket{\Psi^{\pm}_{(\mathbf{a}, \mathbf{b})}} \bra{\Psi^{\pm}_{(\mathbf{a}, \mathbf{b})}}$ can then be estimated by calculating the estimator $T^{-1}\sum_{t = 1}^{T} \lambda_{(s_t, \mathbf{a}_t,\mathbf{b}_t)} $.
To realize such a measurement, we refer to construction in~\cite{liu2024auxiliary} and noted that 
\begin{align}
    \ket{\Psi^{s}_{\mathbf{a}, \mathbf{b}}}&= \mathrm{CNOT}_{34}\mathrm{CNOT}_{12}\mathrm{H}_3\mathrm{H}_1 \cdot
\begin{cases}
    2^{-1/2} \left[{\ket{\mathbf{a}_x \mathbf{a}_z\mathbf{b}_x\mathbf{b}_z} +s \ket{\mathbf{b}_x \mathbf{b}_z\mathbf{a}_x\mathbf{a}_z}}\right] &\text{if~} \mathbf{a} < \mathbf{b}\\
    \ket{\mathbf{a}_x \mathbf{a}_z\mathbf{b}_x\mathbf{b}_z}&\text{if~} \mathbf{a} = \mathbf{b}
\end{cases},\\
&= \mathrm{CNOT}_{34}\mathrm{CNOT}_{12}\mathrm{H}_3\mathrm{H}_1 \mathrm{CNOT}_{13}\mathrm{CNOT}_{24}\cdot
\begin{cases}
    2^{-1/2} \left[{\ket{\mathbf{a}_x \mathbf{a}_z} +s \ket{\mathbf{b}_x \mathbf{b}_z}} \right]\ket{\mathbf{a} + \mathbf{b}} &\text{if~} \mathbf{a} < \mathbf{b}\\
    \ket{\mathbf{a}_x \mathbf{a}_z} \ket{\mathbf{0} \mathbf{0}}&\text{if~} \mathbf{a} = \mathbf{b}
\end{cases},\\
&= \mathrm{CNOT}_{34}\mathrm{CNOT}_{12}\mathrm{H}_3\mathrm{H}_1 \mathrm{CNOT}_{13}\mathrm{CNOT}_{24}\cdot
\begin{cases}
    2^{-1/2} \left[{\ket{\mathbf{a}^{I}} +s \ket{\bar{\mathbf{a}}^{I}}}\right] \ket{\mathbf{a}^{E}}\ket{\mathbf{a} + \mathbf{b}} &\text{if~} \mathbf{a} < \mathbf{b}\\
    \ket{\mathbf{a}_x \mathbf{a}_z} \ket{\mathbf{0} \mathbf{0}}&\text{if~} \mathbf{a} = \mathbf{b}
\end{cases}.
\end{align}
This disentangled the qubit on the third and the forth copy. Thus, we can locally implement the gate sequence $ \mathrm{CNOT}_{13}\mathrm{CNOT}_{24}\mathrm{H}_3\mathrm{H}_1\mathrm{CNOT}_{12}\mathrm{CNOT}_{34}$ and then measure the $3$-rd and $4$-th copies on the computational basis $\ket{\mathbf{s}_{3}\mathbf{s}_4}$. This gives the difference of vector $\mathbf{a}$ and $\mathbf{b}$ through $\mathbf{s}_{3} = \mathbf{a}_x + \mathbf{b}_x,\mathbf{s}_{4} = \mathbf{a}_z + \mathbf{b}_z$. 
This information then indicates a separtiton of the $1$-st and $2$-nd copies of qubits $\{1,\cdots,2n\} = E \cup I$, with $E := \{i \in [2n]| a_i = b_i\} $ and $I:= \{i \in [2n]| \bar{a}_i = b_i\} := \{i_1<i_2<\cdots<i_{|I|}\}$. 
We then implement the computational basis measurement $\ket{\mathbf{s}^{E}}$ on the qubits in set $E$ and a GHZ state measurement $s^{I}$ formulized by circuit:
\begin{align}
    \frac{1}{\sqrt{2}} \left[{\ket{\mathbf{a}^{I}} + s\ket{\bar{\mathbf{a}}^{I}}} \right] &= 
    \begin{cases}
        \frac{1}{\sqrt{2}} \left(\bigotimes_{i \in I} X_i + Z_{i_1}\right) \ket{\bar{\mathbf{a}}^{I}}  &\text{if~} s = -1,\\
        \frac{1}{\sqrt{2}} \left(\bigotimes_{i \in I} X_i + Z_{i_1}\right) \ket{{\mathbf{a}}^{I}}  &\text{if~} s = +1,
    \end{cases}\\
    &=\left(\prod_{i = i_1}^{i_{|I|-1}}\mathrm{CNOT}_{i,i+1}\right)^{\dagger} H_{i_1} \left(\prod_{i = i_1}^{i_{|I|-1}}\mathrm{CNOT}_{i,i+1}\right) \begin{cases}\ket{\bar{\mathbf{a}}^{I}} &\text{if~}s = -1,\\\ket{{\mathbf{a}}^{I}} &\text{if~}s = 1.\end{cases}
\end{align}
Here we use the fact that $\mathbf{a}< \mathbf{b}$, thus $a_{i_1} = 0$ is satisfied.
Note that the last $\text{CNOT}$ gates on computational basis can be simulated by classical post-processing. This 
This then gives the $\mathbf{a}$ vector through $\mathbf{a}^{E} = {\mathbf{s}^{E}}$ and $s = +1, \mathbf{a}^{I} = {s}^{I}$ if $a_{i_1} = 0$, $s = -1, \mathbf{a}^{I} = \bar{s}^{I}$ if $a_{i_1} = 1$.
\begin{figure}[htbp!]
    \centering
    \includegraphics[width=0.5\linewidth]{Figures/4rep_strategy.pdf}
    \caption{The demonstration of 4-replica nonlocal strategy that realized the non-linear Pauli shadows for $|\mathrm{tr}(P \rho^2)|$, for all Pauli operators $P$.}
    \label{fig:4rep_strategy}
\end{figure}

\section{Lower bounds for nonlinear Pauli shadow}\label{app:LB_nonlin_P_shad}

\subsubsection{Introduction to Haar-assembled ensembles}
The lower bound of many problem are find by first reduce this problem to distinguishing one specific state with a state ensemble.
One useful state ensembles are constructed from reducing Haar random state.
\begin{definition}[\textbf{RH} ensembles~\cite{gong2024sample}]
    Given a $D = 2^{n + \log(1/\varepsilon)} = d\cdot \varepsilon^{-1}$ dimensional Haar random state $\ket{H} \sim \mathrm{Haar}(D)$. The $2^n$-dimensional reduced Haar-induced (\textbf{RH}) state is defined as $\mathcal{E}_{\varepsilon}:=\{\text{tr}_{>n} \ket{H}\!\bra{H}\}_{\ket{H} \sim \mathrm{Haar}(D)}$.
\end{definition}
Depending on the scale of reduced parts, the $RH$-ensemble has tunable purity.
We then show in Lemma~\ref{lemma:purity_of_RH} that the random variable $\mathrm{tr}(\rho^2), \rho \sim \mathcal{E}^{\mathrm{RH}}_{\varepsilon}$ is well concentrated and separate from maximum mixed state purity $1/d$.
This statement is useful to proving the reduction of $\mathrm{tr}(\rho^2)$ problem. 
\begin{lemma}[Purity concentration of $\mathcal{E}_{\varepsilon}$]
    \label{lemma:purity_of_RH}
    The $c$-copy average of $\rho \in \mathcal{E}_{\varepsilon}$ writes:
    \begin{align}
        \mathbb{E}_{\rho \sim \mathcal{E}_{\varepsilon}} \rho^{\otimes c} = \frac{1}{D^{\uparrow c}}\sum_{\pi \in S_c} \varepsilon^{-c(\pi)}\mathbb{P}_{\pi}^{d}.
    \end{align}
    Furthermore, its purity average and variance writes:
    \begin{align}
        \mathbb{E}_{\rho \sim \mathcal{E}_{\varepsilon}} \mathrm{tr}(\rho^2) = \frac{2^n\varepsilon + 1}{2^n + \varepsilon},\quad \mathrm{Var}_{\rho \sim \mathcal{E}_{\varepsilon}} [\mathrm{tr}(\rho^2)] = \frac{2(2^{2n} - 1)(1/\varepsilon^2-1)}{(2^n/\varepsilon + 1)^2 (2^n/\varepsilon+2)(2^n/\varepsilon+3)}.
    \end{align}
\end{lemma}
\begin{proof}
    \textcolor{red}{[Proof to be completed]}
    \begin{align}
        \mathbb{E}_{\rho \sim \mathcal{E}_{\varepsilon}}\rho^{\otimes c}&=\mathbb{E}_{H \sim \mathrm{Haar}(D)}(\text{tr}_{>n} \ket{H}\!\bra{H})^{\otimes c} = \frac{1}{D^{\uparrow c}}\mathrm{tr}_{>n}(\sum_{\pi \in S_c} \mathbb{P}_{\pi}^{d} \otimes \mathbb{P}_{\pi}^{\varepsilon^{-1}})=\frac{1}{D^{\uparrow c}}\sum_{\pi \in S_c} \varepsilon^{-c(\pi)}\mathbb{P}_{\pi}^{d}.
    \end{align}
    Then the average of purity can be directly calculate as
    \begin{align}
    \mathbb{E}_{\rho \sim \mathcal{E}_{\varepsilon}} \mathrm{tr}(\rho^2)
    &=\mathrm{tr}\left(\frac{\varepsilon^{-2} \mathbb{I}_{n}^{\otimes 2}+\varepsilon^{-1} \mathrm{SWAP}^{\otimes n}}{D(D+1)} \cdot \mathrm{SWAP}^{\otimes n}\right) = \frac{\varepsilon^{-2} 2^n + \varepsilon^{-1} 2^{2n}}{2^n\varepsilon^{-1}(2^n \varepsilon^{-1}+1)}= \frac{2^n\varepsilon + 1}{2^n + \varepsilon},\\
    \mathbb{E}_{\rho \sim \mathcal{E}_{\varepsilon}} [\mathrm{tr}(\rho^2)]^2&=\mathrm{tr}\left(\frac{ \sum_{\pi \in S_4}{\varepsilon^{-c(\pi)}\mathbb{P}_{\pi}^{d}}}{D(D+1)(D+2)(D+3)}\mathbb{P}_{(12)(34)}^{d}\right)=\frac{\sum_{\pi \in S_4}d^{c(\pi \cdot (12)(34))}(1/\varepsilon)^{c(\pi)}}{D(D+1)(D+2)(D+3)}.
    \end{align}
    Here $c(\pi)$ count the number of cycles for permutation $\pi$. To calculate this value, we classify the $S_4$ group as shown in Table~\ref{tab:S4_terms}. 
    Through denoting $d_A = 2^n, d_B = 1 / \varepsilon$, $D = d_Ad_B$. We get the results
    \begin{align}
        \mathrm{Var}_{\rho \sim \mathcal{E}_{\varepsilon}} [\mathrm{tr}(\rho^2)] &= \frac{d_A^2 d_B^4+2 d_A^3 d_B^3+5 d_A d_B^3+10 d_A^2 d_B^2+5 d_A^3 d_B+d_A^4 d_B^2}{(d_Ad_B)(d_Ad_B+1)(d_Ad_B+2)(d_Ad_B+3)} - \frac{(d_A + d_B)^2}{(d_Ad_B+1)^2}\notag\\
        &= \frac{2d_B\left[1 + d_A^3 - 2d_A -d_B +d_Ad_B\right]}{(d_Ad_B + 1)^2 (d_Ad_B+2)(d_Ad_B+3)} = \frac{2/\varepsilon\left[1 + 2^{3n} - 2^{2n+1} - 1/\varepsilon +2^n/\varepsilon\right]}{(2^n/\varepsilon + 1)^2 (2^n/\varepsilon+2)(2^n/\varepsilon+3)}
    \end{align} 
    \begin{table}[htbp!]
        \centering
        \begin{tabular}{|l|c|c|l|c|l|}
        \hline
        \textbf{$\pi$} & \textbf{number} & \textbf{$c(\pi)$} & \textbf{$\pi \cdot (12)(34)$} & \textbf{$c(\pi \cdot (12)(34))$} & \textbf{summation} \\ \hline
        $e$ & 1 & 4 & (12)(34) & 2 & $d_A^2 d_B^4$ \\ \hline
        (12),(34), & $C_4^2$ & 3 & (34),(12) & 3 & $2 d_A^3 d_B^3$ \\ \cline{4-6}
        (13),(14),(23),(24)&  &  & (1432),(1342),(1243),(1342) & 1 & $4 d_A d_B^3$ \\ \hline
        $\substack{(123),(321),(124),(421),\\(134),(431),(234),(432)}$& $2C_4^3$ & 2 & $\substack{(243),(143),(213),(421),\\(134),(431),(234),(432)}$ & 2 & $8 d_A^2 d_B^2$ \\ \hline
        $\substack{(1234),(1243),(1324),\\(1342),(1423),(1432)}$ & (4-1)! & 1 & (13) & 3 & $6 d_A^3 d_B$ \\ \hline
        $\substack{(12)(34),}$ & 3 & 2 & $e$ & 4 & $d_A^4 d_B^2$ \\ \cline{4-6}
         $\substack{(13)(24),(14)(23)}$& & & $(14)(23)$ & 2 & $2 d_A^2 d_B^2$ \\ \hline
        \end{tabular}
        \caption{Classification of $S_4$ to calculate $\sum_{\pi \in S_4}(d_A)^{c(\pi \cdot (12)(34))}(d_B)^{c(\pi)}$ }
        \label{tab:S4_terms}
        \end{table}
\end{proof}


\subsubsection{Subspace \textbf{RH}-ensemble test and its hardness}
\begin{problem}\label{prob:RHtest}
    Given \textbf{RH} ensembles $\mathcal{E}_{1 - R\delta}, \mathcal{E}_{L\delta}$ (as in Sec.~\ref{app:LB_nonlin_P_shad}), with $R = (2^{n-1} + 1)/(2^{n-1} - 1 +\delta)$ and $L = (1 - 4^{-n+1} - 2^{-n+1} \delta)^{-1}$, try to design an algorithm that $\forall C \in \mathcal{C}$, discriminates two subspace-\textbf{RH} ensembles below:
    \begin{align}
        &\mathcal{E}_{(1,0),C} := \left\{\sigma_{(1,0),C}:= C \left[\frac{1}{2} \ket{0}\!\bra{0} \otimes \psi + \frac{1}{2} \ket{1}\!\bra{1} \otimes (\mathbb{I}/2^{n-1}) \right] C^{\dagger}\right\}_{\psi \sim \mathrm{Haar}_{n-1}},\\
        &\mathcal{E}_{(1 - R\delta,U\delta),C} := \left\{ \sigma_{(1 - R\delta,L\delta),C} :=C \left[\frac{1}{2} \ket{0}\!\bra{0} \otimes \rho_{1 - R\delta} + \frac{1}{2} \ket{1}\!\bra{1} \otimes \rho_{L\delta} \right] C^{\dagger}\right\}_{\rho_{\varepsilon} \sim \mathcal{E}_{\varepsilon}}.
    \end{align}
\end{problem}
This construction ensures that the average purity of $\mathcal{E}_{(1,0)}$ and $\mathcal{E}_{(1 - R\delta,L\delta)} $ equals each others, as manifest below:
\begin{align}
    &\mathbb{E}_{\rho \sim \mathcal{E}_{(1,0),C}} \mathrm{tr}\left(\rho^{2}\right) = 1/4 +  (1/4) \cdot (1/2^{n-1}) =  \frac{1}{4}\left(1+\frac{1}{2^{n-1}}\right), \\
    &\mathbb{E}_{\rho \sim \mathcal{E}_{(1 - R\delta,U\delta),C}} \mathrm{tr}\left(\rho^{2}\right) = \frac{1}{4}\frac{2^{n-1}(1 - R\delta) + 1}{2^{n-1} + 1 - R\delta} + \frac{1}{4}\frac{2^{n-1} L\delta + 1}{2^{n-1} + L\delta} = \frac{1}{4}(1 - \delta) + \frac{1}{4}(\frac{1}{2^{n-1}} + \delta) = \frac{1}{4}\left(1+\frac{1}{2^{n-1}}\right).
\end{align}

Given a $T$-round $2$-replica measurement strategy $M_s = \bigotimes_{t = 1}^{T} F_{s_t}^{u_t}$, where $\{F_{s}^{u}\}_{s}$ is a POVM on space $\mathcal{H}(d^2)$.
Now the \textbf{TVD} under this measurement takes the value:
\begin{align}
    \delta_{\mathrm{tvd}} = \left|\mathrm{tr}\left(M_s  \mathbb{E}_{C}C^{\otimes 2T}\mathbb{E}_{\psi}\left(\sigma_{(1,0),I}^{\otimes 2T} - \sigma_{(1 - R\delta,L\delta),I}^{\otimes 2T}\right)(C^\dagger)^{\otimes 2T}\right)\right|,
\end{align}
where the Haar average takes the value:
\begin{align}
    \mathbb{E}_{\psi}\sigma_{(1,0),I}^{\otimes 2T} &= \left[\frac{1}{2}\ket{0}\!\bra{0} \otimes \psi  + \frac{1}{2}\ket{1}\!\bra{1} \otimes 2^{-n+1}\mathbb{I}_{n-1}\right]^{\otimes 2T} = \sum_{\mathbf{a} \in \mathbb{F}_2^{2T}} {2^{-2T}\ket{\bar{\mathbf{a}}}\!\bra{\bar{\mathbf{a}}} }\otimes\frac{\sum_{\pi \in S_{2T}^{\mathbf{a}}} \mathbb{P}_{\pi}^{d}}{d^{\uparrow |\mathbf{a}|} d^{2T - |\mathbf{a}|}},\\
    \mathbb{E}_{\psi}\sigma_{(1 - R\delta,L\delta),I}^{\otimes 2T} &= \left[\frac{1}{2}\ket{0}\!\bra{0} \otimes \rho_{1 - R\delta} + \frac{1}{2}\ket{1}\!\bra{1} \otimes 2^{-n+1}\rho_{L\delta}\right]^{\otimes 2T} \\
    &= \sum_{\mathbf{a} \in \mathbb{F}_2^{2T}} {2^{-2T}\ket{\bar{\mathbf{a}}}\!\bra{\bar{\mathbf{a}}}} \otimes \frac{\sum_{\pi \in S_{2T}^{\mathbf{a}}, \sigma \in S_{2T}^{\bar{\mathbf{a}}}} \mathbb{P}_{\pi \sigma}^{d} (1 - R \delta)^{-c(\pi)} (L\delta)^{-c(\sigma)}}{[d/(1 - R\delta)]^{\uparrow |\mathbf{a}|} [d/L\delta]^{\uparrow |\bar{\mathbf{a}}|}},
\end{align}
where $d = 2^{n-1}$ and $S_{2T}^{\mathbf{a}}$ denoting the permutation contained in position $\{i|a_i = 1\}$.
Then define their difference $\Delta_{2T} :=\mathbb{E}_{\psi}\sigma_{(1,0),I}^{\otimes 2T}  -  \mathbb{E}_{\psi}\sigma_{(1 - R\delta,L\delta),I}^{\otimes 2T} $. Specifically, when $T = 1$, we conclude that:
\begin{align}
  \Delta_{2} =  4^{-1}  &{\ket{11}\!\bra{11}} 
     \otimes \left[\frac{\mathbb{I}^{\otimes 2}}{ d^{2}} - \frac{\mathbb{I}^{\otimes 2} (L\delta)^{-2} + \mathbb{P}^{d}_{(12)}(L\delta)^{-1}}{[d/L\delta][d/L\delta + 1] }\right] \\
     &+ 4^{-1}  {\ket{00}\!\bra{00}} 
     \otimes \left[\frac{\mathbb{I}^{\otimes 2}  + \mathbb{P}^{d}_{(12)}}{d[d + 1] } - \frac{\mathbb{I}^{\otimes 2} (1 -R\delta)^{-2} + \mathbb{P}^{d}_{(12)}(1 -R\delta)^{-1}}{[d/(1 -R\delta)][d/(1 -R\delta) + 1] }\right].
\end{align}
We now build uniform distribution on all $n$-qubit Clifford gates $\mathcal{D}$ and average $\Delta_{2}$ over $2$-fold Clifford. Since the average $\Phi_{\mathrm{Cl},2}[\Delta_2]$ is invariant under the two-fold Clifford $X_1 \otimes X_1$, we conclude that 
\begin{align}
    \Phi_{\mathrm{Cl},2}[\Delta_2] &= \Phi_{\mathrm{Cl},2}\left[4^{-1} {\ket{11}\!\bra{11}} 
     \otimes \left[-\frac{\mathbb{I}^{\otimes 2} + \mathbb{P}^{d}_{(12)}L\delta}{d[d + L\delta] } +  \frac{\mathbb{I}^{\otimes 2}}{ d^{2}} + \frac{\mathbb{I}^{\otimes 2}  + \mathbb{P}^{d}_{(12)}}{d[d + 1] } -\frac{\mathbb{I}^{\otimes 2} + \mathbb{P}^{d}_{(12)}(1 -R\delta)}{d[d+ 1 -R\delta] } \right]\right]\\
     &=\Phi_{\mathrm{Cl},2}\left[4^{-1} {\ket{11}\!\bra{11}} 
     \otimes \left[\left(-\frac{1}{d[d + L\delta] } +  \frac{1}{ d^{2}} + \frac{1}{d[d + 1] } -\frac{1}{d[d+ 1 -R\delta] } \right)\mathbb{I}^{\otimes 2} +\right.\right.\notag\\
    & \hspace{13em}\left.\left. \left(-\frac{L\delta}{d[d + L\delta] } + \frac{1}{d[d + 1] } -\frac{ 1-R\delta}{d[d+ 1 -R\delta] } \right)\mathbb{P}^{d}_{(12)}\right]\right]\\
    &=\Phi_{\mathrm{Cl},2}\left[4^{-1} {\ket{11}\!\bra{11}} 
     \otimes \left[\left(\frac{L\delta}{d^2(d + L \delta)} -\frac{R\delta}{d(d+1)(d+1 - R\delta)} \right)\mathbb{I}^{\otimes 2} +\right.\right.\notag\\
    & \hspace{13em}\left.\left. \left(-\frac{L\delta}{d[d + L\delta] } + \frac{R\delta d}{d(d+1)(d+1-R\delta)}\right)\mathbb{P}^{d}_{(12)}\right]\right]\\
    &=\Phi_{\mathrm{Cl},2}\left[4^{-1} {\ket{11}\!\bra{11}}
     \otimes \left[\left(\frac{\delta}{d(d^2 -1)} -\frac{\delta}{d(d+1)(d - 1)} \right)\mathbb{I}^{\otimes 2} +\right.\right.\notag\\
    & \hspace{13em}\left.\left. \left(-\frac{\delta}{d^2 - 1} + \frac{d\delta}{d(d+1)(d-1)}\right)\mathbb{P}^{d}_{(12)}\right]\right] = 0. \label{eq:t=1_Delta}
\end{align}
Here in the last line we use the identity $(d-1)^{-1}\delta = (d + 1 - R\delta)^{-1} R\delta$ and  $d(d^2 - 1)^{-1}\delta = L\delta (d+L\delta)^{-1}$.
Since Clifford gate forms a 2-design, we can directly write the average based on the value of average purity
\begin{align}
    &\rho_0 := \mathbb{E}_{\rho \sim \mathcal{E}_{(1,0),C\sim \mathcal{D}}} \rho^{\otimes 2}=\mathbb{E}_{\rho \sim \mathcal{E}_{(1 - R\delta,U\delta),C\sim \mathcal{D}}} \rho^{\otimes 2} = a\cdot \mathbb{I}_{2n}/2^{2n} + b\cdot \mathrm{SWAP}^{\otimes n}/2^{n}, \\
    & a + b = 1, \quad 2^{-n} a + 2^{n} b = \text{Avg. Purity} = 4^{-1}(1 + 2^{-n+1}).
\end{align}
Solve the equation for $a$ and $b$, we get the result:
\begin{align}
    \rho_0 = (1 - \mu)  \frac{\mathbb{I}_n^{\otimes 2}}{2^{2n}} + \mu \frac{\mathbb{I}_n^{\otimes 2} + \mathrm{SWAP}^{\otimes n}}{2^n(2^{n} + 1)}, \text{~where~} \mu = 4^{-1} \frac{2^n -2}{2^n - 1}.
\end{align}
This prevents the Lemma~\ref{lemma:hard_necess1} giving an efficient 2-replica strategy that discriminates $\mathcal{E}_{(1,0),C} $ and $\mathcal{E}_{(1 - R\delta,U\delta),C} $.
However, given a fixed $C \in \mathcal{C}$, the average of $\mathrm{tr}(P \rho^2)$ for $P = C^\dagger Z_1 C$ has a constant gap $\delta/2$:
\begin{align}
    &\mathbb{E}_{\rho \sim \mathcal{E}_{(1,0),C}} \mathrm{tr}\left(P\rho^{2}\right) = 1/4 -  (1/4) \cdot (1/2^{n-1}) = \frac{1}{4}\left(1- \frac{1}{2^{n-1}}\right) ,\\
    &\mathbb{E}_{\rho \sim \mathcal{E}_{(1 - R\delta,U\delta),C}} \mathrm{tr}\left(P\rho^{2}\right)  = \frac{1}{4}(1 - \delta) - \frac{1}{4}(\frac{1}{2^{n-1}} + \delta) = \frac{1}{4}\left(1- \frac{1}{2^{n-1}}\right) - \frac{\delta}{2}.
\end{align}
This means that if one can use strategy $\mathcal{A}$ to estimate $\mathrm{tr}(P \rho^2), \forall P \in \{C^\dagger Z_1 C\}_{C \in \mathcal{C}}$ within additive error $\delta/4$, he can use $\mathcal{A}$ to discriminate $\mathcal{E}_{(1,0),C} $ and $\mathcal{E}_{(1 - R\delta,U\delta),C} $, as rigorously proved below:
\begin{lemma}\label{lemma:hard_necess1}
    An algorithm that estimates $\mathrm{tr}(P \rho^2), \forall P \in \{C Z_1 C^\dagger\}_{C \in \mathcal{C}}$ can be used to distinguish $\mathcal{E}_{(1,0),C} $ from $\mathcal{E}_{(1 - R\delta,U\delta),C} $ for arbitrary distribution on $\mathcal{C}$.
\end{lemma}
\begin{proof}
    \textcolor{red}{[Proof to be completed]}
\end{proof}
Now we prove that two-replica strategy that discriminate $\sigma_{(1,0),C}^{\otimes 2T}$ from $\sigma_{(1 - R\delta,L\delta),C}^{\otimes 2T} $ require $\Omega(2^{n})$ rounds. We first calculate the $2T$-fold Clifford average of their difference. Since $X_1^{\otimes 2T} $ is a valid  $2T$-fold Clifford, $\Phi_{\mathrm{Cl},2T}[\Delta_{2T}]$ is invariant under the conjugates of $X_1^{\otimes 2T} $. 
% \begin{align}
%  &\leq \Phi_{\mathrm{Cl},2T}\left[\sum_{\mathbf{a} \in \mathbb{F}_2^{2T}} {2^{-2T}\ket{\bar{\mathbf{a}}}\!\bra{\bar{\mathbf{a}}}} \otimes \left( \frac{\sum_{\pi \in S_{2T}} \mathbb{P}_{\pi}^{d}}{d^{\uparrow 2T} } + \frac{\mathbb{I}(d)^{\otimes 2T}}{ d^{2T}} \right. \right.-\left.\left.\frac{\sum_{\pi \in S_{2T}} \mathbb{P}_{\pi}^{d} (1 - R \delta)^{-c(\pi)} }{[d/(1 - R\delta)]^{\uparrow 2T} }   - 
%    \frac{\sum_{\pi \in S_{2T}} \mathbb{P}_{\pi }^{d}  (L\delta)^{-c(\pi)}}{[d/L\delta]^{\uparrow 2T}}\right)\right]
%   \Delta_{2T} =   \sum_{\mathbf{a} \in \mathbb{F}_2^{2T}} {2^{-2T}\ket{\bar{\mathbf{a}}}\!\bra{\bar{\mathbf{a}}}} 
%     \otimes \frac{ \sum_{\pi \in S_{|\mathbf{a}|}}\mathbb{P}_{\pi}^{d}}{d^{\uparrow |\mathbf{a}|}} \otimes \left[\frac{ d^{\uparrow |\mathbf{a}|}(1 - R \delta)^{-c(\pi)}}{[d/L\delta]^{\uparrow |\bar{\mathbf{a}}|}[d/(1 - R\delta)]^{\uparrow |\mathbf{a}|} }{\sum_{ \sigma \in S_{|\bar{\mathbf{a}}|}}\mathbb{P}_{\sigma}^{d}   (L\delta)^{-c(\sigma)}} -  \frac{\mathbb{I}_{n-1}^{\otimes |\bar{\mathbf{a}}|}}{ d^{|\bar{\mathbf{a}}|}}\right].
% \end{align}
\begin{align}
    &\Phi_{\mathrm{Cl},2T}[\Delta_{2T}] = \Phi_{\mathrm{Cl},2T}\left[\sum_{\mathbf{a} \in \mathbb{F}_2^{2T}} {2^{-2T}\ket{\bar{\mathbf{a}}}\!\bra{\bar{\mathbf{a}}}} \otimes \left( \frac{\sum_{\pi \in S_{2T}^{\mathbf{a}}} \mathbb{P}_{\pi}^{d}}{d^{\uparrow |\mathbf{a}|} d^{2T - |\mathbf{a}|}} + \frac{\sum_{\pi \in S_{2T}^{\bar{\mathbf{a}}}} \mathbb{P}_{\pi}^{d}}{d^{\uparrow |\bar{\mathbf{a}}|} d^{2T - |\bar{\mathbf{a}}|}} \right. \right.\\
   &-\left.\left.\frac{\sum_{\pi \in S_{|\mathbf{a}|}} \mathbb{P}_{\pi }^{d} (1 - R \delta)^{-c(\pi)}}{[d/(1 - R\delta)]^{\uparrow |\mathbf{a}|} } \otimes \frac{ \sum_{\sigma \in S_{|\bar{\mathbf{a}}|}} \mathbb{P}_{ \sigma}^{d}(L\delta)^{-c(\sigma)}}{[d/L\delta]^{\uparrow |\bar{\mathbf{a}}|}} -\frac{\sum_{\pi \in S_{|\mathbf{a}|}} \mathbb{P}_{\pi }^{d} (L \delta)^{-c(\pi)}}{[d/(L\delta)]^{\uparrow |\mathbf{a}|} } \otimes \frac{ \sum_{\sigma \in S_{|\bar{\mathbf{a}}|}} \mathbb{P}_{ \sigma}^{d}(1 - R\delta)^{-c(\sigma)}}{[d/(1 - R\delta)]^{\uparrow |\bar{\mathbf{a}}|}} 
   \right)\right].\\
    &\leq \Phi_{\mathrm{Cl},2T}\left[\sum_{\mathbf{a} \in \mathbb{F}_2^{2T}} {2^{-2T}\ket{\bar{\mathbf{a}}}\!\bra{\bar{\mathbf{a}}}} \otimes \left( \frac{\sum_{\pi \in S_{2T}} \mathbb{P}_{\pi}^{d}}{d^{\uparrow 2T} } + \frac{\mathbb{I}(d)^{\otimes 2T}}{ d^{2T}} \right. \right.-\left.\left.\frac{\sum_{\pi \in S_{2T}} \mathbb{P}_{\pi}^{d} (1 - R \delta)^{-c(\pi)} }{[d/(1 - R\delta)]^{\uparrow 2T} }   - 
   \frac{\sum_{\pi \in S_{2T}} \mathbb{P}_{\pi }^{d}  (L\delta)^{-c(\pi)}}{[d/L\delta]^{\uparrow 2T}}\right)\right]
\end{align}
Besides, since $Z_1^{\otimes 2T} $is a valid  $2T$-fold Clifford either, we also consider the vector $\mathbf{a}$ with $|\mathbf{a}| \mathrm{~mod~} 2 = 0$.


Then use Leaves likelihood, we conclude that $ \sum_{s}|\mathrm{tr}(M_s \Phi_{\mathrm{Cl,2T}}[\Delta_{2T}])|  \leq 1 -  L(s)$, where
\begin{align}
    &L(s) = \min_{M'_s}  \frac{\mathrm{tr}\left[M'_s \cdot\left(\frac{\sum_{\pi \in S_{2T}} \mathbb{P}_{\pi}^{d} (1 - R \delta)^{-c(\pi)} }{[d/(1 - R\delta)]^{\uparrow 2T} }   +\frac{\sum_{\pi \in S_{2T}} \mathbb{P}_{\pi }^{d}  (L\delta)^{-c(\pi)}}{[d/L\delta]^{\uparrow 2T}}\right)\right]}{\mathrm{tr}\left[M'_s\cdot\left(\frac{\sum_{\pi \in S_{2T}} \mathbb{P}_{\pi}^{d}}{d^{\uparrow 2T} }  + \frac{\mathbb{I}(d)^{\otimes 2T}}{ d^{2T}} \right)\right]} \\
    &\geq \frac{\mathrm{tr}\left[M'_s \cdot\left(\frac{[\sum_{\pi \in S_{2}} \mathbb{P}_{\pi}^{d} (1 - R \delta)^{-c(\pi)}]^{\otimes T} }{[d/(1 - R\delta)]^{\uparrow 2T} }   +\frac{[\sum_{\pi \in S_{2}} \mathbb{P}_{\pi }^{d}  (L\delta)^{-c(\pi)}]^{\otimes T}}{[d/L\delta]^{\uparrow 2T}}\right)\right]}{\mathrm{tr}\left[M'_s\cdot\left(\frac{\sum_{\pi \in S_{2T}} \mathbb{P}_{\pi}^{d}}{d^{\uparrow 2T} }  + \frac{\mathbb{I}(d)^{\otimes 2T}}{ d^{2T}} \right)\right]} 
    % &\leq 1 - d^{2T}\min_{M'_s}  \frac{\mathrm{tr}\left[M'_s\left(-\frac{\sum_{\pi \in S_{2}^{\otimes T}} \mathbb{P}_{\pi}^{d}}{d^{\uparrow 2T} } +\frac{\sum_{\pi \in S_{2}^{\otimes T}} \mathbb{P}_{\pi}^{d} (1 - R \delta)^{-c(\pi)} }{[d/(1 - R\delta)]^{\uparrow 2T} }   +\frac{\sum_{\pi \in S_{2}^{\otimes T}} \mathbb{P}_{\pi }^{d}  (L\delta)^{-c(\pi)}}{[d/L\delta]^{\uparrow 2T}}\right)\right]}{\mathrm{tr}\left[M'_s\right]}
\end{align}
Here we use the identity below as proved in Eq.\eqref{eq:t=1_Delta}:
\begin{align}
    &\frac{\left[\sum_{\pi \in S_2} \mathbb{P}_{\pi}^d (1 - R\delta)^{-c(\pi)}\right]^{\otimes T}}{[d/(1 - R\delta)]^{\uparrow 2T} }= \frac{d^T(d+1-R\delta)^T}{[d/(1 - R\delta)]^{\uparrow 2T} }\left[\frac{\mathbb{I}^{\otimes 2} + \mathbb{P}^{d}_{(12)}(1 -R\delta)}{d[d+ 1 -R\delta] }\right]^{\otimes T} \\
    &=d^T(d+1-R\delta)^T\left[-\frac{\mathbb{I}^{\otimes 2} + \mathbb{P}^{d}_{(12)}L\delta}{d[d + L\delta] } +  \frac{\mathbb{I}^{\otimes 2}}{ d^{2}} + \frac{\mathbb{I}^{\otimes 2}  + \mathbb{P}^{d}_{(12)}}{d[d + 1] } \right]^{\otimes T}
\end{align}
If we defined the operator $\mathbb{S}^{\mathbf{a}} = \sum_{\pi \in S_{2T}^{\mathbf{a}}} \mathbb{P}_{\pi}^{d}$ and $D = 2d = 2^n$, we conclude that
\begin{align}
    \Phi_{\mathrm{U},2T}\left[ \ket{{\mathbf{a}}}\!\bra{{\mathbf{a}}}\otimes\mathbb{S}^{\mathbf{a}}\right] = \sum_{\rho \sigma \in S_{2T}} \mathrm{Wg}_{\rho\sigma}^{D} \mathrm{tr}\left[\ket{{\mathbf{a}}}\!\bra{{\mathbf{a}}}\otimes\mathbb{S}^{\mathbf{a}} \cdot \mathrm{SWAP}_{\sigma} \otimes \mathbb{P}_{\sigma}^{d}\right] \mathbb{P}_{\rho}^{D}.
\end{align}
The trace term can be explicitly calcluated as
\begin{align}
    &\mathrm{tr}\left[\ket{{\mathbf{a}}}\!\bra{{\mathbf{a}}}\otimes\mathbb{S}^{\mathbf{a}} \cdot \mathrm{SWAP}_{\sigma} \otimes \mathbb{P}_{\sigma}^{d}\right]= \delta_{\mathbf{a}, \sigma{\mathbf{a}}} \sum_{\pi \in S_{2T}^{\mathbf{a}}} \mathrm{tr}\left(\mathbb{P}_{\pi \sigma}^{d}\right) = \delta_{\sigma, \beta \gamma} \sum_{\pi \in S_{2T}^{{\mathbf{a}}}} d^{c(\pi\gamma)} = \delta_{\sigma, \beta \gamma} \frac{d^{\uparrow |\mathbf{a}|}}{d^{|{\mathbf{a}}|}}d^{c(\gamma)}.
\end{align}
Here in the second equality we use the fact that $\mathbf{a} = \sigma(\mathbf{a})$ if and only if  $\sigma = \beta \cdot \gamma$ for $\beta\in S_{2T}^{\mathbf{a}}$ and $\gamma \in S_{2T}^{\bar{\mathbf{a}}}$. In the last equality, we use the identity 
\begin{align}
   \forall \gamma \in S_{2T}^{\bar{\mathbf{a}}}, \sum_{\pi\in S_{2T}^{\mathbf{a}}} d^{c(\pi\gamma )} =  \frac{d^{\uparrow |\mathbf{a}|}}{d^{|{\mathbf{a}}|}}d^{c(\gamma)}
\end{align}

Noted that the unitary average of $\sum_{\mathbf{a}}2^{-2T}\ket{{\mathbf{a}}}\!\bra{{\mathbf{a}}} \otimes \rho_{\mathbf{a}}$ and $\sum_{\mathbf{a}}2^{-2T}\ket{\bar{\mathbf{a}}}\!\bra{\bar{\mathbf{a}}}\otimes \rho_{\mathbf{a}}$ are identical since $\ket{\mathbf{a}}$ can be transformed into $\ket{\bar{\mathbf{a}}}$ through $2T$-fold unitary $X^{2T}$ and Haar random unitary is shift invariant. 
Thus, the average state $\rho_0$ is the linear combination of permutation operators:
\begin{align}
    &\Phi_{\mathrm{U},2T}\left[ \mathbb{E}_{\psi}\sigma_{(1,0),I}^{\otimes 2T} \right] = \Phi_{\mathrm{U},2T}\left[ \mathbb{E}_{\psi}\sigma_{(0,1),I}^{\otimes 2T} \right] :=\rho_0 =  \Phi_{\mathrm{U},2T}\left[ \sum_{\mathbf{a} \in \mathbb{F}_2^{2T}}\frac{ \ket{{\mathbf{a}}}\!\bra{{\mathbf{a}}}\otimes\mathbb{S}^{\mathbf{a}}}{2^{2T}d^{\uparrow |\mathbf{a}|} d^{2T - |\mathbf{a}|}}\right]\\
    & = \sum_{\mathbf{a} \in \mathbb{F}_2^{2T}}\sum_{\sigma,\rho \in S_{2T}}\frac{\mathrm{Wg}_{\rho \sigma}^{D}\cdot \delta_{\sigma,\beta\gamma}\frac{d^{\uparrow |\mathbf{a}|}}{d^{|{\mathbf{a}}|}}d^{c(\gamma)}}{2^{2T}d^{\uparrow |\mathbf{a}|} d^{2T - |\mathbf{a}|}} \mathbb{P}_{\rho}^{D} 
    =\sum_{\mathbf{a} \in \mathbb{F}_2^{2T}}\sum_{\rho \in S_{2T}}\sum_{\beta \in S_{2T}^{{\mathbf{a}}},\gamma\in S_{2T}^{\bar{\mathbf{a}}}}\frac{\mathrm{Wg}_{\rho(\beta\gamma)}^{D}\cdot d^{c(\gamma)}}{2^{2T} d^{2T}} \mathbb{P}_{\rho}^{D} 
\end{align}
Similarly, we calculate 
\begin{align}
    &\mathrm{tr}\left[\left(\ket{\bar{\mathbf{a}}}\!\bra{\bar{\mathbf{a}}}\otimes \frac{\sum_{\pi \in S_{2T}^{\mathbf{a}}, \nu \in S_{2T}^{\bar{\mathbf{a}}}} \mathbb{P}_{\pi \nu}^{d} (1 - R \delta)^{-c(\pi)} (L\delta)^{-c(\nu)}}{[d/(1 - R\delta)]^{\uparrow |\mathbf{a}|} [d/L\delta]^{\uparrow |\bar{\mathbf{a}}|}}\right)\cdot \mathrm{SWAP}_{\sigma} \otimes \mathbb{P}_{\sigma}^{d}\right]\\
    &= \delta_{\sigma, \beta\gamma}\sum_{\pi \in S_{2T}^{\mathbf{a}}, \nu \in S_{2T}^{\bar{\mathbf{a}}}}\frac{ d^{c(\pi \nu)} (1 - R \delta)^{-c(\pi \beta^{-1})} (L\delta)^{-c(\nu\gamma^{-1})}}{[d/(1 - R\delta)]^{\uparrow |\mathbf{a}|} [d/L\delta]^{\uparrow |\bar{\mathbf{a}}|}}.
\end{align}
When $\delta = 0$, this expression reduced to $\delta_{\sigma,\beta\gamma}d^{-2T} d^{c(\gamma)}$. We thus Taylor expands this expression in terms of $\delta$. 
\begin{align}
    &\Phi_{\mathrm{U},2T}\left[ \mathbb{E}_{\psi}\sigma_{(1-R\delta,L\delta),I}^{\otimes 2T} \right]  = \sum_{\mathbf{a} \in \mathbb{F}_2^{2T}}\sum_{\rho \in S_{2T}}\sum_{\beta \in S_{2T}^{{\mathbf{a}}},\gamma\in S_{2T}^{\bar{\mathbf{a}}}}\mathrm{Wg}_{\rho(\beta\gamma)}^{D}\cdot \frac{ d^{c(\pi \nu)} (1 - R \delta)^{-c(\pi \beta^{-1})} (L\delta)^{-c(\nu\gamma^{-1})}}{2^{2T}[d/(1 - R\delta)]^{\uparrow |\mathbf{a}|} [d/L\delta]^{\uparrow |\bar{\mathbf{a}}|}} \mathbb{P}_{\rho}^{D} 
\end{align}
The difference of two average can then be reform into the discrimination between $\mathbb{I}_{n}^{\otimes 2T}/d^{2T}$ from the mixed state
\begin{align}
    \Phi_{\mathrm{U},2T}\left[ \mathbb{E}_{\psi}\sigma_{(1,0),I}^{\otimes 2T} \right] -\Phi_{\mathrm{U},2T}\left[ \mathbb{E}_{\psi}\sigma_{(1-R\delta,L\delta),I}^{\otimes 2T} \right] =  \mathbb{I}_{n}^{\otimes 2T}/D^{2T} + \left[\sum_{\rho \in S_{2T}} \lambda_{\rho} \mathbb{P}_{\rho}^{D} \right]
\end{align}
This \textbf{TVD} is bounded by the leaves likelihood for $M_s = \otimes_{t =1}^{T} F_{s_t}^{u_t}$ as
\begin{align}
    L(s) = \frac{D^{2T}\mathrm{tr}\left(M_s\left[\sum_{\rho \in S_{2T}} \lambda_{\rho} \mathbb{P}_{\rho}^{D} \right]\right)}{\mathrm{tr}\left(M_s\right)} 
    \geq  D^{2T} \prod_{t = 1}^{2T} \frac{\mathrm{tr}\left(F_{s_t}^{u_t}\left[\sum_{\rho \in S_{2}} \lambda'_{\rho} \mathbb{P}_{\rho}^{D} \right]\right)}{\mathrm{tr}\left(F_{s_t}^{u_t}\right)}
\end{align}
This further reduced the problem to the discrimination between $\mathbb{I}(D)^{\otimes 2}/D^2$ and $\sum_{\rho \in S_{2}} \lambda'_{\rho} \mathbb{P}_{\rho}^{D} $

\begin{lemma}[Martingale distinguishing lower bound]\label{theo:mart_distinguish_lb}
    \textcolor{red}{[To be stated]}
\end{lemma}

Based on Lemma~\ref{theo:mart_distinguish_lb}, we calculate its chi-square divergence first:



\subsubsection{Subspace \textbf{MH}-ensemble test}
\begin{problem}
    Given the state ensemble $\mathcal{E}_{1,\mathrm{Cl}}:= \{\frac{1}{2} \mathrm{Cl}\left[\ket{0}\!\bra{0} \otimes \psi + \ket{1}\!\bra{1} \otimes \mathbb{I}/{2^{n-1}}\right]\mathrm{Cl}^\dagger\}_{\psi \sim \mathrm{Haar}_{n-1}}$ and the mixed state ensemble $\mathcal{E}_{2,\mathrm{Cl}}:=\left\{\mathrm{Cl}\left[{\mathbb{I}}/{2} \otimes \left(\frac{1}{2}\psi_1 + \frac{1}{2}\psi_2\right)\right]\mathrm{Cl}^\dagger\right\}_{\psi_1,\psi_2 \sim \mathrm{Haar}_{n-1}}$, try to design an algorithm using restrained measurements to distinguish these two ensembles with success rate $1/2 + \delta$.
\end{problem}
Given a fixed $n$-qubit Clifford gate $\mathrm{Cl} = \mathbb{I}$, these two ensembles have two-replica average as
\begin{align}
    &\mathbb{E}_{\rho \sim \mathcal{E}_1}\rho^{\otimes 2} = \frac{1}{4}\left[\ket{01}\!\bra{01} + \ket{10}\!\bra{10} + \ket{00}\!\bra{00}\right] \otimes \left(\frac{\mathbb{I}_{n-1}}{d}\right)^{\otimes 2} + \frac{1}{4}\ket{11}\!\bra{11} \otimes \left[\frac{\mathbb{I}_{n-1} \otimes \mathbb{I}_{n-1} + \mathrm{SWAP}^{\otimes (n-1)}}{d(d+1)}\right]\\
    &\mathbb{E}_{\rho \sim \mathcal{E}_2} \rho^{\otimes 2} = \left(\frac{\mathbb{I}}{2}\right)^{\otimes 2} \otimes \frac{1}{2} \left[\frac{\mathbb{I}_{n-1} \otimes \mathbb{I}_{n-1} + \mathrm{SWAP}^{\otimes (n-1)}}{d(d+1)} + \frac{\mathbb{I}_{n-1}}{d} \otimes \frac{\mathbb{I}_{n-1}}{d}\right]\\
    &\Delta = \frac{\ket{00}\!\bra{00} + \ket{10}\!\bra{10} +\ket{01}\!\bra{01} -\ket{11}\!\bra{11}}{4} \otimes \frac{1}{2}\left[\frac{\mathbb{I}_{n-1} \otimes \mathbb{I}_{n-1}}{d^2(d+1)} - \frac{\mathrm{SWAP}^{\otimes (n-1)}}{d(d+1)} \right]
\end{align}
In the case $\mathrm{Cl} = \mathbb{I}$, we have 2-replica localized strategy $\{\ket{00},\ket{01},\ket{10},\ket{11}\}$ on the first qubits and Bell measurement on the last $(n-1)$ qubits such that $\langle \Delta \rangle \sim \mathcal{O}(1)$.  
However, in the case $\mathrm{Cl}$ maps $Z_1$ to a weight $w$ Pauli operator, we show that $\mathrm{tr}(\mathrm{Cl}^{\otimes 2} \Delta \mathrm{Cl}^{\dagger,\otimes 2} \cdot \bigotimes_{j=1}^n\Phi_j) \leq d^2 2^{-\alpha \cdot w}$. This means that non-adaptive 2-replica local measurement strategy at least repeats $2^w$ times to discirminate $\mathcal{E}_{1,\mathrm{Cl}}$ from $\mathcal{E}_{2,\mathrm{Cl}}$
In the case of $T$-round of $2$-replica strategy, we analysis the $2T$-fold average as:
\begin{align}
    \mathbb{E}_{\rho \sim \mathcal{E}_1}\rho^{\otimes 2T} &= \frac{1}{2^{2T}} \sum_{\mathbf{a} \in \mathbb{F}_2^{2T}}\ket{\mathbf{a}}\!\bra{\mathbf{a}} \otimes  \frac{S^{\mathbf{a}}}{d^{\uparrow \mathbf{a}}} \otimes \frac{\mathbb{I}^{\bar{\mathbf{a}}}}{d^{|\bar {\mathbf{a}}|}},\\
    \mathbb{E}_{\rho \sim \mathcal{E}_2}\rho^{\otimes 2T} &= \frac{1}{2^{2T}} \ket{\mathbf{a}}\!\bra{\mathbf{a}} \otimes \sum_{\mathbf{a}} \frac{S^{\mathbf{a}}}{d^{\uparrow \mathbf{a}}}
\end{align}
In the case Clifford gate selected from uniform distribution, we conclude that the average \textcolor{red}{[Section to be completed]}


%%%%%%%%%%%%%%%%%%%%%%%%%%%%%%%%%%%%%%%%%%%%%%%%%%%%%%%%%%%%%%%%%%%%%%%%%%%%%%%%%%%%%%%%%%%%%%%%%%%
%%%%%%%%%%%%%%%%%%%%%%%%%%%%%%%%%%%%%%%%%%%%%%%%%%%%%%%%%%%%%%%%%%%%%%%%%%%%%%%%%%%%%%%%%%%%%%%%%%%


\end{document}